% docs/manual/latex/user_manual.tex - main file of the LaTeX edition of the book "user_manual" (SIMC++ documentation).
%   Compile: cmake --build <build> --target pdf, or tools/build_docs_pdf.py (figures, xelatex passes,
%   makeindex with docs/common/latex/simcpp.ist); by hand, see docs/README.md. XeLaTeX, class
%   docs/common/latex/simcpp.cls (packages: fontspec, tikz, xcolor, hyperref, listings, longtable, booktabs,
%   titlesec, titletoc, fancyhdr, makeidx, preview, ...); fonts TeX Gyre Pagella, TeX Gyre Heros, DejaVu Sans
%   Mono. The figures are built first from docs/{manual,translation}/figures/*.tex (preview.sty) into PDFs.
%   Structure (parts, chapter order) is edited HERE; docs/common/latex/books.json only drives
%   tools/md2latex.py when chapters are re-converted from the Markdown.
\documentclass{simcpp}
\simcsetmeta{SIMC++ User Manual}
  {Installing, configuring and running the Monte Carlo simulation of electron-scattering experiments with magnetic spectrometers}
  {\simcbuildversion}{\simcbuilddate}{SIMC++ documentation}
\begin{document}
\simcfrontmatter
% docs/manual/latex/user/index.tex - chapter of the LaTeX edition of the book "user_manual" (SIMC++ documentation).
%   Converted once from docs/manual/markdown/user/index.md by tools/md2latex.py; from then on THIS FILE is the source of the
%   LaTeX edition and may be edited by hand. Re-running tools/md2latex.py --write-sources for this
%   chapter OVERWRITES the hand edits (see docs/README.md).
%   Compile: cmake --build <build> --target pdf, or tools/build_docs_pdf.py (figures, xelatex passes,
%   makeindex with docs/common/latex/simcpp.ist); by hand, see docs/README.md. XeLaTeX, class
%   docs/common/latex/simcpp.cls (packages: fontspec, tikz, xcolor, hyperref, listings, longtable, booktabs,
%   titlesec, titletoc, fancyhdr, makeidx, preview, ...); fonts TeX Gyre Pagella, TeX Gyre Heros, DejaVu Sans
%   Mono. The figures are built first from docs/{manual,translation}/figures/*.tex (preview.sty) into PDFs.
\chapter*{About this manual}\addcontentsline{toc}{chapter}{About this manual}\markboth{About this manual}{About this manual}\label{user_manual-index}
How to install and run SIMC++. No knowledge of the legacy Fortran SIMC is needed. The developer manual
(../\allowbreak{}developer/\allowbreak{}index.\allowbreak{}md\footnote{SIMC++ Developer Manual, ``SIMC++ developer manual'' (\code{docs/\allowbreak{}manual/\allowbreak{}markdown/\allowbreak{}developer/\allowbreak{}index.\allowbreak{}md})}) covers maintaining and extending the code.

\begin{simctable}{small}{2}{>{\raggedright\arraybackslash}p{\dimexpr 0.298\linewidth-2\tabcolsep\relax}>{\raggedright\arraybackslash}p{\dimexpr 0.655\linewidth-2\tabcolsep\relax}}
\caption{About this manual}\\%
\toprule\rowcolor{simchead}\simcth{chapter} & \simcth{content} \\\midrule\endfirsthead
\multicolumn{2}{@{}l}{\simccontinued}\\\toprule\rowcolor{simchead}\simcth{chapter} & \simcth{content} \\\midrule\endhead
\midrule\multicolumn{2}{r@{}}{\simccontinues}\\\endfoot
\bottomrule\endlastfoot
\hyperref[user_manual-01u5fintroduction]{1. Introduction} & what SIMC++ simulates, run modes, limitations \\
\hyperref[user_manual-02u5finstallation]{2. Installation} & requirements,\allowbreak{} build, install, data directory (\code{SIMC\_\allowbreak{}DATA\_\allowbreak{}DIR}, \code{-{}-{}print-{}data-{}dir}\simcidx{Command-line options!print-data-dir@\code{-{}-{}print-{}data-{}dir}}), tests \\
\hyperref[user_manual-03u5fquicku5fstart]{3. Quick start} & run an example input, the outputs, look at the tree in ROOT, the same run from JSON \\
\hyperref[user_manual-04u5finputs]{4. Inputs} & \code{.\allowbreak{}inp} syntax, old names, JSON, plug-in selection, \textbf{complete parameter reference} \\
\hyperref[user_manual-05u5freactions]{5. Reactions} & (e,e'p), pion, kaon, Delta, SIDIS, rho, phase space: what each computes and its limitations \\
\hyperref[user_manual-06u5fspectrometers]{6. Spectrometers} & HMS, SHMS, SOS, HRS, calorimeter; spectrometer numbers; \hyperref[param:spectu5fmode]{\code{spect\_\allowbreak{}mode}}\simcidx{Parameters!spect_mode@\code{spect\_mode}}; target field \\
\hyperref[user_manual-07u5foutputs]{7. Outputs} & tree branches per reaction, \code{.\allowbreak{}hist}, \code{.\allowbreak{}geni}, \code{.\allowbreak{}gen}, summary JSON, exit status \\
\hyperref[user_manual-08u5frunningu5fandu5freproducibility]{8. Running modes and reproducibility} & seeds, mt64 vs RANLUX, the legacy-compatible mode, reproducing legacy SIMC results exactly, \code{-{}-{}threads}\simcidx{Command-line options!threads@\code{-{}-{}threads}} \\
\hyperref[user_manual-09u5fperformance]{9. Performance tips} & speed, threads, memory \\
\hyperref[user_manual-10u5ffaq]{10. FAQ and troubleshooting} & common stops, old inputs, data not found, NaN columns, ... \\
Example inputs\footnotemark{} & \code{h\_\allowbreak{}elastic\_\allowbreak{}rad.\allowbreak{}inp} (chapter 3), \code{legacy\_\allowbreak{}names.\allowbreak{}inp} (chapter 4.1.3) \\
\hyperref[user_manual-VERIFIEDu5fEXAMPLES]{Verified examples} & transcript of every command of this manual, as run on the stated commit \\
\end{simctable}
\addtocounter{footnote}{-1}
\stepcounter{footnote}\footnotetext{\code{docs/\allowbreak{}manual/\allowbreak{}markdown/\allowbreak{}user/\allowbreak{}examples/}}
Reference documents: \code{docs/\allowbreak{}manual/\allowbreak{}markdown/\allowbreak{}reference/\allowbreak{}INP\_\allowbreak{}JSON\_\allowbreak{}MAPPING.\allowbreak{}md} (the \code{.\allowbreak{}inp} → JSON table), \code{docs/\allowbreak{}manual/\allowbreak{}markdown/\allowbreak{}reference/\allowbreak{}OUTPUT\_\allowbreak{}FORMAT.\allowbreak{}md} (every output value),
\code{docs/\allowbreak{}manual/\allowbreak{}markdown/\allowbreak{}reference/\allowbreak{}DRIVER.\allowbreak{}md}, \code{docs/\allowbreak{}manual/\allowbreak{}markdown/\allowbreak{}reference/\allowbreak{}PERFORMANCE.\allowbreak{}md}. Comparison with the legacy Fortran SIMC and the list of differences:
\code{docs/\allowbreak{}translation/}.

Conventions: commands are shown as typed in a shell with \code{\$SIMCPP\_\allowbreak{}PREFIX/\allowbreak{}bin} on the \code{PATH} (directory layout: \hyperref[user_manual-02u5finstallation:21-directory-layout-used-in-this-manual]{2.1}). \code{\$EX} is the directory of the example inputs (\code{\$SIMCPP\_\allowbreak{}PREFIX/\allowbreak{}share/\allowbreak{}simc/\allowbreak{}docs/\allowbreak{}manual/\allowbreak{}markdown/\allowbreak{}user/\allowbreak{}examples}).
\code{\$SIMC\_\allowbreak{}REFERENCE} (the project's reference runs of the legacy program) is needed only for the example of chapter 8.4.


% docs/manual/latex/front/user_reading.tex - front matter of the user manual (PDF only): how to read this book.
% Hand-written LaTeX (no Markdown counterpart); \input by the main .tex of the book.
%   Compile: cmake --build <build> --target pdf, or python3 -I tools/build_docs_pdf.py --source-root . --build-dir DIR
%   (figures, then xelatex passes + makeindex -s docs/common/latex/simcpp.ist); by hand see docs/README.md.
%   XeLaTeX; class docs/common/latex/simcpp.cls (TeX Live packages listed there); fonts TeX Gyre Pagella,
%   TeX Gyre Heros, DejaVu Sans Mono. The figures are built first from docs/{manual,translation}/figures/*.tex
%   (preview.sty) into PDFs.
\section*{How to read this book}
The manual is in four parts. \textbf{Getting started} (chapters~\ref{user_manual-01u5fintroduction}
to~\ref{user_manual-03u5fquicku5fstart}) is enough to install SIMC++ and make a first run. \textbf{Configuring a
simulation} explains the input files and contains the complete parameter reference
(section~\ref{user_manual-04u5finputs:46-parameter-reference}), then what each reaction and each spectrometer does.
\textbf{Results and reproducibility} describes every output file, the random-number generators, the
legacy-compatible mode and multi-threaded runs. \textbf{Help} collects the common messages and problems. The appendix
is a transcript of every command of the manual, run on the version printed on the title page.

\input{conventions}
\simcmainmatter

\simcpart{Getting started}
  {What SIMC++ simulates, how to build and install it, and a first run from an example input.}
% docs/manual/latex/user/01_introduction.tex - chapter of the LaTeX edition of the book "user_manual" (SIMC++ documentation).
%   Converted once from docs/manual/markdown/user/01_introduction.md by tools/md2latex.py; from then on THIS FILE is the source of the
%   LaTeX edition and may be edited by hand. Re-running tools/md2latex.py --write-sources for this
%   chapter OVERWRITES the hand edits (see docs/README.md).
%   Compile: cmake --build <build> --target pdf, or tools/build_docs_pdf.py (figures, xelatex passes,
%   makeindex with docs/common/latex/simcpp.ist); by hand, see docs/README.md. XeLaTeX, class
%   docs/common/latex/simcpp.cls (packages: fontspec, tikz, xcolor, hyperref, listings, longtable, booktabs,
%   titlesec, titletoc, fancyhdr, makeidx, preview, ...); fonts TeX Gyre Pagella, TeX Gyre Heros, DejaVu Sans
%   Mono. The figures are built first from docs/{manual,translation}/figures/*.tex (preview.sty) into PDFs.
\chapter{\texorpdfstring{Introduction}{Introduction}}\label{user_manual-01u5fintroduction}
\section{\texorpdfstring{What SIMC++ is}{What SIMC++ is}}\label{user_manual-01u5fintroduction:11-what-simc-is}
SIMC++ is a Monte Carlo for coincidence electron-scattering experiments A(e,e'X) with magnetic spectrometers (Jefferson Lab
Hall C, and the Hall A HRS). It generates events over a phase space matched to the spectrometer acceptances and weights
them with a physics model. Each event goes through:

\begin{itemize}
\item \textbf{beam and target}: beam energy spread, beam size and raster, energy loss and multiple scattering in the target and its
windows, Coulomb correction;
\item \textbf{reaction}: event generation and the cross section of the chosen physics model;
\item \textbf{radiative corrections}: radiative tails of the incoming electron, the scattered electron and the hadron;
\item \textbf{spectrometers}: optics maps, apertures, decay in flight, multiple scattering, detector resolution and reconstruction;
\item \textbf{normalisation}: luminosity, generation volume and event weight, so that the weighted events give a yield for the
simulated beam charge.
\end{itemize}
SIMC++ is not a detector simulation (GEANT4) and not a general background generator: events are generated only where the
spectrometers can see them.

It is a C++20 / ROOT 6 program, \code{simc++}. It reads the classic SIMC input files (\code{.\allowbreak{}inp}) as they are, or a JSON
configuration that can also choose the components ("plug-ins") of the simulation. It writes a ROOT tree with one entry per
accepted event, text reports and a JSON run summary (\hyperref[user_manual-07u5foutputs]{chapter 7}).

(SIMC++ is a translation of the Fortran SIMC, \code{simc\_\allowbreak{}gfortran}. The Fortran code is not needed to build or run it. How the
translation was validated is described in \code{docs/\allowbreak{}translation/}.)

\textbf{Which Fortran SIMC.} SIMC++ reproduces the official Fortran SIMC of Jefferson Lab (\code{Jefferson\allowbreak{}Lab/\allowbreak{}simc\_\allowbreak{}gfortran}) at
master commit \textbf{59ccdf6} (2026-10-01),\allowbreak{} the version recorded in the file \code{UPSTREAM\_\allowbreak{}SIMC} of the source tree. The default
column layout of the ROOT tree follows 59ccdf6; \code{-{}-{}tree-{}layout\ legacy}\simcidx{Command-line options!tree-layout@\code{-{}-{}tree-{}layout}} writes the layout of the older version 294f662
(\hyperref[user_manual-07u5foutputs]{7.1}), and \code{-{}-{}fortran-{}compat\ -{}-{}compat-{}target\ 294f662}\simcidx{Command-line options!fortran-compat@\code{-{}-{}fortran-{}compat}} reproduces the behaviour of 294f662
(\hyperref[user_manual-08u5frunningu5fandu5freproducibility:83-default-and-legacy-compatible-mode]{8.3}). The SIMC++ release that reproduced
294f662 is the git tag \code{v1.\allowbreak{}0-{}fortran294f662} (\code{git\ checkout\ v1.\allowbreak{}0-{}fortran294f662}, \hyperref[user_manual-02u5finstallation:21-directory-layout-used-in-this-manual]{2.1});
the release that reproduces 59ccdf6 is the tag \code{v2.\allowbreak{}0-{}fortran59ccdf6}.

\begin{figure}[htbp]
\simcfigure{run_flow}
\caption[A run of SIMC++]{A run of SIMC++: the inputs are read into SimcParameters and RunSetup, the pipeline creates the stages by name, the event loop calls them once per try, and finish() normalises and writes the outputs through the sinks}\label{fig:runu5fflow}
\end{figure}
\section{\texorpdfstring{What it simulates}{What it simulates}}\label{user_manual-01u5fintroduction:12-what-it-simulates}
\begin{simctable}{small}{3}{>{\raggedright\arraybackslash}p{\dimexpr 0.559\linewidth-2\tabcolsep\relax}>{\raggedright\arraybackslash}p{\dimexpr 0.232\linewidth-2\tabcolsep\relax}>{\raggedright\arraybackslash}p{\dimexpr 0.141\linewidth-2\tabcolsep\relax}}
%
\toprule\rowcolor{simchead}\simcth{reaction (input flags)} & \simcth{targets} & \simcth{chapter} \\\midrule\endfirsthead
\multicolumn{3}{@{}l}{\simccontinued}\\\toprule\rowcolor{simchead}\simcth{reaction (input flags)} & \simcth{targets} & \simcth{chapter} \\\midrule\endhead
\midrule\multicolumn{3}{r@{}}{\simccontinues}\\\endfoot
\bottomrule\endlastfoot
elastic and quasi-elastic (e,e'p) (no reaction flag) & H, D, nuclei (theory or Benhar spectral functions) & \hyperref[user_manual-05u5freactions:52-heep-elastic-deep-aeep-and-aeea]{5.2} \\
exclusive pion production (\hyperref[param:doingu5fpion]{\code{doing\_\allowbreak{}pion}}\simcidx{Parameters!doing_pion@\code{doing\_pion}}): π+, π−, π0, Δ final states, coherent & H, D, nuclei & \hyperref[user_manual-05u5freactions:53-exclusive-pion-production-doingu5fpion]{5.3} \\
kaon electroproduction (\hyperref[param:doingu5fkaon]{\code{doing\_\allowbreak{}kaon}}\simcidx{Parameters!doing_kaon@\code{doing\_kaon}}): K+Λ, K+Σ0, K+Σ−, hypernuclei & H, D, nuclei & \hyperref[user_manual-05u5freactions:54-kaon-electroproduction-doingu5fkaon]{5.4} \\
(e,e'p)π in the Delta region (\hyperref[param:doingu5fdelta]{\code{doing\_\allowbreak{}delta}}\simcidx{Parameters!doing_delta@\code{doing\_delta}}) & H, D, nuclei & \hyperref[user_manual-05u5freactions:55-eepu3c0-in-the-delta-region-doingu5fdelta]{5.5} \\
semi-inclusive (e,\allowbreak{}e'π±/\allowbreak{}π0)X and (e,e'K±)X (\hyperref[param:doingu5fsemi]{\code{doing\_\allowbreak{}semi}}\simcidx{Parameters!doing_semi@\code{doing\_semi}}) & H, D, nuclei & \hyperref[user_manual-05u5freactions:56-semi-inclusive-production-doingu5fsemi]{5.6} \\
exclusive ρ0 → π+π− (\hyperref[param:doingu5frho]{\code{doing\_\allowbreak{}rho}}\simcidx{Parameters!doing_rho@\code{doing\_rho}}) & H, D, ³He & \hyperref[user_manual-05u5freactions:57-exclusive-u3c10-production-doingu5frho]{5.7} \\
phase-space (acceptance) runs (\hyperref[param:doingu5fphsp]{\code{doing\_\allowbreak{}phsp}}\simcidx{Parameters!doing_phsp@\code{doing\_phsp}}) & any & \hyperref[user_manual-05u5freactions:58-phase-space-and-acceptance-runs-doingu5fphsp]{5.8} \\
\end{simctable}
Spectrometers: HMS, SHMS, SOS, HRS-R, HRS-L and an NPS-like calorimeter arm with π0 → γγ detection; tracking through the
magnetic field of a polarized target (\hyperref[user_manual-06u5fspectrometers]{chapter 6}).

\section{\texorpdfstring{Run modes}{Run modes}}\label{user_manual-01u5fintroduction:13-run-modes}
\begin{simctable}{small}{3}{>{\raggedright\arraybackslash}p{\dimexpr 0.209\linewidth-2\tabcolsep\relax}>{\raggedright\arraybackslash}p{\dimexpr 0.324\linewidth-2\tabcolsep\relax}>{\raggedright\arraybackslash}p{\dimexpr 0.399\linewidth-2\tabcolsep\relax}}
%
\toprule\rowcolor{simchead}\simcth{mode} & \simcth{how} & \simcth{use} \\\midrule\endfirsthead
\multicolumn{3}{@{}l}{\simccontinued}\\\toprule\rowcolor{simchead}\simcth{mode} & \simcth{how} & \simcth{use} \\\midrule\endhead
\midrule\multicolumn{3}{r@{}}{\simccontinues}\\\endfoot
\bottomrule\endlastfoot
default & nothing to do & physics runs \\
legacy-compatible & \code{-{}-{}fortran-{}compat}\simcidx{Command-line options!fortran-compat@\code{-{}-{}fortran-{}compat}}, optionally with the \code{ranlux} random-number generator & reproduce results of the legacy Fortran SIMC (59ccdf6; 294f662 with \code{-{}-{}compat-{}target\ 294f662}\simcidx{Command-line options!compat-target@\code{-{}-{}compat-{}target}}) exactly, including its known bugs \\
\end{simctable}
Both modes and the other run options (seeds, generators, threads) are described in
\hyperref[user_manual-08u5frunningu5fandu5freproducibility]{chapter 8}. The log line \code{\ SIMC++:\ fortran\_\allowbreak{}compat\ =\ false\ (\allowbreak{}.\allowbreak{}.\allowbreak{}.\allowbreak{})} tells which mode a run
used.

\section{\texorpdfstring{Validation in one paragraph}{Validation in one paragraph}}\label{user_manual-01u5fintroduction:14-validation-in-one-paragraph}
Every physics routine was checked number for number against the original code, and complete runs of all the example
inputs distributed with SIMC reproduce the original output files exactly in the legacy-compatible mode. Default-mode runs
agree statistically,\allowbreak{} except where a known bug of the original is corrected. The test suite keeps this true for every
change (\hyperref[user_manual-02u5finstallation:26-running-the-tests]{2.6}); the evidence is in \code{docs/\allowbreak{}translation/}.

\begin{figure}[htbp]
\simcfigure{validation_ladder}
\caption[The validation ladder]{The validation ladder: single routines on fixed inputs, the whole event chain on the random numbers of the original program, complete runs identical bit for bit, and complete runs compared statistically}\label{fig:validationu5fladder}
\end{figure}
\section{\texorpdfstring{Limitations}{Limitations}}\label{user_manual-01u5fintroduction:15-limitations}
\begin{itemize}
\item \textbf{Delta-region runs have no cross section in the weight} (no Delta model is implemented): the weight is phase space ×
virtual-photon flux × a jacobian, so Delta yields are not physical rates (\hyperref[user_manual-05u5freactions:55-eepu3c0-in-the-delta-region-doingu5fdelta]{5.5}).
\item \textbf{Acceptance runs on hydrogen} keep the elastic kinematics, and radiation stays on unless \code{using\_\allowbreak{}rad\ =\ 0}
(\hyperref[user_manual-05u5freactions:58-phase-space-and-acceptance-runs-doingu5fphsp]{5.8}).
\item \textbf{Model limitations}, for example: coherent pion and kaon production use the nucleon models; the ρ0 on D/³He has no
Fermi motion; a π0 sent into the HMS/SHMS is transported like a charged particle. Each reaction section of
\hyperref[user_manual-05u5freactions]{chapter 5} lists them.
\item No \code{.\allowbreak{}bin} file and no PAW ntuple: the ROOT tree is written directly.
\item The \code{debug(\allowbreak{}2.\allowbreak{}.\allowbreak{}6)} diagnostic print-outs are not implemented (\hyperref[user_manual-04u5finputs:46-parameter-reference]{4.6}).
\end{itemize}

% docs/manual/latex/user/02_installation.tex - chapter of the LaTeX edition of the book "user_manual" (SIMC++ documentation).
%   Converted once from docs/manual/markdown/user/02_installation.md by tools/md2latex.py; from then on THIS FILE is the source of the
%   LaTeX edition and may be edited by hand. Re-running tools/md2latex.py --write-sources for this
%   chapter OVERWRITES the hand edits (see docs/README.md).
%   Compile: cmake --build <build> --target pdf, or tools/build_docs_pdf.py (figures, xelatex passes,
%   makeindex with docs/common/latex/simcpp.ist); by hand, see docs/README.md. XeLaTeX, class
%   docs/common/latex/simcpp.cls (packages: fontspec, tikz, xcolor, hyperref, listings, longtable, booktabs,
%   titlesec, titletoc, fancyhdr, makeidx, preview, ...); fonts TeX Gyre Pagella, TeX Gyre Heros, DejaVu Sans
%   Mono. The figures are built first from docs/{manual,translation}/figures/*.tex (preview.sty) into PDFs.
\chapter{\texorpdfstring{Installation}{Installation}}\label{user_manual-02u5finstallation}
\simcabstract{\textbf{SIMC++ is self-contained.} You do not need the Fortran SIMC, CERNLIB or any Fortran compiler to build, install or run it.
The model tables, optics maps and spectral functions it needs are part of SIMC++ (\code{data/}). The Fortran code and its
reference runs are used only by optional validation tests, which are skipped when those files are absent
(\hyperref[user_manual-02u5finstallation:26-running-the-tests]{2.6}).}
\section{\texorpdfstring{Directory layout used in this manual}{Directory layout used in this manual}}\label{user_manual-02u5finstallation:21-directory-layout-used-in-this-manual}
All commands of this manual use these locations (adapt them to your machine):

\begin{simctable}{small}{3}{>{\raggedright\arraybackslash}p{\dimexpr 0.209\linewidth-2\tabcolsep\relax}>{\raggedright\arraybackslash}p{\dimexpr 0.228\linewidth-2\tabcolsep\relax}>{\raggedright\arraybackslash}p{\dimexpr 0.495\linewidth-2\tabcolsep\relax}}
%
\toprule\rowcolor{simchead}\simcth{variable} & \simcth{example} & \simcth{content} \\\midrule\endfirsthead
\multicolumn{3}{@{}l}{\simccontinued}\\\toprule\rowcolor{simchead}\simcth{variable} & \simcth{example} & \simcth{content} \\\midrule\endhead
\midrule\multicolumn{3}{r@{}}{\simccontinues}\\\endfoot
\bottomrule\endlastfoot
\code{\$SIMCPP\_\allowbreak{}SRC} & \code{\textasciitilde{}/\allowbreak{}src/\allowbreak{}simcpp} & the git checkout of SIMC++ (the source tree) \\
\code{\$SIMCPP\_\allowbreak{}BUILD} & \code{\textasciitilde{}/\allowbreak{}build/\allowbreak{}simcpp} & the CMake build directory \\
\code{\$SIMCPP\_\allowbreak{}PREFIX} & \code{\textasciitilde{}/\allowbreak{}opt/\allowbreak{}simcpp} & the installation prefix (\code{bin/}, \code{lib/}, \code{include/}, \code{share/\allowbreak{}simc/\allowbreak{}data}) \\
\code{\$SIMC\_\allowbreak{}FORTRAN} & \code{\textasciitilde{}/\allowbreak{}src/\allowbreak{}simc\_\allowbreak{}gfortran} & optional: the legacy Fortran SIMC, only for the validation tests \\
\code{\$SIMC\_\allowbreak{}REFERENCE} & \code{\textasciitilde{}/\allowbreak{}simc\_\allowbreak{}reference} & optional: reference runs and tables of the legacy program, only for validation \\
\code{\$ROOTSYS} & \code{\textasciitilde{}/\allowbreak{}opt/\allowbreak{}root} & the ROOT 6.40 installation \\
\end{simctable}
\begin{simcshell}[basicstyle=\ttfamily\scriptsize]
source ~/opt/root/bin/thisroot.sh       # ROOT's set-up script: sets ROOTSYS, PATH and LD_LIBRARY_PATH
export SIMCPP_SRC=~/src/simcpp SIMCPP_BUILD=~/build/simcpp SIMCPP_PREFIX=~/opt/simcpp
export PATH=$SIMCPP_PREFIX/bin:$PATH
\end{simcshell}
Source ROOT's \code{thisroot.\allowbreak{}sh} in every new shell before configuring,\allowbreak{} building or running SIMC++ (it is the usual ROOT
set-up). Setting \code{ROOTSYS}, \code{LD\_\allowbreak{}LIBRARY\_\allowbreak{}PATH=\allowbreak{}\$ROOTSYS/\allowbreak{}lib} and \code{PATH=\allowbreak{}\$ROOTSYS/\allowbreak{}bin:\allowbreak{}\$PATH} by hand works too.

\textbf{Which version to check out.} The \code{main} branch reproduces the Fortran SIMC at master commit 59ccdf6 (the file
\code{UPSTREAM\_\allowbreak{}SIMC}; the default tree layout is that of 59ccdf6, \code{-{}-{}tree-{}layout\ legacy}\simcidx{Command-line options!tree-layout@\code{-{}-{}tree-{}layout}} gives the 294f662 layout). The
release that reproduces 59ccdf6 is the tag \code{v2.\allowbreak{}0-{}fortran59ccdf6}; the release that reproduced the older Fortran version
294f662 is the tag \code{v1.\allowbreak{}0-{}fortran294f662}:

\begin{simcshell}[basicstyle=\ttfamily\scriptsize]
git -C $SIMCPP_SRC checkout v2.0-fortran59ccdf6   # the release for Fortran SIMC 59ccdf6
git -C $SIMCPP_SRC checkout v1.0-fortran294f662   # the release for Fortran SIMC 294f662
git -C $SIMCPP_SRC checkout main                  # back to the current version (Fortran SIMC 59ccdf6)
\end{simcshell}
\section{\texorpdfstring{Requirements}{Requirements}}\label{user_manual-02u5finstallation:22-requirements}
\begin{simctable}{small}{3}{>{\raggedright\arraybackslash}p{\dimexpr 0.154\linewidth-2\tabcolsep\relax}>{\raggedright\arraybackslash}p{\dimexpr 0.393\linewidth-2\tabcolsep\relax}>{\raggedright\arraybackslash}p{\dimexpr 0.385\linewidth-2\tabcolsep\relax}}
%
\toprule\rowcolor{simchead}\simcth{item} & \simcth{version} & \simcth{note} \\\midrule\endfirsthead
\multicolumn{3}{@{}l}{\simccontinued}\\\toprule\rowcolor{simchead}\simcth{item} & \simcth{version} & \simcth{note} \\\midrule\endhead
\midrule\multicolumn{3}{r@{}}{\simccontinues}\\\endfoot
\bottomrule\endlastfoot
C++ compiler & C++20 (GCC 16.2 and clang 22 are used by the project; both build without warnings with \code{-{}Wall\ -{}Wextra\ -{}pedantic}) &  \\
CMake & ≥ 3.15 &  \\
ROOT & ≥ 6.40, built with C++20; components Core, Tree, Hist, RIO (the comparison tool also needs Gpad, Graf, MathCore) & the project uses ROOT 6.40.04 \\
nlohmann/\allowbreak{}json & ≥ 3.9 & downloaded automatically (FetchContent,\allowbreak{} v3.11.2) when it is not installed; this needs network access at configure time \\
Python 3 & any recent & only for the \code{docs} tests; they are not registered without Python \\
\end{simctable}
Runtime library path: \code{LD\_\allowbreak{}LIBRARY\_\allowbreak{}PATH} must contain the \code{lib} directory of the ROOT that SIMC++ was built with, and no
directory with another ROOT or another \code{libtbb.\allowbreak{}so} (an older ROOT in \code{/\allowbreak{}usr/\allowbreak{}local/\allowbreak{}lib}, for example, breaks ROOT 6.40).

\section{\texorpdfstring{Building}{Building}}\label{user_manual-02u5finstallation:23-building}
\begin{simcshell}[basicstyle=\ttfamily\scriptsize]
cmake -S $SIMCPP_SRC -B $SIMCPP_BUILD -DCMAKE_PREFIX_PATH=$ROOTSYS -DCMAKE_BUILD_TYPE=Release
cmake --build $SIMCPP_BUILD -j8
\end{simcshell}
The programs are in \code{\$SIMCPP\_\allowbreak{}BUILD/\allowbreak{}bin/}:

\begin{simctable}{small}{2}{>{\raggedright\arraybackslash}p{\dimexpr 0.234\linewidth-2\tabcolsep\relax}>{\raggedright\arraybackslash}p{\dimexpr 0.719\linewidth-2\tabcolsep\relax}}
%
\toprule\rowcolor{simchead}\simcth{program} & \simcth{purpose} \\\midrule\endfirsthead
\multicolumn{2}{@{}l}{\simccontinued}\\\toprule\rowcolor{simchead}\simcth{program} & \simcth{purpose} \\\midrule\endhead
\midrule\multicolumn{2}{r@{}}{\simccontinues}\\\endfoot
\bottomrule\endlastfoot
\code{simc++} & the Monte Carlo \\
\code{inp2json} & converts a \code{.\allowbreak{}inp} input (and its \hyperref[param:extrau5fdbaseu5ffile]{\code{extra\_\allowbreak{}dbase\_\allowbreak{}file}}\simcidx{Parameters!extra_dbase_file@\code{extra\_dbase\_file}}) into a SIMC++ JSON configuration (\hyperref[user_manual-04u5finputs:42-the-json-format]{4.2}) \\
\code{hist2json} & converts a \code{.\allowbreak{}hist} + \code{.\allowbreak{}geni} report into the run-summary JSON format \\
\code{simc\_\allowbreak{}root\_\allowbreak{}identity} & compares two output ROOT files entry by entry and branch by branch (bitwise) \\
\end{simctable}
CMake options:

\begin{simctable}{footnotesize}{3}{>{\raggedright\arraybackslash}p{\dimexpr 0.218\linewidth-2\tabcolsep\relax}>{\raggedright\arraybackslash}p{\dimexpr 0.219\linewidth-2\tabcolsep\relax}>{\raggedright\arraybackslash}p{\dimexpr 0.501\linewidth-2\tabcolsep\relax}}
%
\toprule\rowcolor{simchead}\simcth{option} & \simcth{default} & \simcth{meaning} \\\midrule\endfirsthead
\multicolumn{3}{@{}l}{\simccontinued}\\\toprule\rowcolor{simchead}\simcth{option} & \simcth{default} & \simcth{meaning} \\\midrule\endhead
\midrule\multicolumn{3}{r@{}}{\simccontinues}\\\endfoot
\bottomrule\endlastfoot
\code{CMAKE\_\allowbreak{}BUILD\_\allowbreak{}TYPE} & (none) & use \code{Release} (\code{-{}O3}). Do \textbf{not} add \code{-{}march=\allowbreak{}native} or \code{-{}ffast-{}math}: they change the floating-point results and the exact-reproduction tests then fail \\
\code{SIMC\_\allowbreak{}BUILD\_\allowbreak{}EXAMPLES} & \code{ON} & build the example plug-ins (\code{ideal\_\allowbreak{}box}) into \code{simc++} \\
\code{SIMC\_\allowbreak{}REFERENCE\_\allowbreak{}DIR}, \code{SIMC\_\allowbreak{}FORTRAN\_\allowbreak{}DIR}, \code{SIMC\_\allowbreak{}BASELINE\_\allowbreak{}DIR} & the environment variables \code{SIMC\_\allowbreak{}REFERENCE}, \code{SIMC\_\allowbreak{}FORTRAN}, \code{SIMC\_\allowbreak{}BASELINES}, else empty & \textbf{validation tests only}: the legacy reference data, the legacy source tree and the regression baselines. Not needed otherwise; when empty, the tests that need them are skipped (\hyperref[user_manual-02u5finstallation:26-running-the-tests]{2.6}) \\
\end{simctable}
The statistical comparison tool \code{simc\_\allowbreak{}compare} is a separate, optional CMake project in \code{validation/} (see
\code{docs/\allowbreak{}translation/\allowbreak{}markdown/\allowbreak{}user\_\allowbreak{}comparing\_\allowbreak{}with\_\allowbreak{}fortran.\allowbreak{}md}).

\section{\texorpdfstring{Installing}{Installing}}\label{user_manual-02u5finstallation:24-installing}
\begin{simcshell}[basicstyle=\ttfamily\footnotesize]
cmake --install $SIMCPP_BUILD --prefix $SIMCPP_PREFIX
\end{simcshell}
This installs:

\begin{itemize}
\item \code{\$SIMCPP\_\allowbreak{}PREFIX/\allowbreak{}bin/} with the four programs;
\item \code{\$SIMCPP\_\allowbreak{}PREFIX/\allowbreak{}share/\allowbreak{}simc/\allowbreak{}data/} with the data files (optics maps, spectral functions, pion/\allowbreak{}kaon/\allowbreak{}SIDIS model tables, the
target-field map, example configurations; each directory has an \code{MD5SUMS} file);
\item \code{\$SIMCPP\_\allowbreak{}PREFIX/\allowbreak{}share/\allowbreak{}simc/\allowbreak{}docs/} (this manual and its example inputs in \code{docs/\allowbreak{}manual/\allowbreak{}markdown/\allowbreak{}user/\allowbreak{}examples/}),
\code{\$SIMCPP\_\allowbreak{}PREFIX/\allowbreak{}share/\allowbreak{}simc/\allowbreak{}examples/} and \code{\$SIMCPP\_\allowbreak{}PREFIX/\allowbreak{}share/\allowbreak{}simc/\allowbreak{}README.\allowbreak{}md};
\item the headers, libraries and the CMake package \code{SIMC} (\code{find\_\allowbreak{}package(\allowbreak{}SIMC)}), needed only to write your own plug-ins
(developer manual).
\end{itemize}
With \code{\$SIMCPP\_\allowbreak{}PREFIX/\allowbreak{}bin} on the \code{PATH} the commands of this manual can be typed as \code{simc++\ .\allowbreak{}.\allowbreak{}.}. Without installing, use
\code{\$SIMCPP\_\allowbreak{}BUILD/\allowbreak{}bin/\allowbreak{}simc++}; both find their data automatically (\hyperref[user_manual-02u5finstallation:25-where-the-data-files-are-found]{2.5}).

\begin{figure}[htbp]
\simcfigure{source_layout}
\caption[The source tree and the installed layout]{The source tree and the installed layout: cmake -{}-build compiles the code, cmake -{}-install puts programs, libraries, headers, data and documentation under the prefix}\label{fig:sourceu5flayout}
\end{figure}
\section{\texorpdfstring{Where the data files are found}{Where the data files are found}}\label{user_manual-02u5finstallation:25-where-the-data-files-are-found}
\code{simc++} reads its tables from one data directory, the first of these that applies (\code{include/\allowbreak{}simc/\allowbreak{}core/\allowbreak{}Data\allowbreak{}Path.\allowbreak{}h}):

\begin{enumerate}
\item the environment variable \code{SIMC\_\allowbreak{}DATA\_\allowbreak{}DIR}, if it is set and not empty;
\item for a program built against the installed CMake package (\code{find\_\allowbreak{}package(\allowbreak{}SIMC)}): that package's
\code{<prefix>\allowbreak{}/\allowbreak{}share/\allowbreak{}simc/\allowbreak{}data};
\item for an installed program (\code{<prefix>\allowbreak{}/\allowbreak{}bin/\allowbreak{}simc++}): \code{<prefix>\allowbreak{}/\allowbreak{}share/\allowbreak{}simc/\allowbreak{}data}, if it exists;
\item the source tree the program was built from (\code{<source>\allowbreak{}/\allowbreak{}data}), for in-tree builds and the tests.
\end{enumerate}
\code{simc++\ -{}-{}print-{}data-{}dir} prints the directory it will use and exits:

\begin{simcshell}[basicstyle=\ttfamily\footnotesize]
$ simc++ --print-data-dir
~/opt/simcpp/share/simc/data
$ SIMC_DATA_DIR=/data/simc simc++ --print-data-dir
/data/simc
\end{simcshell}
Two more ways to choose data files:

\begin{itemize}
\item \code{-{}-{}data-{}dir\ DIR}\simcidx{Command-line options!data-dir@\code{-{}-{}data-{}dir}}: the directory of the spectral-function files of D/A(e,e'p), and the first place where a target-field
map (\hyperref[param:tgtu5ffieldu5ffile]{\code{tgt\_\allowbreak{}field\_\allowbreak{}file}}\simcidx{Parameters!tgt_field_file@\code{tgt\_field\_file}}) is looked for (then the current directory, then the data directory's \code{field/});
\item the JSON keys \code{"data\_\allowbreak{}dir"} (pion, kaon and SIDIS models) and \code{"matrix\_\allowbreak{}dir"} (spectrometers) of a pipeline stage
(\hyperref[user_manual-04u5finputs:43-the-pipeline-section-choosing-plug-ins]{4.3}).
\end{itemize}
\begin{simcnote}
a missing data file stops the run at start-up, not in the middle of a run.
\end{simcnote}
\section{\texorpdfstring{Running the tests}{Running the tests}}\label{user_manual-02u5finstallation:26-running-the-tests}
\begin{simcshell}[basicstyle=\ttfamily\footnotesize]
cd $SIMCPP_BUILD
ctest -L docs                  # the user-manual checks (seconds)
ctest -L unit -j8              # the unit tests (no external data needed)
ctest -j8                      # everything
\end{simcshell}
\begin{simctable}{small}{3}{>{\raggedright\arraybackslash}p{\dimexpr 0.191\linewidth-2\tabcolsep\relax}>{\raggedright\arraybackslash}p{\dimexpr 0.506\linewidth-2\tabcolsep\relax}>{\raggedright\arraybackslash}p{\dimexpr 0.235\linewidth-2\tabcolsep\relax}}
\caption{Running the tests}\\%
\toprule\rowcolor{simchead}\simcth{label} & \simcth{what it tests} & \simcth{needs the validation data} \\\midrule\endfirsthead
\multicolumn{3}{@{}l}{\simccontinued}\\\toprule\rowcolor{simchead}\simcth{label} & \simcth{what it tests} & \simcth{needs the validation data} \\\midrule\endhead
\midrule\multicolumn{3}{r@{}}{\simccontinues}\\\endfoot
\bottomrule\endlastfoot
\code{unit} & unit tests of the components & partly (many compare with reference tables) \\
\code{docs} & this manual: every input parameter and command-line option is documented, every link resolves & no \\
\code{examples}, \code{package} & the example plug-in; the installed CMake package used by an external project & yes (they run reference inputs) \\
\code{l1}, \code{replay} & each physics routine and complete runs against the original code's outputs & yes \\
\code{fortran\_\allowbreak{}diff} & each documented difference with the original code, in both modes & partly \\
\code{cli} & the command line (input formats, rejected inputs) & partly \\
\code{regression}, \code{threads} & outputs of standard runs unchanged from the stored baselines; \code{-{}-{}threads}\simcidx{Command-line options!threads@\code{-{}-{}threads}} reproducibility & yes \\
\code{regression\_\allowbreak{}shipped}, \code{slow} & complete runs of the example inputs, identical to the original program's outputs (long) & yes \\
\end{simctable}
Tests whose reference data, legacy source tree or baselines are missing report "Skipped", never "Failed": a build
configured without them runs the remaining tests (component self-tests, checksums of the shipped data files, the \code{docs}
tests) and skips the rest. To run the validation tests, export \code{SIMC\_\allowbreak{}REFERENCE}, \code{SIMC\_\allowbreak{}FORTRAN} and \code{SIMC\_\allowbreak{}BASELINES} before
configuring (or pass the \code{-{}D} options of \hyperref[user_manual-02u5finstallation:23-building]{2.3}).

The project script \code{tools/\allowbreak{}check\_\allowbreak{}all.\allowbreak{}sh\ [-{}-{}commit\ REV]\ [-{}-{}label\ NAME]} makes a clean build of SIMC++ and of the comparison
tool (with \code{ROOTSYS}, \code{SIMC\_\allowbreak{}REFERENCE}, \code{SIMC\_\allowbreak{}FORTRAN}, \code{SIMC\_\allowbreak{}BASELINES} from the environment),\allowbreak{} runs both test suites and logs the result in a run folder (about 75 minutes, most of it in the
\code{slow} tests).


% docs/manual/latex/user/03_quick_start.tex - chapter of the LaTeX edition of the book "user_manual" (SIMC++ documentation).
%   Converted once from docs/manual/markdown/user/03_quick_start.md by tools/md2latex.py; from then on THIS FILE is the source of the
%   LaTeX edition and may be edited by hand. Re-running tools/md2latex.py --write-sources for this
%   chapter OVERWRITES the hand edits (see docs/README.md).
%   Compile: cmake --build <build> --target pdf, or tools/build_docs_pdf.py (figures, xelatex passes,
%   makeindex with docs/common/latex/simcpp.ist); by hand, see docs/README.md. XeLaTeX, class
%   docs/common/latex/simcpp.cls (packages: fontspec, tikz, xcolor, hyperref, listings, longtable, booktabs,
%   titlesec, titletoc, fancyhdr, makeidx, preview, ...); fonts TeX Gyre Pagella, TeX Gyre Heros, DejaVu Sans
%   Mono. The figures are built first from docs/{manual,translation}/figures/*.tex (preview.sty) into PDFs.
\chapter{\texorpdfstring{Quick start}{Quick start}}\label{user_manual-03u5fquicku5fstart}
\simcabstract{The commands below were run as shown on the commit of this manual. The full transcript is in
\hyperref[user_manual-VERIFIEDu5fEXAMPLES]{VERIFIED\_\allowbreak{}EXAMPLES.md}. They assume the environment of
\hyperref[user_manual-02u5finstallation:21-directory-layout-used-in-this-manual]{2.1} and:}
\begin{simcshell}[basicstyle=\ttfamily\tiny]
export EX=$SIMCPP_PREFIX/share/simc/docs/manual/markdown/user/examples    # the example inputs of this manual
\end{simcshell}
\begin{simctip}
without installing, use \code{\$SIMCPP\_\allowbreak{}BUILD/\allowbreak{}bin/\allowbreak{}simc++} and \code{EX=\allowbreak{}\$SIMCPP\_\allowbreak{}SRC/\allowbreak{}docs/\allowbreak{}manual/\allowbreak{}markdown/\allowbreak{}user/\allowbreak{}examples}.
\end{simctip}
\section{\texorpdfstring{Run an input file}{Run an input file}}\label{user_manual-03u5fquicku5fstart:31-run-an-input-file}
\code{examples/\allowbreak{}h\_\allowbreak{}elastic\_\allowbreak{}rad.\allowbreak{}inp}\footnote{\code{docs/\allowbreak{}manual/\allowbreak{}markdown/\allowbreak{}user/\allowbreak{}examples/\allowbreak{}h\_\allowbreak{}elastic\_\allowbreak{}rad.\allowbreak{}inp}} is H(e,e'p) at Ebeam = 10.6 GeV with the SHMS on the electron
and the HMS on the proton, radiative corrections on, \code{random\_\allowbreak{}seed\ =\ 12345}, 200000 events requested. For a quick look, ask
for 5000:

\begin{simcshell}[basicstyle=\ttfamily\tiny]
$ simc++ $EX/h_elastic_rad.inp -o out --ngen 5000 > out.log; echo "exit status $?"
exit status 0
$ grep -E "fortran_compat|random seed" out.log
 Use random seed = 12345
 SIMC++: fortran_compat = false (clear legacy-SIMC bugs corrected, docs/translation/markdown/06_catalogue.md)
\end{simcshell}
The log shows the set-up messages and progress lines ("Generating Event ... successes so far - Monitor: ...") and ends
with \code{Come\ back\ soon!}. A run takes about 16 µs per try on one core for this input; here
102995 tries gave 5000 accepted events (\hyperref[user_manual-09u5fperformance]{chapter 9}).

\section{\texorpdfstring{What it writes}{What it writes}}\label{user_manual-03u5fquicku5fstart:32-what-it-writes}
\begin{simcshell}[basicstyle=\ttfamily\footnotesize]
$ ls out
h_elastic_rad.gen
h_elastic_rad.geni
h_elastic_rad.gen.root
h_elastic_rad.hist
h_elastic_rad.root
h_elastic_rad.summary.json
h_elastic_rad.summary.txt
\end{simcshell}
\begin{simctable}{small}{2}{>{\raggedright\arraybackslash}p{\dimexpr 0.174\linewidth-2\tabcolsep\relax}>{\raggedright\arraybackslash}p{\dimexpr 0.779\linewidth-2\tabcolsep\relax}}
%
\toprule\rowcolor{simchead}\simcth{file} & \simcth{content} \\\midrule\endfirsthead
\multicolumn{2}{@{}l}{\simccontinued}\\\toprule\rowcolor{simchead}\simcth{file} & \simcth{content} \\\midrule\endhead
\midrule\multicolumn{2}{r@{}}{\simccontinues}\\\endfoot
\bottomrule\endlastfoot
\code{.\allowbreak{}hist} & the run report: inputs, central event, normalisation (\code{wtcontr}, \code{normfac}, luminosity, \code{genvol}), limits \\
\code{.\allowbreak{}geni} & the spectrometer stop counters (where the tries were lost) \\
\code{.\allowbreak{}gen} & diagnostic histograms (text); \code{.\allowbreak{}gen.\allowbreak{}root} the same as ROOT TH1D \\
\code{.\allowbreak{}root} & the event tree \code{h10}, one entry per accepted event (written when \code{Nntu\ >\ 0}) \\
\code{.\allowbreak{}summary.\allowbreak{}json}, \code{.\allowbreak{}summary.\allowbreak{}txt} & the run summary in machine-readable form (the input of \code{simc\_\allowbreak{}compare}) \\
\end{simctable}
The normalisation is in the \code{.\allowbreak{}hist}:

\begin{simcshell}[basicstyle=\ttfamily\footnotesize]
$ grep -E "^ *(Ngen|Ntried|Ncontribute)|wtcontr|normfac" out/h_elastic_rad.hist
            Ngen (request) =       5000
            Ntried         =     102995
            Ncontribute    =       5000
               MeV: wtcontr=   0.20906661E+03
                   normfac =     0.407176E+06
\end{simcshell}
\hyperref[user_manual-07u5foutputs]{Chapter 7} explains every output and every branch.

\section{\texorpdfstring{Look at the tree in ROOT}{Look at the tree in ROOT}}\label{user_manual-03u5fquicku5fstart:33-look-at-the-tree-in-root}
The \code{Weight} branch is the event weight. \code{Weight} × normfac / nevent is the yield of one event for the
simulated charge (nevent = Ncontribute when ngen > 0) (\hyperref[user_manual-07u5foutputs:711-weights-and-normalisation]{7.1.1}). A small macro:

\begin{simccode}[basicstyle=\ttfamily\tiny,style=simccpp]
// peek_h10.C: open a SIMC++ tree, print its size and the weighted Em distribution
void peek_h10(const char* file = "out/h_elastic_rad.root") {
  TFile f(file);
  TTree* t = (TTree*)f.Get("h10");
  std::cout << t->GetEntries() << " events, " << t->GetListOfBranches()->GetEntries() << " branches" << std::endl;
  t->Scan("hsdelta:ssdelta:W:Em:Weight", "", "", 5);
  TH1D h("h", "E_{m} (weighted);E_{m} [GeV]", 100, -0.02, 0.2);
  t->Project("h", "Em", "Weight");
  std::cout << "sum of weights in [-0.02, 0.2] GeV: " << h.Integral() << ", mean Em " << h.GetMean() << " GeV" << std::endl;
}
\end{simccode}
\begin{simcshell}[basicstyle=\ttfamily\footnotesize]
$ root -l -b -q peek_h10.C
Processing peek_h10.C...
5000 events, 46 branches
************************************************************************
*    Row   *   hsdelta *   ssdelta *         W *        Em *    Weight *
************************************************************************
*        0 * 6.8000464 * -2.967325 * 1.0069015 * 0.0488890 * 0.0002927 *
*        1 * 2.4381718 * -0.939918 * 0.9532427 * 0.0063532 * 0.0004068 *
*        2 * -2.744465 * -9.402965 * 1.4675965 * 0.8642277 * 0.0004998 *
*        3 * -1.449645 * 0.4658077 * 0.9453506 * 0.0016173 * 0.0005580 *
*        4 * 0.9987937 * -1.216251 * 0.9978361 * 0.0712616 * 0.0004664 *
************************************************************************
sum of weights in [-0.02, 0.2] GeV: 2.24626, mean Em 0.0176472 GeV
\end{simcshell}
Here the electron is in the SHMS (\code{electron\_\allowbreak{}arm\ =\ 5}, a beam-left spectrometer),\allowbreak{} so \code{hs*} holds the proton (HMS) and \code{ss*}
the electron (\hyperref[user_manual-07u5foutputs:712-arm-prefixes-hs-and-ss]{7.1.2}). Event 2 has radiated: W = 1.47 GeV, Em = 0.86 GeV.

\section{\texorpdfstring{The same run from JSON}{The same run from JSON}}\label{user_manual-03u5fquicku5fstart:34-the-same-run-from-json}
\begin{simcshell}[basicstyle=\ttfamily\tiny]
$ inp2json -o run.json $EX/h_elastic_rad.inp
$ simc++ run.json -o out_json --ngen 5000 > out_json.log && diff <(sed 2,3d out/h_elastic_rad.hist) <(sed 2,3d out_json/run.hist) && echo ".hist identical (apart from the time lines)"
.hist identical (apart from the time lines)
$ simc_root_identity out/h_elastic_rad.root out_json/run.root
  h10                      TTree  5000 entries x 46 leaves: identical
out/h_elastic_rad.root vs out_json/run.root: IDENTICAL
\end{simcshell}
(The output names follow the input file name, here \code{run}; use \code{-{}-{}run-{}name}\simcidx{Command-line options!run-name@\code{-{}-{}run-{}name}} or \code{"run":\ \{"name":\ .\allowbreak{}.\allowbreak{}.\allowbreak{}\}} to change them.)
Edit \code{run.\allowbreak{}json} to change parameters, add a \code{pipeline} section to change plug-ins, or add \code{"run":\ \{"fortran\_\allowbreak{}compat":\ true\}} for the legacy-compatible mode (\hyperref[user_manual-04u5finputs]{chapter 4}).

The same run is also shipped as a ready JSON configuration,\allowbreak{} \code{\$SIMCPP\_\allowbreak{}PREFIX/\allowbreak{}share/\allowbreak{}simc/\allowbreak{}data/\allowbreak{}config/\allowbreak{}h\_\allowbreak{}elastic\_\allowbreak{}rad.\allowbreak{}json}, with the
plug-in choice written out; \code{simc++\ \$SIMCPP\_\allowbreak{}PREFIX/\allowbreak{}share/\allowbreak{}simc/\allowbreak{}data/\allowbreak{}config/\allowbreak{}h\_\allowbreak{}elastic\_\allowbreak{}rad.\allowbreak{}json\ -{}o\ out\_\allowbreak{}cfg\ -{}-{}ngen\ 5000} runs it.

\section{\texorpdfstring{Next steps}{Next steps}}\label{user_manual-03u5fquicku5fstart:35-next-steps}
\begin{itemize}
\item More events, faster: \code{-{}-{}threads\ N}\simcidx{Command-line options!threads@\code{-{}-{}threads}} (\hyperref[user_manual-08u5frunningu5fandu5freproducibility]{chapter 8}, \hyperref[user_manual-09u5fperformance]{chapter 9}).
\item Reproduce results of the legacy SIMC exactly: \hyperref[user_manual-08u5frunningu5fandu5freproducibility:84-reproducing-legacy-simc-results-exactly]{8.4}.
\item An old input stops with \code{bad\ targ.\allowbreak{}can\ value}: \hyperref[user_manual-04u5finputs:413-old--names]{4.1.3}.
\end{itemize}


\simcpart{Configuring a simulation}
  {The input files, every input parameter, the reactions and the spectrometers.}
% docs/manual/latex/user/04_inputs.tex - chapter of the LaTeX edition of the book "user_manual" (SIMC++ documentation).
%   Converted once from docs/manual/markdown/user/04_inputs.md by tools/md2latex.py; from then on THIS FILE is the source of the
%   LaTeX edition and may be edited by hand. Re-running tools/md2latex.py --write-sources for this
%   chapter OVERWRITES the hand edits (see docs/README.md).
%   Compile: cmake --build <build> --target pdf, or tools/build_docs_pdf.py (figures, xelatex passes,
%   makeindex with docs/common/latex/simcpp.ist); by hand, see docs/README.md. XeLaTeX, class
%   docs/common/latex/simcpp.cls (packages: fontspec, tikz, xcolor, hyperref, listings, longtable, booktabs,
%   titlesec, titletoc, fancyhdr, makeidx, preview, ...); fonts TeX Gyre Pagella, TeX Gyre Heros, DejaVu Sans
%   Mono. The figures are built first from docs/{manual,translation}/figures/*.tex (preview.sty) into PDFs.
\chapter{\texorpdfstring{Inputs}{Inputs}}\label{user_manual-04u5finputs}
\simcabstract{\code{simc++} takes one input file, a \code{.\allowbreak{}inp} or a \code{.\allowbreak{}json} (chosen by the file extension):}
\begin{simcplain}[basicstyle=\ttfamily\footnotesize]
simc++ <file.inp | file.json> [options]
\end{simcplain}
Both formats set the same 106 input variables, with the same names, units and defaults (\hyperref[user_manual-04u5finputs:46-parameter-reference]{4.6}).
The JSON format can also choose the plug-ins and the random-number generator, and set the run options.

\section{\texorpdfstring{The \code{.\allowbreak{}inp} format}{The .inp format}}\label{user_manual-04u5finputs:41-the-inp-format}
\subsection{\texorpdfstring{Syntax}{Syntax}}\label{user_manual-04u5finputs:411-syntax}
\code{.\allowbreak{}inp} files are made of named blocks of \code{name\ =\ value} lines (the "CTP parm" syntax of SIMC input files):

\begin{simcplain}[basicstyle=\ttfamily\footnotesize]
; comment
begin parm experiment
  ngen = 50000			;  POS: # of successes; NEG: # of tries
  EXPER%charge = 1.0		;  total charge (mC)
  doing_pion = 1
  extra_dbase_file = 'extra_nucpi'   ; more variables in extra_nucpi.inp
end parm experiment
begin parm debug
  debug(1) = 0
end parm debug
\end{simcplain}
\begin{itemize}
\item Only the lines between \code{begin\ parm\ <name>} and \code{end\ parm} are read. The block name does not matter.
\item \code{;} starts a comment, except inside quotes. Tabs count as spaces.
\item Variable names are \textbf{case-insensitive} (\code{ebeam} = \hyperref[param:Ebeam]{\code{Ebeam}}\simcidx{Parameters!Ebeam@\code{Ebeam}}). Structure members are separated by \code{\%} (\hyperref[param:targu25thick]{\code{targ\%\allowbreak{}thick}}\simcidx{Parameters!targ.thick@\code{targ\%thick}}).
\item \code{name(\allowbreak{}i)} is 1-based, \code{name[i]} 0-based. \code{name\ =\ a,\ b,\ c} fills consecutive array elements.
\item Values may be expressions (\code{+\ -{}\ *\ /\ (\ )}). A real value given to an integer variable is rounded.
\item Flags such as \hyperref[param:doingu5fpion]{\code{doing\_\allowbreak{}pion}}\simcidx{Parameters!doing_pion@\code{doing\_pion}} are integers in the file; a value > 0 means true.
\item \textbf{Every real value is stored as a 32-bit float first.} For example \code{ctau\ =\ 780.\allowbreak{}4} becomes 780.4000244140625. This keeps
the inputs identical to those of the legacy program. \code{inp2json} writes the rounded values.
\end{itemize}
The example inputs of this manual are in examples/\footnote{\code{docs/\allowbreak{}manual/\allowbreak{}markdown/\allowbreak{}user/\allowbreak{}examples/}} (installed in \code{\$SIMCPP\_\allowbreak{}PREFIX/\allowbreak{}share/\allowbreak{}simc/\allowbreak{}docs/\allowbreak{}manual/\allowbreak{}markdown/\allowbreak{}user/\allowbreak{}examples/}).

\subsection{\texorpdfstring{\code{extra\_\allowbreak{}dbase\_\allowbreak{}file}}{extra\_\allowbreak{}dbase\_\allowbreak{}file}}\label{user_manual-04u5finputs:412-extrau5fdbaseu5ffile}
\code{extra\_\allowbreak{}dbase\_\allowbreak{}file\ =\ 'name'} loads a second file after the main one; its values override the main file's. \code{.\allowbreak{}inp} is
added when the name has no \code{.}. The file is looked for in the \textbf{directory of the main input file}, or in the directory
given with \code{-{}-{}infiles-{}dir\ DIR}\simcidx{Command-line options!infiles-dir@\code{-{}-{}infiles-{}dir}}.

\subsection{\texorpdfstring{Old \code{.} names}{Old . names}}\label{user_manual-04u5finputs:413-old--names}
Old input files write structure names with \code{.} instead of \code{\%} (\code{targ.\allowbreak{}A}, \code{spec.\allowbreak{}e.\allowbreak{}P}, \code{SPedge.\allowbreak{}e.\allowbreak{}delta.\allowbreak{}min}, ...). These
names are \textbf{not} input variables any more: each one is reported, its value is ignored, and the run uses the defaults,
which usually makes it stop. SIMC++ lists the ignored names. With the example examples/\allowbreak{}legacy\_\allowbreak{}names.\allowbreak{}inp\footnote{\code{docs/\allowbreak{}manual/\allowbreak{}markdown/\allowbreak{}user/\allowbreak{}examples/\allowbreak{}legacy\_\allowbreak{}names.\allowbreak{}inp}}:

\begin{simcshell}[basicstyle=\ttfamily\tiny]
$ simc++ $EX/legacy_names.inp -o legacy > legacy.log 2> legacy.err; echo "exit status $?"; cat legacy.err
exit status 3
simc++: WARNING: 46 input variable(s) are not SIMC variables and are IGNORED (legacy SIMC behaviour; their defaults are used):
  EXPER.charge, spec.e.P, spec.e.theta, spec.p.P, spec.p.theta, targ.A
  ...
  Legacy '.' names (targ.A, spec.e.P, ...) must be converted to '%' names with simc_gfortran/infiles/convert_inputfile.sh.
simc++: STOP during setup (legacy SIMC behaviour): STOP bad targ.can value
simc++: note: the input has 46 ignored variable(s) (listed above), so their defaults were used — the likely cause of this STOP. ...
\end{simcshell}
Convert such a file by replacing the \code{.} of these names with \code{\%} (keep a copy of the original). This \code{sed} command does
it for every structure name of the inputs (\code{targ}, \code{spec}, \code{SPedge}, \code{gen}, \code{cuts}, \code{EXPER}):

\begin{simcshell}[basicstyle=\ttfamily\tiny]
sed -E ':a; s/^(\s*(targ|spec|SPedge|gen|cuts|EXPER)[A-Za-z_%]*)\./\1%/; ta' legacy_names.inp > converted.inp
\end{simcshell}
Convert the \hyperref[param:extrau5fdbaseu5ffile]{\code{extra\_\allowbreak{}dbase\_\allowbreak{}file}}\simcidx{Parameters!extra_dbase_file@\code{extra\_dbase\_file}} too. (The SIMC distribution has an equivalent script, \code{infiles/\allowbreak{}convert\_\allowbreak{}inputfile.\allowbreak{}sh}.)
The converted file runs normally (\code{simc++\ converted.\allowbreak{}inp}).

\subsection{\texorpdfstring{Unknown names}{Unknown names}}\label{user_manual-04u5finputs:414-unknown-names}
Any name that is not one of the 106 variables, such as a typo or an old name, is reported once per name
(\code{<name>\ not\ registered,\ registering\ as\ int}) and ignored. \code{inp2json\ -{}-{}strict} turns these names into an error (exit
status 3).

\section{\texorpdfstring{The JSON format}{The JSON format}}\label{user_manual-04u5finputs:42-the-json-format}
A SIMC++ JSON input has up to three sections:

\begin{simccode}[basicstyle=\ttfamily\footnotesize,style=simcjson]
{
  "parameters": { ... the input variables ... },
  "pipeline":   { ... optional: plug-ins, random-number generator, threads ... },
  "run":        { "name": "myrun", "output_dir": "out", "fortran_compat": false }
}
\end{simccode}
\subsection{\texorpdfstring{\code{parameters}}{parameters}}\label{user_manual-04u5finputs:421-parameters}
\begin{itemize}
\item The key path is the variable name split on \code{\%}: \hyperref[param:targu25thick]{\code{targ\%\allowbreak{}thick}}\simcidx{Parameters!targ.thick@\code{targ\%thick}} becomes \code{\{"targ":\ \{"thick":\ 723.\allowbreak{}0\}\}}, and
\hyperref[param:SPedgeu25eu25deltau25min]{\code{SPedge\%\allowbreak{}e\%\allowbreak{}delta\%\allowbreak{}min}}\simcidx{Parameters!SPedge.e.delta.min@\code{SPedge\%e\%delta\%min}} becomes \code{\{"SPedge":\ \{"e":\ \{"delta":\ \{"min":\ -{}10.\allowbreak{}0\}\}\}\}}. Arrays are JSON arrays
(\code{"lambda":\ [0,\ 0,\ 0]}, \code{"debug":\ [false,\ false,\ false,\ false,\ false,\ false]}).
\item Values are in the \textbf{\code{.\allowbreak{}inp} units} (MeV, MeV/c, degrees, mrad, mg/cm², \%). \code{inp2json\ -{}-{}table} prints the values after
the internal unit conversions.
\item Flags accept \code{true}/\code{false} or numbers (> 0 is true).
\item Keys are matched case-insensitively. Keys that start with \code{\_} (for example \code{"\_\allowbreak{}comment"}) are ignored. Unknown keys
produce \code{simc++:\ warning:\ unknown\ parameter\ '<key>\allowbreak{}'\ .\allowbreak{}.\allowbreak{}.\allowbreak{},\ ignored}.
\item A missing key keeps the default (\hyperref[user_manual-04u5finputs:46-parameter-reference]{4.6}).
\end{itemize}
The easiest way to write a JSON input is to convert a \code{.\allowbreak{}inp}:

\begin{simcshell}[basicstyle=\ttfamily\footnotesize]
inp2json -o run.json $EX/h_elastic_rad.inp
simc++ run.json -o out_json
\end{simcshell}
The converted file holds all 106 variables (with \hyperref[param:extrau5fdbaseu5ffile]{\code{extra\_\allowbreak{}dbase\_\allowbreak{}file}}\simcidx{Parameters!extra_dbase_file@\code{extra\_dbase\_file}} already merged) and a \code{\_\allowbreak{}source} block that records
the converted files and any names that were not recognised. A \code{.\allowbreak{}inp} and its conversion give identical outputs
(\hyperref[user_manual-03u5fquicku5fstart:34-the-same-run-from-json]{3.4}; CTest \code{cli\_\allowbreak{}json\_\allowbreak{}vs\_\allowbreak{}inp\_\allowbreak{}*}).

Other \code{inp2json} options: \code{-{}-{}infiles-{}dir\ DIR}\simcidx{Command-line options!infiles-dir@\code{-{}-{}infiles-{}dir}} (where to find \hyperref[param:extrau5fdbaseu5ffile]{\code{extra\_\allowbreak{}dbase\_\allowbreak{}file}}\simcidx{Parameters!extra_dbase_file@\code{extra\_dbase\_file}}), \code{-{}-{}no-{}extra} (do not load it),
\code{-{}-{}table[=\allowbreak{}FILE]} (also print the converted values), \code{-{}-{}mapping} (print the \code{.\allowbreak{}inp} → JSON table of all variables), \code{-{}-{}strict},
\code{-{}q}.

Example configurations come with SIMC++ in \code{data/\allowbreak{}config/} (installed in \code{\$SIMCPP\_\allowbreak{}PREFIX/\allowbreak{}share/\allowbreak{}simc/\allowbreak{}data/\allowbreak{}config/}):
\code{h\_\allowbreak{}elastic\_\allowbreak{}norad.\allowbreak{}json}, \code{h\_\allowbreak{}elastic\_\allowbreak{}rad.\allowbreak{}json} and \code{acceptance\_\allowbreak{}phsp\_\allowbreak{}hms\_\allowbreak{}shms.\allowbreak{}json}, with the three sections written out.

\subsection{\texorpdfstring{\code{run}}{run}}\label{user_manual-04u5finputs:422-run}
\begin{simctable}{small}{4}{>{\raggedright\arraybackslash}p{\dimexpr 0.153\linewidth-2\tabcolsep\relax}>{\raggedright\arraybackslash}p{\dimexpr 0.268\linewidth-2\tabcolsep\relax}>{\raggedright\arraybackslash}p{\dimexpr 0.163\linewidth-2\tabcolsep\relax}>{\raggedright\arraybackslash}p{\dimexpr 0.328\linewidth-2\tabcolsep\relax}}
%
\toprule\rowcolor{simchead}\simcth{key} & \simcth{command-line equivalent} & \simcth{default} & \simcth{meaning} \\\midrule\endfirsthead
\multicolumn{4}{@{}l}{\simccontinued}\\\toprule\rowcolor{simchead}\simcth{key} & \simcth{command-line equivalent} & \simcth{default} & \simcth{meaning} \\\midrule\endhead
\midrule\multicolumn{4}{r@{}}{\simccontinues}\\\endfoot
\bottomrule\endlastfoot
\code{name} & \code{-{}-{}run-{}name\ NAME}\simcidx{Command-line options!run-name@\code{-{}-{}run-{}name}} & the input file name without its extension & stem of the output files \code{<name>\allowbreak{}.\allowbreak{}root}, \code{<name>\allowbreak{}.\allowbreak{}hist}, ... \\
\code{output\_\allowbreak{}dir} & \code{-{}o\ DIR}, \code{-{}-{}output-{}dir\ DIR} & \code{.} & output directory (created if needed) \\
\code{fortran\_\allowbreak{}compat} & \code{-{}-{}fortran-{}compat}\simcidx{Command-line options!fortran-compat@\code{-{}-{}fortran-{}compat}} & \code{false} & \code{true}: the legacy-compatible mode (\hyperref[user_manual-08u5frunningu5fandu5freproducibility:83-default-and-legacy-compatible-mode]{8.3}) \\
\end{simctable}
A command-line option wins over the JSON value. A \code{.\allowbreak{}inp} input runs in the default mode unless you give \code{-{}-{}fortran-{}compat}\simcidx{Command-line options!fortran-compat@\code{-{}-{}fortran-{}compat}}.

\section{\texorpdfstring{The \code{pipeline} section: choosing plug-ins}{The pipeline section: choosing plug-ins}}\label{user_manual-04u5finputs:43-the-pipeline-section-choosing-plug-ins}
Without a \code{pipeline} section, SIMC++ chooses the components from the input flags. \code{-{}-{}print-{}pipeline}\simcidx{Command-line options!print-pipeline@\code{-{}-{}print-{}pipeline}} prints the resolved
choice when the run starts (the run then continues). Each entry is a plug-in name, or an object
\code{\{"type":\ "<name>\allowbreak{}",\ <options>\allowbreak{}\}}:

\begin{simctable}{small}{3}{>{\raggedright\arraybackslash}p{\dimexpr 0.188\linewidth-2\tabcolsep\relax}>{\raggedright\arraybackslash}p{\dimexpr 0.256\linewidth-2\tabcolsep\relax}>{\raggedright\arraybackslash}p{\dimexpr 0.488\linewidth-2\tabcolsep\relax}}
\caption{The \code{pipeline} section: choosing plug-ins}\\%
\toprule\rowcolor{simchead}\simcth{key} & \simcth{default (from the flags)} & \simcth{registered plug-ins} \\\midrule\endfirsthead
\multicolumn{3}{@{}l}{\simccontinued}\\\toprule\rowcolor{simchead}\simcth{key} & \simcth{default (from the flags)} & \simcth{registered plug-ins} \\\midrule\endhead
\midrule\multicolumn{3}{r@{}}{\simccontinues}\\\endfoot
\bottomrule\endlastfoot
\code{rng} & \code{mt64}, seeded with \hyperref[param:randomu5fseed]{\code{random\_\allowbreak{}seed}}\simcidx{Parameters!random_seed@\code{random\_seed}} & \code{mt64} (std::\allowbreak{}mt19937\_\allowbreak{}64),\allowbreak{} \code{ranlux} (the legacy SIMC generator); options \code{"seed"}, and for \code{ranlux} \code{"state"} (25 integers) (\hyperref[user_manual-08u5frunningu5fandu5freproducibility]{chapter 8}) \\
\code{target} & \code{hallc\_\allowbreak{}cryo} & \code{hallc\_\allowbreak{}cryo} (cryotarget and solid targets: windows, energy loss, multiple scattering) \\
\code{beam} & \code{standard} & \code{standard} \\
\code{reaction} & from the flags (\hyperref[user_manual-05u5freactions]{chapter 5}) & \code{hydrogen\_\allowbreak{}elastic}, \code{deuterium\_\allowbreak{}eep}, \code{nucleus\_\allowbreak{}eep}, \code{pion\_\allowbreak{}production}, \code{kaon\_\allowbreak{}production}, \code{delta\_\allowbreak{}production}, \code{sidis\_\allowbreak{}production}, \code{rho\_\allowbreak{}production}, \code{phase\_\allowbreak{}space} \\
\code{physics} & from the flags & \code{sigep}, \code{deforest} (and \code{deforest\_\allowbreak{}cc1}, \code{deforest\_\allowbreak{}cc1\_\allowbreak{}onshell}, \code{deforest\_\allowbreak{}cc2}), \code{pion}, \code{kaon}, \code{delta}, \code{sidis}, \code{rho} \\
\code{radiator} & \code{simc\_\allowbreak{}rad} & \code{simc\_\allowbreak{}rad} \\
\code{spectrometers} & \code{\{"electron":\ .\allowbreak{}.\allowbreak{}.\allowbreak{},\ "hadron":\ .\allowbreak{}.\allowbreak{}.\allowbreak{}\}} from \hyperref[param:electronu5farm]{\code{electron\_\allowbreak{}arm}}\simcidx{Parameters!electron_arm@\code{electron\_arm}}/\hyperref[param:hadronu5farm]{\code{hadron\_\allowbreak{}arm}}\simcidx{Parameters!hadron_arm@\code{hadron\_arm}} & \code{hms}, \code{shms}, \code{sos}, \code{hrsr}, \code{hrsl}, \code{calo}; the example \code{ideal\_\allowbreak{}box} (\hyperref[user_manual-06u5fspectrometers]{chapter 6}) \\
\code{sinks} & \code{root\_\allowbreak{}tree} (only if \code{Nntu\ >\ 0}), \code{run\_\allowbreak{}summary}, \code{gen\_\allowbreak{}hist}, \code{fortran\_\allowbreak{}hist} & the list \textbf{replaces} the default list (\hyperref[user_manual-07u5foutputs]{chapter 7}) \\
\code{threads} & 1 & number of event-loop threads; \code{-{}-{}threads\ N}\simcidx{Command-line options!threads@\code{-{}-{}threads}} wins (\hyperref[user_manual-08u5frunningu5fandu5freproducibility]{chapter 8}) \\
\end{simctable}
Options of the built-in plug-ins (all optional):

\begin{simctable}{small}{2}{>{\raggedright\arraybackslash}p{\dimexpr 0.212\linewidth-2\tabcolsep\relax}>{\raggedright\arraybackslash}p{\dimexpr 0.741\linewidth-2\tabcolsep\relax}}
%
\toprule\rowcolor{simchead}\simcth{plug-in} & \simcth{options} \\\midrule\endfirsthead
\multicolumn{2}{@{}l}{\simccontinued}\\\toprule\rowcolor{simchead}\simcth{plug-in} & \simcth{options} \\\midrule\endhead
\midrule\multicolumn{2}{r@{}}{\simccontinues}\\\endfoot
\bottomrule\endlastfoot
all stages & \code{"verbose":\ false} silences the stage messages (\code{-{}-{}quiet}\simcidx{Command-line options!quiet@\code{-{}-{}quiet}} sets it for every stage) \\
\code{reaction} of D/A(e,e'p) & \code{"spectral\_\allowbreak{}function":\ \{"type":\ "theory"\ |\ "benhar",\ "file":\ "<file\ in\ -{}-{}data-{}dir>\allowbreak{}",\ "transparency":\ x\}} (\hyperref[user_manual-05u5freactions:52-heep-elastic-deep-and-aeep]{5.2}) \\
\code{pion}, \code{kaon}, \code{sidis} models & \code{"data\_\allowbreak{}dir"}: directory of the model tables (default: the data directory's \code{pion/}, \code{kaon/}, \code{sidis/}) \\
\code{hms}, \code{shms}, \code{sos}, \code{hrsr}, \code{hrsl} & \code{"matrix\_\allowbreak{}dir"} (optics maps), and expert overrides of the per-arm switches: \code{"ms\_\allowbreak{}flag"}, \code{"wcs\_\allowbreak{}flag"} (multiple scattering and wire-chamber smearing, default \hyperref[param:mcu5fsmear]{\code{mc\_\allowbreak{}smear}}\simcidx{Parameters!mc_smear@\code{mc\_smear}}), \code{"decay"} (default \hyperref[param:doingu5fdecay]{\code{doing\_\allowbreak{}decay}}\simcidx{Parameters!doing_decay@\code{doing\_decay}} for the hadron arm), \code{"ctau"}, \code{"use\_\allowbreak{}sieve"}; \code{hms}: \code{"using\_\allowbreak{}HMScoll"}; \code{shms}: \code{"using\_\allowbreak{}SHMScoll"}, \code{"use\_\allowbreak{}coll"}, \code{"skip\_\allowbreak{}hb"}, \code{"use\_\allowbreak{}2sieve"} \\
\code{calo} & \code{"drift\_\allowbreak{}to\_\allowbreak{}cal"}, \code{"ctau"} \\
sinks & \code{root\_\allowbreak{}tree}: \code{"file"}, \code{"tree"}; \code{run\_\allowbreak{}summary}, \code{gen\_\allowbreak{}hist}: \code{"file"}; \code{fortran\_\allowbreak{}hist}: \code{"hist"}, \code{"geni"}, \code{"gen"} (file names) \\
\end{simctable}
An example: transport the hadron through the ideal box acceptance of the example plug-in instead of the real
spectrometer. \code{-{}-{}pipeline\ FILE}\simcidx{Command-line options!pipeline@\code{-{}-{}pipeline}} merges a pipeline file over the input's own pipeline:

\begin{simcshell}[basicstyle=\ttfamily\tiny]
$ simc++ run.json -o out_ideal --ngen 200 --pipeline $SIMCPP_PREFIX/share/simc/examples/ideal_box/hadron_arm.json --print-pipeline --quiet
  ...
  "spectrometers": {
    "electron": { "type": "" },
    "hadron": { "delta_max": 8.0, "delta_min": -8.0, "type": "ideal_box", "xptar_max": 0.05, ... }
  },
\end{simcshell}
(\code{"type":\ ""} means "the default for this arm code".) Two arms with the same plug-in and the same options share one
spectrometer object, so they share one set of stop counters.

\section{\texorpdfstring{Run-time messages about the inputs}{Run-time messages about the inputs}}\label{user_manual-04u5finputs:44-run-time-messages-about-the-inputs}
\begin{simctable}{small}{2}{>{\raggedright\arraybackslash}p{\dimexpr 0.492\linewidth-2\tabcolsep\relax}>{\raggedright\arraybackslash}p{\dimexpr 0.462\linewidth-2\tabcolsep\relax}}
\caption{Run-time messages about the inputs}\\%
\toprule\rowcolor{simchead}\simcth{message} & \simcth{meaning} \\\midrule\endfirsthead
\multicolumn{2}{@{}l}{\simccontinued}\\\toprule\rowcolor{simchead}\simcth{message} & \simcth{meaning} \\\midrule\endhead
\midrule\multicolumn{2}{r@{}}{\simccontinues}\\\endfoot
\bottomrule\endlastfoot
\code{<name>\ not\ registered,\ registering\ as\ int} & \code{.\allowbreak{}inp} name unknown: ignored \\
\code{simc++:\ WARNING:\ N\ input\ variable(\allowbreak{}s)\ are\ not\ SIMC\ variables\ and\ are\ IGNORED\ .\allowbreak{}.\allowbreak{}.} & summary of the above, with the hint to convert old names \\
\code{simc++:\ warning:\ unknown\ parameter\ '<key>\allowbreak{}'\ .\allowbreak{}.\allowbreak{}.} & JSON key unknown: ignored \\
\code{WARNING:\ forcing\ target\ mass\ to\ Mp!!!} and similar & the set-up adjusted an input (for example the H/D masses) \\
\code{NOTE:\ random\_\allowbreak{}state\_\allowbreak{}file\ '.\allowbreak{}.\allowbreak{}.\allowbreak{}'\ is\ not\ used\ by\ SIMC++} & see \hyperref[param:randomu5fstateu5ffile]{\code{random\_\allowbreak{}state\_\allowbreak{}file}}\simcidx{Parameters!random_state_file@\code{random\_state\_file}} below \\
a WARNING when \hyperref[param:ptu5fbu5fparam]{\code{pt\_\allowbreak{}b\_\allowbreak{}param}}\simcidx{Parameters!pt_b_param@\code{pt\_b\_param}} or \code{sigc\_\allowbreak{}*} are set & these SIDIS inputs are read and ignored \\
a WARNING for \hyperref[param:spectu5fmode]{\code{spect\_\allowbreak{}mode}}\simcidx{Parameters!spect_mode@\code{spect\_mode}} outside \{0, 1, −1, −2\} & any other value means "both arms on" (\hyperref[user_manual-06u5fspectrometers:64-spectu5fmode-switching-the-spectrometer-monte-carlo-off]{6.4}) \\
\code{simc++:\ input\ error:\ .\allowbreak{}.\allowbreak{}.} (exit 2) & a rejected input, for example arm code 6 \\
\code{simc++:\ STOP\ during\ setup\ .\allowbreak{}.\allowbreak{}.} (exit 3) & an input check of the set-up failed (the message comes from the original code; \hyperref[user_manual-10u5ffaq:simc-stop-during-setup-legacy-simc-behaviour-stop--exit-3]{10} lists the common ones) \\
\end{simctable}
\section{\texorpdfstring{Units and conventions}{Units and conventions}}\label{user_manual-04u5finputs:45-units-and-conventions}
\begin{itemize}
\item Energies and momenta in MeV and MeV/c; angles in degrees (spectrometer angles, target angle, field angles);
spectrometer slopes (\code{xptar}, \code{yptar}) in mrad; delta in \%; lengths in cm; target thickness in mg/cm²; density in
g/cm³; charge in mC.
\item Spectrometer angles are taken as their absolute value. The side of the beam is fixed by the spectrometer number
(\hyperref[user_manual-06u5fspectrometers:61-spectrometer-numbers]{6.1}).
\item Target coordinates: +x = beam right, +y = up, +z = downstream.
\item Defaults: 0 / false / blank, except the \code{SPedge\%\allowbreak{}*} limits and \code{cuts\%\allowbreak{}Em} (±1e10).
\end{itemize}
\section{\texorpdfstring{Parameter reference}{Parameter reference}}\label{user_manual-04u5finputs:46-parameter-reference}
All 106 input variables, grouped by topic. \textbf{Name} is the \code{.\allowbreak{}inp} name. The JSON path is the same name with \code{\%} replaced
by nested objects (\hyperref[user_manual-04u5finputs:421-parameters]{4.2.1}). \textbf{Default} is the value used when the variable is absent. \textbf{Used by}
says which runs read it: \emph{all}, \emph{eep} (H/\allowbreak{}D/\allowbreak{}A(e,\allowbreak{}e'p)),\allowbreak{} \emph{D/A eep}, \emph{pion}, \emph{kaon}, \emph{delta}, \emph{semi} (SIDIS), \emph{rho},
\emph{phsp}, \emph{rad} (with \hyperref[param:usingu5frad]{\code{using\_\allowbreak{}rad}}\simcidx{Parameters!using_rad@\code{using\_rad}}), \emph{calo} (arms 7/8), \emph{field} (\hyperref[param:usingu5ftgtu5ffield]{\code{using\_\allowbreak{}tgt\_\allowbreak{}field}}\simcidx{Parameters!using_tgt_field@\code{using\_tgt\_field}}). The test
\code{docs\_\allowbreak{}user\_\allowbreak{}manual\_\allowbreak{}parameters} fails if an input variable of SIMC++ is missing here.

\subsection*{\texorpdfstring{Run control}{Run control}}\label{user_manual-04u5finputs:run-control}
\begin{simcparams}
\simcparam{ngen}{param:ngen}{events}{0}{all}{> 0: number of \textbf{accepted} events to generate; < 0: number of \textbf{tries} (|ngen|). 0 writes nothing. \code{-{}-{}ngen\ N}\simcidx{Command-line options!ngen@\code{-{}-{}ngen}} overrides it. With ngen < 0 the \code{.\allowbreak{}hist} wtcontr and AVERAGE.sigcc are divided by the tries}{ngen@\code{ngen}}
\simcparam{random\_\allowbreak{}seed}{param:randomu5fseed}{integer}{0}{all}{seed of the random-number generator. −1: the current time. For \code{mt64} a value ≤ 0 means the fixed seed 5489; for \code{ranlux} it seeds RANLUX at luxury level 3 (0 = RANLUX's default seed). \code{-{}-{}seed\ N}\simcidx{Command-line options!seed@\code{-{}-{}seed}} overrides it}{random_seed@\code{random\_seed}}
\simcparam{random\_\allowbreak{}state\_\allowbreak{}file}{param:randomu5fstateu5ffile}{path}{blank}{—}{\textbf{Not used} (a note is printed); to start RANLUX from a saved state use \code{"rng":\ \{"type":\ "ranlux",\ "state":\ [25\ integers]\}}}{random_state_file@\code{random\_state\_file}}
\simcparam{extra\_\allowbreak{}dbase\_\allowbreak{}file}{param:extrau5fdbaseu5ffile}{file stem}{blank}{all}{second input file, loaded after the main one (\hyperref[user_manual-04u5finputs:412-extrau5fdbaseu5ffile]{4.1.2})}{extra_dbase_file@\code{extra\_dbase\_file}}
\simcparam{Nntu}{param:Nntu}{flag}{0}{all}{> 0: write the event tree \code{<run>\allowbreak{}.\allowbreak{}root}. Without it no \code{.\allowbreak{}root} file is written}{Nntu@\code{Nntu}}
\simcparam{EXPER\%\allowbreak{}charge}{param:EXPERu25charge}{mC}{0}{all}{beam charge of the simulated run; the luminosity and every normalised yield are proportional to it}{EXPER.charge@\code{EXPER\%charge}}
\simcparam{debug}{param:debug}{6 flags}{all F}{—}{diagnostic print-outs: \code{debug(\allowbreak{}1)} prints the bremsstrahlung calculation; \code{debug(\allowbreak{}2.\allowbreak{}.\allowbreak{}6)} are accepted but print nothing}{debug@\code{debug}}
\simcparam{hard\_\allowbreak{}cuts}{param:hardu5fcuts}{flag}{F}{all}{apply the \code{SPedge} limits (and, for (e,e'p), \hyperref[param:cutsu25Emu25max]{\code{cuts\%\allowbreak{}Em\%\allowbreak{}max}}\simcidx{Parameters!cuts.Em.max@\code{cuts\%Em\%max}}) as cuts on the reconstructed events: events outside are not accepted and are not written. Without it the limits only decide which events count in the \code{.\allowbreak{}hist} wtcontr (pass\_\allowbreak{}cuts)}{hard_cuts@\code{hard\_cuts}}
\end{simcparams}
\subsection*{\texorpdfstring{Reaction selection and physics models}{Reaction selection and physics models}}\label{user_manual-04u5finputs:reaction-selection-and-physics-models}
\begin{simcparams}
\simcparam{doing\_\allowbreak{}phsp}{param:doingu5fphsp}{flag}{F}{phsp}{phase-space (acceptance) run: the cross section is 1 and normfac = 1. On hydrogen the elastic kinematics are kept (\hyperref[user_manual-05u5freactions:58-phase-space-and-acceptance-runs-doingu5fphsp]{5.8})}{doing_phsp@\code{doing\_phsp}}
\simcparam{doing\_\allowbreak{}pion}{param:doingu5fpion}{flag}{F}{pion, semi}{exclusive pion production; with \hyperref[param:doingu5fsemi]{\code{doing\_\allowbreak{}semi}}\simcidx{Parameters!doing_semi@\code{doing\_semi}}: SIDIS pions}{doing_pion@\code{doing\_pion}}
\simcparam{which\_\allowbreak{}pion}{param:whichu5fpion}{code}{0}{pion}{0: π+ (π0 with \hyperref[param:doingu5fpizero]{\code{doing\_\allowbreak{}pizero}}\simcidx{Parameters!doing_pizero@\code{doing\_pizero}}) on a proton, quasifree; 1: π− (π0 from the neutron), on a neutron; 2: π+Δ0 (π0Δ+); 3: π−Δ++ (π0Δ0 from the neutron); 10: coherent π+ A→A'; 11: coherent π−}{which_pion@\code{which\_pion}}
\simcparam{doing\_\allowbreak{}pizero}{param:doingu5fpizero}{flag}{F}{pion, semi}{the pion is a π0 (decays to γγ, detected in the calorimeter arm 7/8 or transported as a π0 through the HMS/SHMS, \hyperref[user_manual-06u5fspectrometers:66-u3c00-in-the-hms-or-shms]{6.6})}{doing_pizero@\code{doing\_pizero}}
\simcparam{pizero\_\allowbreak{}ngamma}{param:pizerou5fngamma}{count}{0}{pion, semi (π0)}{for π0 runs: 1 = at least one photon, 2 = both photons must hit the calorimeter. Must be 1 or 2 with \hyperref[param:doingu5fpizero]{\code{doing\_\allowbreak{}pizero}}\simcidx{Parameters!doing_pizero@\code{doing\_pizero}} (else the run stops)}{pizero_ngamma@\code{pizero\_ngamma}}
\simcparam{doing\_\allowbreak{}kaon}{param:doingu5fkaon}{flag}{F}{kaon, semi}{kaon electroproduction; with \hyperref[param:doingu5fsemi]{\code{doing\_\allowbreak{}semi}}\simcidx{Parameters!doing_semi@\code{doing\_semi}}: SIDIS kaons}{doing_kaon@\code{doing\_kaon}}
\simcparam{which\_\allowbreak{}kaon}{param:whichu5fkaon}{code}{0}{kaon}{0: K+Λ; 1: K+Σ0; 2: K+Σ− (on a neutron); 10/11/12: coherent production of the Λ/Σ0/Σ− hypernucleus}{which_kaon@\code{which\_kaon}}
\simcparam{doing\_\allowbreak{}delta}{param:doingu5fdelta}{flag}{F}{delta}{(e,e'p)π in the Delta region; the proton is detected. \textbf{The weight has no cross section} (\hyperref[user_manual-05u5freactions:55-eepu3c0-in-the-delta-region-doingu5fdelta]{5.5})}{doing_delta@\code{doing\_delta}}
\simcparam{doing\_\allowbreak{}semi}{param:doingu5fsemi}{flag}{F}{semi}{semi-inclusive (e,e'h)X; needs \hyperref[param:doingu5fpion]{\code{doing\_\allowbreak{}pion}}\simcidx{Parameters!doing_pion@\code{doing\_pion}} or \hyperref[param:doingu5fkaon]{\code{doing\_\allowbreak{}kaon}}\simcidx{Parameters!doing_kaon@\code{doing\_kaon}} (and \hyperref[param:doingu5fhplus]{\code{doing\_\allowbreak{}hplus}}\simcidx{Parameters!doing_hplus@\code{doing\_hplus}})}{doing_semi@\code{doing\_semi}}
\simcparam{doing\_\allowbreak{}hplus}{param:doingu5fhplus}{flag}{F}{semi, rho}{SIDIS and rho: the detected hadron is positive (1) or negative (0)}{doing_hplus@\code{doing\_hplus}}
\simcparam{doing\_\allowbreak{}rho}{param:doingu5frho}{flag}{F}{rho}{exclusive ρ0 → π+π−, one pion detected (H, D, ³He)}{doing_rho@\code{doing\_rho}}
\simcparam{doing\_\allowbreak{}decay}{param:doingu5fdecay}{flag}{F}{pion, kaon, semi, rho}{decay in flight of the hadron (π, K) in the spectrometer. Off: kaon runs multiply the weight by the survival probability}{doing_decay@\code{doing\_decay}}
\simcparam{ctau}{param:ctau}{cm}{0}{pion, kaon, semi, rho}{cτ of the hadron for the decay / survival probability (780.4 for pions, 371.3 for kaons)}{ctau@\code{ctau}}
\simcparam{do\_\allowbreak{}fermi}{param:dou5ffermi}{flag}{F}{semi}{SIDIS Fermi motion (D and A only; switched off for hydrogen with a WARNING)}{do_fermi@\code{do\_fermi}}
\simcparam{use\_\allowbreak{}benhar\_\allowbreak{}sf}{param:useu5fbenharu5fsf}{flag}{F}{D/A eep}{A(e,e'p): use the Benhar spectral function tables instead of the theory files}{use_benhar_sf@\code{use\_benhar\_sf}}
\simcparam{transparency}{param:transparency}{factor}{0}{D/A eep}{A(e,e'p) with \hyperref[param:useu5fbenharu5fsf]{\code{use\_\allowbreak{}benhar\_\allowbreak{}sf}}\simcidx{Parameters!use_benhar_sf@\code{use\_benhar\_sf}}: the nuclear transparency in the weight (SF\_weight = Z × transparency × SF). With the Benhar tables it must be set (0 gives weight 0)}{transparency@\code{transparency}}
\simcparam{deForest\_\allowbreak{}flag}{param:deForestu5fflag}{code}{0}{D/A eep}{off-shell e-p cross section of D/A(e,e'p): 0 = σcc1, 1 = σcc2, −1 = σcc1 with on-shell kinematics; |value| > 1 STOPs}{deForest_flag@\code{deForest\_flag}}
\simcparam{pt\_\allowbreak{}b\_\allowbreak{}param}{param:ptu5fbu5fparam}{GeV\textsuperscript{−}²}{0}{—}{SIDIS pT slope: \textbf{read and ignored} (the slope is b = 1/(0.120 z² + 0.200)); a WARNING is printed when it is set}{pt_b_param@\code{pt\_b\_param}}
\simcparam{sigc\_\allowbreak{}flag}{param:sigcu5fflag}{code}{0}{—}{SIDIS "central cross section" binning (0 = z, 1 = pT²): \textbf{unused} (the central cross-section calculation is disabled)}{sigc_flag@\code{sigc\_flag}}
\simcparam{sigc\_\allowbreak{}nbin}{param:sigcu5fnbin}{count}{0}{—}{unused (as \hyperref[param:sigcu5fflag]{\code{sigc\_\allowbreak{}flag}}\simcidx{Parameters!sigc_flag@\code{sigc\_flag}})}{sigc_nbin@\code{sigc\_nbin}}
\simcparam{sigc\_\allowbreak{}kin\_\allowbreak{}min}{param:sigcu5fkinu5fmin}{z or GeV²}{0}{—}{unused (as \hyperref[param:sigcu5fflag]{\code{sigc\_\allowbreak{}flag}}\simcidx{Parameters!sigc_flag@\code{sigc\_flag}})}{sigc_kin_min@\code{sigc\_kin\_min}}
\simcparam{sigc\_\allowbreak{}kin\_\allowbreak{}max}{param:sigcu5fkinu5fmax}{z or GeV²}{0}{—}{unused (as \hyperref[param:sigcu5fflag]{\code{sigc\_\allowbreak{}flag}}\simcidx{Parameters!sigc_flag@\code{sigc\_flag}})}{sigc_kin_max@\code{sigc\_kin\_max}}
\simcparam{sigc\_\allowbreak{}kin\_\allowbreak{}ind}{param:sigcu5fkinu5find}{GeV² or z}{0}{—}{unused (as \hyperref[param:sigcu5fflag]{\code{sigc\_\allowbreak{}flag}}\simcidx{Parameters!sigc_flag@\code{sigc\_flag}})}{sigc_kin_ind@\code{sigc\_kin\_ind}}
\end{simcparams}
\subsection*{\texorpdfstring{Beam and spectrometer settings}{Beam and spectrometer settings}}\label{user_manual-04u5finputs:beam-and-spectrometer-settings}
\begin{simcparams}
\simcparam{Ebeam}{param:Ebeam}{MeV}{0}{all}{beam energy}{Ebeam@\code{Ebeam}}
\simcparam{dEbeam}{param:dEbeam}{\% of Ebeam}{0}{all}{full width of the flat beam-energy spread}{dEbeam@\code{dEbeam}}
\simcparam{electron\_\allowbreak{}arm}{param:electronu5farm}{code}{0}{all}{spectrometer of the electron: 1 HMS, 2 SOS, 3 HRS-R, 4 HRS-L, 5 SHMS, 7/8 calorimeter. \textbf{6 is rejected} (the retired SHMS LSA tune); other values stop the run (\code{I\ dont\ know\ what\ phi\ should\ be\ for\ the\ electron\ arm}) (\hyperref[user_manual-06u5fspectrometers:61-spectrometer-numbers]{6.1})}{electron_arm@\code{electron\_arm}}
\simcparam{hadron\_\allowbreak{}arm}{param:hadronu5farm}{code}{0}{all}{spectrometer of the hadron, same codes}{hadron_arm@\code{hadron\_arm}}
\simcparam{spec\%\allowbreak{}e\%\allowbreak{}P}{param:specu25eu25P}{MeV/c}{0}{all}{central momentum of the electron arm}{spec.e.P@\code{spec\%e\%P}}
\simcparam{spec\%\allowbreak{}e\%\allowbreak{}theta}{param:specu25eu25theta}{deg}{0}{all}{central angle of the electron arm (absolute value)}{spec.e.theta@\code{spec\%e\%theta}}
\simcparam{spec\%\allowbreak{}p\%\allowbreak{}P}{param:specu25pu25P}{MeV/c}{0}{all}{central momentum of the hadron arm}{spec.p.P@\code{spec\%p\%P}}
\simcparam{spec\%\allowbreak{}p\%\allowbreak{}theta}{param:specu25pu25theta}{deg}{0}{all}{central angle of the hadron arm (absolute value)}{spec.p.theta@\code{spec\%p\%theta}}
\simcparam{spect\_\allowbreak{}mode}{param:spectu5fmode}{code}{0}{all}{0: both spectrometer Monte Carlos on; −1: electron arm off (hadron only); −2: hadron arm off (electron only); 1: both off (\hyperref[user_manual-06u5fspectrometers:64-spectu5fmode-switching-the-spectrometer-monte-carlo-off]{6.4})}{spect_mode@\code{spect\_mode}}
\simcparam{use\_\allowbreak{}first\_\allowbreak{}cer}{param:useu5ffirstu5fcer}{flag}{0}{—}{"use first Cherenkov in SHMS": printed in the \code{.\allowbreak{}hist}, not used by the simulation}{use_first_cer@\code{use\_first\_cer}}
\end{simcparams}
\subsection*{\texorpdfstring{Beam spot, raster and target position}{Beam spot, raster and target position}}\label{user_manual-04u5finputs:beam-spot-raster-and-target-position}
\begin{simcparams}
\simcparam{gen\%\allowbreak{}xwid}{param:genu25xwid}{cm}{0}{all}{horizontal beam size, one σ (gaussian)}{gen.xwid@\code{gen\%xwid}}
\simcparam{gen\%\allowbreak{}ywid}{param:genu25ywid}{cm}{0}{all}{vertical beam size, one σ}{gen.ywid@\code{gen\%ywid}}
\simcparam{targ\%\allowbreak{}fr\_\allowbreak{}pattern}{param:targu25fru5fpattern}{code}{0}{all}{fast raster: 1 = old "bedpost" rectangle, 2 = circle, 3 = flat rectangle; other values: no raster}{targ.fr_pattern@\code{targ\%fr\_pattern}}
\simcparam{targ\%\allowbreak{}fr1}{param:targu25fr1}{cm}{0}{all}{raster x half-width (patterns 1, 3) or inner radius (2)}{targ.fr1@\code{targ\%fr1}}
\simcparam{targ\%\allowbreak{}fr2}{param:targu25fr2}{cm}{0}{all}{raster y half-width (patterns 1, 3) or outer radius (2)}{targ.fr2@\code{targ\%fr2}}
\simcparam{targ\%\allowbreak{}xoffset}{param:targu25xoffset}{cm}{0}{all}{target x offset (+x = beam right)}{targ.xoffset@\code{targ\%xoffset}}
\simcparam{targ\%\allowbreak{}yoffset}{param:targu25yoffset}{cm}{0}{all}{target y offset (+y = up)}{targ.yoffset@\code{targ\%yoffset}}
\simcparam{targ\%\allowbreak{}zoffset}{param:targu25zoffset}{cm}{0}{all}{target z offset (+z = downstream)}{targ.zoffset@\code{targ\%zoffset}}
\end{simcparams}
\subsection*{\texorpdfstring{Target}{Target}}\label{user_manual-04u5finputs:target}
\begin{simcparams}
\simcparam{targ\%\allowbreak{}A}{param:targu25A}{nucleons}{0}{all}{target mass number. Chooses H (1), D (2) or a nucleus (≥ 3) for every reaction}{targ.A@\code{targ\%A}}
\simcparam{targ\%\allowbreak{}Z}{param:targu25Z}{charge}{0}{all}{target charge. Z < 2.4 means a liquid (cryogenic) target}{targ.Z@\code{targ\%Z}}
\simcparam{targ\%\allowbreak{}mass\_\allowbreak{}amu}{param:targu25massu5famu}{amu}{0}{all}{target mass (forced to Mp / Md for H / D with a WARNING)}{targ.mass_amu@\code{targ\%mass\_amu}}
\simcparam{targ\%\allowbreak{}mrec\_\allowbreak{}amu}{param:targu25mrecu5famu}{amu}{0}{all}{recoil (A−1 system) mass for (e,e'p); final-nucleus mass for coherent production. Replaced by the program's estimate when it is more than 100 MeV off}{targ.mrec_amu@\code{targ\%mrec\_amu}}
\simcparam{targ\%\allowbreak{}rho}{param:targu25rho}{g/cm³}{0}{all}{target density (with \hyperref[param:targu25thick]{\code{targ\%\allowbreak{}thick}}\simcidx{Parameters!targ.thick@\code{targ\%thick}} it gives the target length)}{targ.rho@\code{targ\%rho}}
\simcparam{targ\%\allowbreak{}thick}{param:targu25thick}{mg/cm²}{0}{all}{target thickness}{targ.thick@\code{targ\%thick}}
\simcparam{targ\%\allowbreak{}angle}{param:targu25angle}{deg}{0}{all}{angle of a solid target (0 = perpendicular to the beam); forced to 0 for cryotargets; sin(angle) > 0.85 STOPs}{targ.angle@\code{targ\%angle}}
\simcparam{targ\%\allowbreak{}abundancy}{param:targu25abundancy}{\%}{0}{all}{isotopic purity of the target. \textbf{Must be set} (100 for a pure target): the luminosity is divided by abundancy/\allowbreak{}100,\allowbreak{} so 0 gives zero luminosity and zero weights}{targ.abundancy@\code{targ\%abundancy}}
\simcparam{targ\%\allowbreak{}can}{param:targu25can}{code}{0}{all}{cryotarget cell: 1 = beer can, 2 = pudding can, 3 = cryo17. Required for Z < 2.4 (\code{bad\ targ.\allowbreak{}can\ value} STOP otherwise)}{targ.can@\code{targ\%can}}
\end{simcparams}
\subsection*{\texorpdfstring{Polarized target and target field}{Polarized target and target field}}\label{user_manual-04u5finputs:polarized-target-and-target-field}
\begin{simcparams}
\simcparam{using\_\allowbreak{}tgt\_\allowbreak{}field}{param:usingu5ftgtu5ffield}{flag}{F}{field}{track the particles through the polarized-target magnetic field (\hyperref[user_manual-06u5fspectrometers:65-polarized-target-field-usingu5ftgtu5ffield]{6.5})}{using_tgt_field@\code{using\_tgt\_field}}
\simcparam{tgt\_\allowbreak{}field\_\allowbreak{}file}{param:tgtu5ffieldu5ffile}{path}{blank}{field}{field map. \code{'0'} = no field; \code{'\ '} = a uniform 5 T test field; otherwise a map file looked for in \code{-{}-{}data-{}dir}\simcidx{Command-line options!data-dir@\code{-{}-{}data-{}dir}}, the current directory, then the data directory's \code{field/} (\code{trg\_\allowbreak{}field\_\allowbreak{}map.\allowbreak{}dat} is included). Unset with \code{using\_\allowbreak{}tgt\_\allowbreak{}field\ =\ 1}: refused at start-up}{tgt_field_file@\code{tgt\_field\_file}}
\simcparam{targ\_\allowbreak{}Bangle}{param:targu5fBangle}{deg}{0}{field}{angle of the target field (+ = beam left); also the polarization direction used for the polarized-target angle columns of the tree (\hyperref[user_manual-07u5foutputs:717-polarized-target-columns-usingu5ftgtu5ffield]{7.1.7})}{targ_Bangle@\code{targ\_Bangle}}
\simcparam{targ\_\allowbreak{}Bphi}{param:targu5fBphi}{deg}{0}{field}{side of the beam the field points to (90 = beam left, 270 = beam right)}{targ_Bphi@\code{targ\_Bphi}}
\simcparam{targ\_\allowbreak{}pol}{param:targu5fpol}{factor}{0}{field}{target polarization (+ = along B); 0 makes the \code{th\_\allowbreak{}tarq} column NaN}{targ_pol@\code{targ\_pol}}
\end{simcparams}
\subsection*{\texorpdfstring{Spectrometer offsets}{Spectrometer offsets}}\label{user_manual-04u5finputs:spectrometer-offsets}
Offsets applied to the event before the spectrometer Monte Carlo of each arm (misalignments).

\begin{simcparams}
\simcparam{spec\%\allowbreak{}e\%\allowbreak{}offset\%\allowbreak{}x}{param:specu25eu25offsetu25x}{cm}{0}{all}{electron arm x offset}{spec.e.offset.x@\code{spec\%e\%offset\%x}}
\simcparam{spec\%\allowbreak{}e\%\allowbreak{}offset\%\allowbreak{}y}{param:specu25eu25offsetu25y}{cm}{0}{all}{electron arm y offset}{spec.e.offset.y@\code{spec\%e\%offset\%y}}
\simcparam{spec\%\allowbreak{}e\%\allowbreak{}offset\%\allowbreak{}z}{param:specu25eu25offsetu25z}{cm}{0}{all}{electron arm z offset}{spec.e.offset.z@\code{spec\%e\%offset\%z}}
\simcparam{spec\%\allowbreak{}e\%\allowbreak{}offset\%\allowbreak{}xptar}{param:specu25eu25offsetu25xptar}{mrad}{0}{all}{electron arm xptar (vertical slope) offset}{spec.e.offset.xptar@\code{spec\%e\%offset\%xptar}}
\simcparam{spec\%\allowbreak{}e\%\allowbreak{}offset\%\allowbreak{}yptar}{param:specu25eu25offsetu25yptar}{mrad}{0}{all}{electron arm yptar (horizontal slope) offset}{spec.e.offset.yptar@\code{spec\%e\%offset\%yptar}}
\simcparam{spec\%\allowbreak{}p\%\allowbreak{}offset\%\allowbreak{}x}{param:specu25pu25offsetu25x}{cm}{0}{all}{hadron arm x offset}{spec.p.offset.x@\code{spec\%p\%offset\%x}}
\simcparam{spec\%\allowbreak{}p\%\allowbreak{}offset\%\allowbreak{}y}{param:specu25pu25offsetu25y}{cm}{0}{all}{hadron arm y offset}{spec.p.offset.y@\code{spec\%p\%offset\%y}}
\simcparam{spec\%\allowbreak{}p\%\allowbreak{}offset\%\allowbreak{}z}{param:specu25pu25offsetu25z}{cm}{0}{all}{hadron arm z offset}{spec.p.offset.z@\code{spec\%p\%offset\%z}}
\simcparam{spec\%\allowbreak{}p\%\allowbreak{}offset\%\allowbreak{}xptar}{param:specu25pu25offsetu25xptar}{mrad}{0}{all}{hadron arm xptar offset}{spec.p.offset.xptar@\code{spec\%p\%offset\%xptar}}
\simcparam{spec\%\allowbreak{}p\%\allowbreak{}offset\%\allowbreak{}yptar}{param:specu25pu25offsetu25yptar}{mrad}{0}{all}{hadron arm yptar offset}{spec.p.offset.yptar@\code{spec\%p\%offset\%yptar}}
\end{simcparams}
\subsection*{\texorpdfstring{Spectrometer limits and missing-energy cuts}{Spectrometer limits and missing-energy cuts}}\label{user_manual-04u5finputs:spectrometer-limits-and-missing-energy-cuts}
The \code{SPedge} limits describe the spectrometer acceptance that you want to simulate. SIMC++ widens them by the
spectrometer "slop" (resolution,\allowbreak{} multiple scattering, energy loss) to build the \textbf{generation limits}.
Reconstructed events inside the unwidened limits count in the \code{.\allowbreak{}hist} wtcontr (pass\_\allowbreak{}cuts); with \hyperref[param:hardu5fcuts]{\code{hard\_\allowbreak{}cuts}}\simcidx{Parameters!hard_cuts@\code{hard\_cuts}} they are
also the only accepted events. The defaults (±1e10) are wide open: \textbf{set them} in every input. A delta minimum
≤ −100 \% is changed to −99.99 \%.

\begin{simcparams}
\simcparam{SPedge\%\allowbreak{}e\%\allowbreak{}delta\%\allowbreak{}min}{param:SPedgeu25eu25deltau25min}{\%}{−1e10}{all}{electron arm delta lower limit}{SPedge.e.delta.min@\code{SPedge\%e\%delta\%min}}
\simcparam{SPedge\%\allowbreak{}e\%\allowbreak{}delta\%\allowbreak{}max}{param:SPedgeu25eu25deltau25max}{\%}{1e10}{all}{electron arm delta upper limit}{SPedge.e.delta.max@\code{SPedge\%e\%delta\%max}}
\simcparam{SPedge\%\allowbreak{}e\%\allowbreak{}yptar\%\allowbreak{}min}{param:SPedgeu25eu25yptaru25min}{mrad}{−1e10}{all}{electron arm yptar lower limit}{SPedge.e.yptar.min@\code{SPedge\%e\%yptar\%min}}
\simcparam{SPedge\%\allowbreak{}e\%\allowbreak{}yptar\%\allowbreak{}max}{param:SPedgeu25eu25yptaru25max}{mrad}{1e10}{all}{electron arm yptar upper limit}{SPedge.e.yptar.max@\code{SPedge\%e\%yptar\%max}}
\simcparam{SPedge\%\allowbreak{}e\%\allowbreak{}xptar\%\allowbreak{}min}{param:SPedgeu25eu25xptaru25min}{mrad}{−1e10}{all}{electron arm xptar lower limit}{SPedge.e.xptar.min@\code{SPedge\%e\%xptar\%min}}
\simcparam{SPedge\%\allowbreak{}e\%\allowbreak{}xptar\%\allowbreak{}max}{param:SPedgeu25eu25xptaru25max}{mrad}{1e10}{all}{electron arm xptar upper limit}{SPedge.e.xptar.max@\code{SPedge\%e\%xptar\%max}}
\simcparam{SPedge\%\allowbreak{}p\%\allowbreak{}delta\%\allowbreak{}min}{param:SPedgeu25pu25deltau25min}{\%}{−1e10}{all}{hadron arm delta lower limit}{SPedge.p.delta.min@\code{SPedge\%p\%delta\%min}}
\simcparam{SPedge\%\allowbreak{}p\%\allowbreak{}delta\%\allowbreak{}max}{param:SPedgeu25pu25deltau25max}{\%}{1e10}{all}{hadron arm delta upper limit (the legacy-compatible mode uses \hyperref[param:SPedgeu25eu25deltau25max]{\code{SPedge\%\allowbreak{}e\%\allowbreak{}delta\%\allowbreak{}max}}\simcidx{Parameters!SPedge.e.delta.max@\code{SPedge\%e\%delta\%max}} here, see \hyperref[user_manual-08u5frunningu5fandu5freproducibility:83-default-and-legacy-compatible-mode]{8.3})}{SPedge.p.delta.max@\code{SPedge\%p\%delta\%max}}
\simcparam{SPedge\%\allowbreak{}p\%\allowbreak{}yptar\%\allowbreak{}min}{param:SPedgeu25pu25yptaru25min}{mrad}{−1e10}{all}{hadron arm yptar lower limit}{SPedge.p.yptar.min@\code{SPedge\%p\%yptar\%min}}
\simcparam{SPedge\%\allowbreak{}p\%\allowbreak{}yptar\%\allowbreak{}max}{param:SPedgeu25pu25yptaru25max}{mrad}{1e10}{all}{hadron arm yptar upper limit}{SPedge.p.yptar.max@\code{SPedge\%p\%yptar\%max}}
\simcparam{SPedge\%\allowbreak{}p\%\allowbreak{}xptar\%\allowbreak{}min}{param:SPedgeu25pu25xptaru25min}{mrad}{−1e10}{all}{hadron arm xptar lower limit}{SPedge.p.xptar.min@\code{SPedge\%p\%xptar\%min}}
\simcparam{SPedge\%\allowbreak{}p\%\allowbreak{}xptar\%\allowbreak{}max}{param:SPedgeu25pu25xptaru25max}{mrad}{1e10}{all}{hadron arm xptar upper limit}{SPedge.p.xptar.max@\code{SPedge\%p\%xptar\%max}}
\simcparam{cuts\%\allowbreak{}Em\%\allowbreak{}min}{param:cutsu25Emu25min}{MeV}{−1e10}{eep}{A(e,e'p): lower end of the generated missing-energy range (widened by the slop, at least 0). Ignored (wide open) for pion, kaon and Delta runs and when min = max}{cuts.Em.min@\code{cuts\%Em\%min}}
\simcparam{cuts\%\allowbreak{}Em\%\allowbreak{}max}{param:cutsu25Emu25max}{MeV}{1e10}{eep}{A(e,e'p): upper end of the generated Em range; with \hyperref[param:hardu5fcuts]{\code{hard\_\allowbreak{}cuts}}\simcidx{Parameters!hard_cuts@\code{hard\_cuts}}, (e,e'p) events with recon Em above it are rejected}{cuts.Em.max@\code{cuts\%Em\%max}}
\end{simcparams}
\subsection*{\texorpdfstring{Simulation switches}{Simulation switches}}\label{user_manual-04u5finputs:simulation-switches}
\begin{simcparams}
\simcparam{mc\_\allowbreak{}smear}{param:mcu5fsmear}{flag}{F}{all}{multiple scattering in the target and spectrometer materials, and wire-chamber smearing (the focal-plane resolution)}{mc_smear@\code{mc\_smear}}
\simcparam{using\_\allowbreak{}Eloss}{param:usingu5fEloss}{flag}{F}{all}{energy loss in the target and windows}{using_Eloss@\code{using\_Eloss}}
\simcparam{correct\_\allowbreak{}Eloss}{param:correctu5fEloss}{flag}{F}{all}{correct the reconstructed momenta for the most probable energy loss (forced off when \hyperref[param:usingu5fEloss]{\code{using\_\allowbreak{}Eloss}}\simcidx{Parameters!using_Eloss@\code{using\_Eloss}} is off)}{correct_Eloss@\code{correct\_Eloss}}
\simcparam{correct\_\allowbreak{}raster}{param:correctu5fraster}{flag}{F}{all}{reconstruct with the raster position (the vertical beam position enters the reconstruction)}{correct_raster@\code{correct\_raster}}
\simcparam{using\_\allowbreak{}Coulomb}{param:usingu5fCoulomb}{flag}{F}{all}{Coulomb correction of the electron energy (forced off for Z = 1)}{using_Coulomb@\code{using\_Coulomb}}
\simcparam{using\_\allowbreak{}HMScoll}{param:usingu5fHMScoll}{flag}{F}{HMS arm}{step pions/muons through the HMS collimator (absorption,\allowbreak{} scattering, the \code{.\allowbreak{}geni} collimator line)}{using_HMScoll@\code{using\_HMScoll}}
\simcparam{using\_\allowbreak{}SHMScoll}{param:usingu5fSHMScoll}{flag}{F}{SHMS arm}{the same for the SHMS collimator}{using_SHMScoll@\code{using\_SHMScoll}}
\simcparam{drift\_\allowbreak{}to\_\allowbreak{}cal}{param:driftu5ftou5fcal}{cm}{0}{calo}{distance from the target to the calorimeter face (arms 7/8). 0 puts the face at the target (not checked)}{drift_to_cal@\code{drift\_to\_cal}}
\simcparam{using\_\allowbreak{}cit\_\allowbreak{}generation}{param:usingu5fcitu5fgeneration}{flag}{F}{—}{accepted for compatibility,\allowbreak{} not used}{using_cit_generation@\code{using\_cit\_generation}}
\end{simcparams}
\subsection*{\texorpdfstring{Radiative corrections}{Radiative corrections}}\label{user_manual-04u5finputs:radiative-corrections}
\begin{simcparams}
\simcparam{using\_\allowbreak{}rad}{param:usingu5frad}{flag}{F}{rad}{radiative corrections: generate the radiative tails}{using_rad@\code{using\_rad}}
\simcparam{rad\_\allowbreak{}flag}{param:radu5fflag}{code}{0}{rad}{radiative option: 0 = best formulas, (ntail, Eγ) basis; 1 = BASICRAD only, (ntail, Eγ) basis; 2, 3 = BASICRAD in the (Eγ1, Eγ2, Eγ3) basis. Set to 0 by \hyperref[param:doingu5fphsp]{\code{doing\_\allowbreak{}phsp}}\simcidx{Parameters!doing_phsp@\code{doing\_phsp}}; other values STOP}{rad_flag@\code{rad\_flag}}
\simcparam{extrad\_\allowbreak{}flag}{param:extradu5fflag}{code}{0}{rad}{external radiation: 1 = basic formulas; 2 = basic × φ; 3 = Friedrich approximation; 0 = default (3 for rad\_flag 0, 1 otherwise); < 0 STOPs}{extrad_flag@\code{extrad\_flag}}
\simcparam{one\_\allowbreak{}tail}{param:oneu5ftail}{code}{0}{rad}{which tails radiate: 0 all; 1 incoming e; 2 scattered e'; 3 hadron; −1 all but the incoming e; −2 all but e'; −3 all but the hadron. |one\_tail| > 3 with \hyperref[param:usingu5frad]{\code{using\_\allowbreak{}rad}}\simcidx{Parameters!using_rad@\code{using\_rad}} STOPs}{one_tail@\code{one\_tail}}
\simcparam{intcor\_\allowbreak{}mode}{param:intcoru5fmode}{code}{0}{rad}{internal-correction formula: 0 = Schwinger; otherwise brem (or bremos with \hyperref[param:useu5foffshellu5frad]{\code{use\_\allowbreak{}offshell\_\allowbreak{}rad}}\simcidx{Parameters!use_offshell_rad@\code{use\_offshell\_rad}}). Normally 1}{intcor_mode@\code{intcor\_mode}}
\simcparam{use\_\allowbreak{}expon}{param:useu5fexpon}{flag}{0}{rad}{1 = exponentiate the corrections (brem); normally 0}{use_expon@\code{use\_expon}}
\simcparam{use\_\allowbreak{}offshell\_\allowbreak{}rad}{param:useu5foffshellu5frad}{flag}{F}{rad}{use the off-shell bremos formula for the hard correction}{use_offshell_rad@\code{use\_offshell\_rad}}
\simcparam{lambda}{param:lambda}{3 numbers}{0, 0, 0}{—}{radiative λ(1:3) for rad\_flag 4: \textbf{overwritten} for every event by the computed values, so the input has no effect}{lambda@\code{lambda}}
\simcparam{Egamma\_\allowbreak{}gen\_\allowbreak{}max}{param:Egammau5fgenu5fmax}{MeV}{0}{rad}{> 0.01: use this fixed maximum photon energy instead of the computed limits}{Egamma_gen_max@\code{Egamma\_gen\_max}}
\simcparam{dE\_\allowbreak{}edge\_\allowbreak{}test}{param:dEu5fedgeu5ftest}{MeV}{0}{all}{widens the energy edges of the generation and radiation limits by this amount (a test of the edges)}{dE_edge_test@\code{dE\_edge\_test}}
\end{simcparams}
\section{\texorpdfstring{Settings that have no \code{.\allowbreak{}inp} variable}{Settings that have no .inp variable}}\label{user_manual-04u5finputs:47-settings-that-have-no-inp-variable}
These have no \code{.\allowbreak{}inp} counterpart:

\begin{simctable}{small}{4}{>{\raggedright\arraybackslash}p{\dimexpr 0.297\linewidth-2\tabcolsep\relax}>{\raggedright\arraybackslash}p{\dimexpr 0.175\linewidth-2\tabcolsep\relax}>{\raggedright\arraybackslash}p{\dimexpr 0.289\linewidth-2\tabcolsep\relax}>{\raggedright\arraybackslash}p{\dimexpr 0.150\linewidth-2\tabcolsep\relax}}
\caption{Settings that have no \code{.\allowbreak{}inp} variable}\\%
\toprule\rowcolor{simchead}\simcth{setting} & \simcth{JSON} & \simcth{command line} & \simcth{default} \\\midrule\endfirsthead
\multicolumn{4}{@{}l}{\simccontinued}\\\toprule\rowcolor{simchead}\simcth{setting} & \simcth{JSON} & \simcth{command line} & \simcth{default} \\\midrule\endhead
\midrule\multicolumn{4}{r@{}}{\simccontinues}\\\endfoot
\bottomrule\endlastfoot
legacy-compatible mode & \code{run.\allowbreak{}fortran\_\allowbreak{}compat} & \code{-{}-{}fortran-{}compat}\simcidx{Command-line options!fortran-compat@\code{-{}-{}fortran-{}compat}} & off \\
output name and directory & \code{run.\allowbreak{}name}, \code{run.\allowbreak{}output\_\allowbreak{}dir} & \code{-{}-{}run-{}name\ NAME}\simcidx{Command-line options!run-name@\code{-{}-{}run-{}name}}, \code{-{}o\ DIR} & input stem, \code{.} \\
plug-ins, random-number generator & \code{pipeline} & \code{-{}-{}pipeline\ FILE}\simcidx{Command-line options!pipeline@\code{-{}-{}pipeline}} & from the flags (\hyperref[user_manual-04u5finputs:43-the-pipeline-section-choosing-plug-ins]{4.3}) \\
event-loop threads & \code{pipeline.\allowbreak{}threads} & \code{-{}-{}threads\ N}\simcidx{Command-line options!threads@\code{-{}-{}threads}} & 1 \\
data directory & — & \code{-{}-{}data-{}dir\ DIR}\simcidx{Command-line options!data-dir@\code{-{}-{}data-{}dir}}, env \code{SIMC\_\allowbreak{}DATA\_\allowbreak{}DIR} & \hyperref[user_manual-02u5finstallation:25-where-the-data-files-are-found]{2.5} \\
directory of \hyperref[param:extrau5fdbaseu5ffile]{\code{extra\_\allowbreak{}dbase\_\allowbreak{}file}}\simcidx{Parameters!extra_dbase_file@\code{extra\_dbase\_file}} & — & \code{-{}-{}infiles-{}dir\ DIR}\simcidx{Command-line options!infiles-dir@\code{-{}-{}infiles-{}dir}} & the main file's directory \\
number of events, seed & \code{parameters.\allowbreak{}ngen}, \code{parameters.\allowbreak{}random\_\allowbreak{}seed} & \code{-{}-{}ngen\ N}\simcidx{Command-line options!ngen@\code{-{}-{}ngen}}, \code{-{}-{}seed\ N}\simcidx{Command-line options!seed@\code{-{}-{}seed}} (override the input; the \code{.\allowbreak{}hist} prints the overridden values) & from the input \\
quiet stages & \code{"verbose":\ false} per stage & \code{-{}-{}quiet}\simcidx{Command-line options!quiet@\code{-{}-{}quiet}} & verbose \\
print the resolved pipeline & — & \code{-{}-{}print-{}pipeline}\simcidx{Command-line options!print-pipeline@\code{-{}-{}print-{}pipeline}} & off \\
print the data directory and exit & — & \code{-{}-{}print-{}data-{}dir}\simcidx{Command-line options!print-data-dir@\code{-{}-{}print-{}data-{}dir}} & — \\
usage & — & \code{-{}h}, \code{-{}-{}help}\simcidx{Command-line options!help@\code{-{}-{}help}} & — \\
\end{simctable}

% docs/manual/latex/user/05_reactions.tex - chapter of the LaTeX edition of the book "user_manual" (SIMC++ documentation).
%   Converted once from docs/manual/markdown/user/05_reactions.md by tools/md2latex.py; from then on THIS FILE is the source of the
%   LaTeX edition and may be edited by hand. Re-running tools/md2latex.py --write-sources for this
%   chapter OVERWRITES the hand edits (see docs/README.md).
%   Compile: cmake --build <build> --target pdf, or tools/build_docs_pdf.py (figures, xelatex passes,
%   makeindex with docs/common/latex/simcpp.ist); by hand, see docs/README.md. XeLaTeX, class
%   docs/common/latex/simcpp.cls (packages: fontspec, tikz, xcolor, hyperref, listings, longtable, booktabs,
%   titlesec, titletoc, fancyhdr, makeidx, preview, ...); fonts TeX Gyre Pagella, TeX Gyre Heros, DejaVu Sans
%   Mono. The figures are built first from docs/{manual,translation}/figures/*.tex (preview.sty) into PDFs.
\chapter{\texorpdfstring{Reactions}{Reactions}}\label{user_manual-05u5freactions}
\section{\texorpdfstring{How the reaction is chosen}{How the reaction is chosen}}\label{user_manual-05u5freactions:51-how-the-reaction-is-chosen}
The reaction follows from the reaction flags and from \hyperref[param:targu25A]{\code{targ\%\allowbreak{}A}}\simcidx{Parameters!targ.A@\code{targ\%A}}. SIMC++ then picks the reaction and physics-model plug-ins
(\code{Pipeline\allowbreak{}Config::\allowbreak{}Defaults}\simcidx{Classes and interfaces!PipelineConfig@\code{PipelineConfig}}, \code{src/\allowbreak{}driver/\allowbreak{}Pipeline.\allowbreak{}cpp}). The first matching row wins:

\begin{simctable}{footnotesize}{5}{>{\raggedright\arraybackslash}p{\dimexpr 0.218\linewidth-2\tabcolsep\relax}>{\raggedright\arraybackslash}p{\dimexpr 0.128\linewidth-2\tabcolsep\relax}>{\raggedright\arraybackslash}p{\dimexpr 0.197\linewidth-2\tabcolsep\relax}>{\raggedright\arraybackslash}p{\dimexpr 0.150\linewidth-2\tabcolsep\relax}>{\raggedright\arraybackslash}p{\dimexpr 0.208\linewidth-2\tabcolsep\relax}}
\caption{How the reaction is chosen}\\%
\toprule\rowcolor{simchead}\simcth{flags} & \simcth{target} & \simcth{reaction plug-in} & \simcth{physics model} & \simcth{tree layout (\hyperref[user_manual-07u5foutputs:71-the-event-tree-h10]{7.1})} \\\midrule\endfirsthead
\multicolumn{5}{@{}l}{\simccontinued}\\\toprule\rowcolor{simchead}\simcth{flags} & \simcth{target} & \simcth{reaction plug-in} & \simcth{physics model} & \simcth{tree layout (\hyperref[user_manual-07u5foutputs:71-the-event-tree-h10]{7.1})} \\\midrule\endhead
\midrule\multicolumn{5}{r@{}}{\simccontinues}\\\endfoot
\bottomrule\endlastfoot
\code{doing\_\allowbreak{}phsp\ =\ 1} & any & \code{phase\_\allowbreak{}space} & none (σ = 1) & as the reaction it replaces \\
none of the flags below & A = 1 & \code{hydrogen\_\allowbreak{}elastic} & \code{sigep} & (e,e'p), 46 columns \\
none & A = 2 & \code{deuterium\_\allowbreak{}eep} & \code{deforest} & (e,e'p) \\
none & A ≥ 3 & \code{nucleus\_\allowbreak{}eep} & \code{deforest} & (e,e'p) \\
\hyperref[param:doingu5fnucu5felast]{\code{doing\_\allowbreak{}nuc\_\allowbreak{}elast}}\simcidx{Parameters!doing_nuc_elast@\code{doing\_nuc\_elast}} & any & \code{nuclear\_\allowbreak{}elastic} & \code{nuc\_\allowbreak{}elastic} & (e,e'p), 46 columns \\
\hyperref[param:doingu5fpion]{\code{doing\_\allowbreak{}pion}}\simcidx{Parameters!doing_pion@\code{doing\_pion}} & any & \code{pion\_\allowbreak{}production} & \code{pion} & pion, 53 (65 with π0) \\
\hyperref[param:doingu5fkaon]{\code{doing\_\allowbreak{}kaon}}\simcidx{Parameters!doing_kaon@\code{doing\_kaon}} & any & \code{kaon\_\allowbreak{}production} & \code{kaon} & kaon, 55 \\
\hyperref[param:doingu5fdelta]{\code{doing\_\allowbreak{}delta}}\simcidx{Parameters!doing_delta@\code{doing\_delta}} & any & \code{delta\_\allowbreak{}production} & \code{delta} & pion layout, 53 \\
\hyperref[param:doingu5fsemi]{\code{doing\_\allowbreak{}semi}}\simcidx{Parameters!doing_semi@\code{doing\_semi}} + \hyperref[param:doingu5fpion]{\code{doing\_\allowbreak{}pion}}\simcidx{Parameters!doing_pion@\code{doing\_pion}} or \hyperref[param:doingu5fkaon]{\code{doing\_\allowbreak{}kaon}}\simcidx{Parameters!doing_kaon@\code{doing\_kaon}} & any & \code{sidis\_\allowbreak{}production} & \code{sidis} & SIDIS, 56 (68 with π0) \\
\hyperref[param:doingu5frho]{\code{doing\_\allowbreak{}rho}}\simcidx{Parameters!doing_rho@\code{doing\_rho}} & A = 1, 2, 3 & \code{rho\_\allowbreak{}production} & \code{rho} & rho, 59 \\
\end{simctable}
\hyperref[param:doingu5fsemi]{\code{doing\_\allowbreak{}semi}}\simcidx{Parameters!doing_semi@\code{doing\_semi}} takes priority: \hyperref[param:doingu5fpion]{\code{doing\_\allowbreak{}pion}}\simcidx{Parameters!doing_pion@\code{doing\_pion}}/\hyperref[param:doingu5fkaon]{\code{doing\_\allowbreak{}kaon}}\simcidx{Parameters!doing_kaon@\code{doing\_kaon}} then only select the SIDIS hadron. With \hyperref[param:usingu5ftgtu5ffield]{\code{using\_\allowbreak{}tgt\_\allowbreak{}field}}\simcidx{Parameters!using_tgt_field@\code{using\_tgt\_field}}, 8
polarized-target columns are added (\hyperref[user_manual-07u5foutputs:718-polarized-target-columns-usingu5ftgtu5ffield]{7.1.7}). The log shows
the reaction title (for example \code{\ ****-{}-{}-{}-{}-{}-{}-{}-{}\ \ H(\allowbreak{}e,\allowbreak{}e'pi+)X\ \ -{}-{}-{}-{}-{}-{}-{}-{}****}).

Common to all reactions: the beam (energy spread \hyperref[param:dEbeam]{\code{dEbeam}}\simcidx{Parameters!dEbeam@\code{dEbeam}}, size \code{gen\%\allowbreak{}xwid/\allowbreak{}ywid}, raster), the target (energy loss,
multiple scattering, Coulomb correction),\allowbreak{} the radiative tails (\hyperref[param:usingu5frad]{\code{using\_\allowbreak{}rad}}\simcidx{Parameters!using_rad@\code{using\_rad}}), the two spectrometer arms and the
normalisation. The generation volume \code{genvol} is the product of the generated ranges: the electron solid angle; also the
hadron solid angle and the electron energy for every reaction except H(e,e'p); also the hadron energy for A(e,e'p) and
SIDIS.

\textbf{Physics notes.} Every formula behind a reaction (generated variables and ranges, event completion, jacobians,
the cross-section model, radiation, inputs and tree columns, validation and Fortran origin) is in its physics note, one
PDF per reaction, bound together as the \emph{SIMC++ Physics Reference}\footnote{\code{docs/physics/pdf/SIMCpp\_\allowbreak{}Physics\_\allowbreak{}Reference.\allowbreak{}pdf}; sources and build: \code{docs/physics/README.md}.}.
The common elements (the plug-in architecture, notation and glossary, weighting and normalisation, energy loss, Coulomb
correction, radiative corrections) are in note 00\footnote{\code{docs/physics/pdf/SIMCpp\_\allowbreak{}Physics\_\allowbreak{}00\_\allowbreak{}common.\allowbreak{}pdf}}. A
copy of each reaction note sits in the source folder of its physics model (\code{src/\allowbreak{}physics/\allowbreak{}<reaction>/}).

\section{\texorpdfstring{H(e,e'p) elastic, D(e,e'p), A(e,e'p) and A(e,e')A}{H(e,e'p) elastic, D(e,e'p), A(e,e'p) and A(e,e')A}}\label{user_manual-05u5freactions:52-heep-elastic-deep-aeep-and-aeea}
\begin{itemize}
\item \textbf{H(e,e'p)} (\code{hydrogen\_\allowbreak{}elastic}, model \code{sigep}): elastic scattering. Only the electron angles are generated; the
proton follows from the kinematics. Elastic e-p cross section with a proton form-factor fit. Physics note 01, H(e,e'p)
elastic\footnote{\code{docs/physics/pdf/SIMCpp\_\allowbreak{}Physics\_\allowbreak{}01\_\allowbreak{}hydrogen\_\allowbreak{}elastic.\allowbreak{}pdf}}.
\item \textbf{D(e,e'p) and A(e,e'p)} (\code{deuterium\_\allowbreak{}eep}, \code{nucleus\_\allowbreak{}eep}, model \code{deforest}): quasi-elastic knock-out with the de Forest
off-shell cross section (\hyperref[param:deForestu5fflag]{\code{deForest\_\allowbreak{}flag}}\simcidx{Parameters!deForest_flag@\code{deForest\_flag}}: σcc1, σcc2 or σcc1 on-shell) times a spectral function:

\begin{itemize}
\item \textbf{theory files} (default): \code{h2.\allowbreak{}theory}, \code{c12.\allowbreak{}theory}, \code{fe56.\allowbreak{}theory}, \code{au197.\allowbreak{}theory} in the data directory's
\code{spectral/}, chosen from A (2, 12, 56, 197). Any other A uses the C-12 file, with a WARNING;
\item \textbf{Benhar spectral functions} (\code{use\_\allowbreak{}benhar\_\allowbreak{}sf\ =\ 1}): \code{benharsf\_\allowbreak{}*.\allowbreak{}dat} (A = 3, 4, 12, 56, 64, 197; Cu-64 since upstream
SIMC a85dc6b), weighted by \hyperref[param:transparency]{\code{transparency}}\simcidx{Parameters!transparency@\code{transparency}}.
\end{itemize}
In a JSON input, \code{"pipeline":\ \{"reaction":\ \{"spectral\_\allowbreak{}function":\ \{"type":\ "benhar",\ "file":\ "benharsf\_\allowbreak{}12.\allowbreak{}dat",\ "transparency":\ 0.\allowbreak{}6\}\}\}} chooses the file explicitly (\code{-{}-{}data-{}dir}\simcidx{Command-line options!data-dir@\code{-{}-{}data-{}dir}} sets its directory). \code{cuts\%\allowbreak{}Em} limits the generated
missing energy. Physics note 02, D(e,e'p) and A(e,e'p)\footnote{\code{docs/physics/pdf/SIMCpp\_\allowbreak{}Physics\_\allowbreak{}02\_\allowbreak{}quasielastic\_\allowbreak{}eep.\allowbreak{}pdf}}.
\item \textbf{A(e,e')A nuclear elastic} (\code{doing\_\allowbreak{}nuc\_\allowbreak{}elast\ =\ 1}, \code{nuclear\_\allowbreak{}elastic}, model \code{nuc\_\allowbreak{}elastic}; upstream SIMC commit
4df7351): elastic scattering from the whole nucleus, the "hadron" is the recoiling nucleus (mass \hyperref[param:targu25massu5famu]{\code{targ\%\allowbreak{}mass\_\allowbreak{}amu}}\simcidx{Parameters!targ.mass_amu@\code{targ\%mass\_amu}}).
Only the electron angles are generated, as for H(e,e'p). Cross section σ\_Mott · (Z·F(Q²))² · recoil factor with a
charge form factor F, a fixed Gaussian (2.4 fm) times a monopole in R (R = 1.07 A\textasciicircum{}1/\allowbreak{}3 fm; de Vries radii for A = 27, 28, 56, 205), and no magnetic term (spin-0
nuclei). W and the ntuple columns \code{corrsing}, \code{Pmx/\allowbreak{}y/\allowbreak{}z} follow the (e,e'p) formulas with the proton mass (\code{corrsing}
and \code{Pmx/\allowbreak{}y/\allowbreak{}z} are 0). \hyperref[param:doingu5fnucu5felast]{\code{doing\_\allowbreak{}nuc\_\allowbreak{}elast}}\simcidx{Parameters!doing_nuc_elast@\code{doing\_nuc\_elast}} cannot be combined with another \code{doing\_\allowbreak{}*} reaction flag or \hyperref[param:doingu5fphsp]{\code{doing\_\allowbreak{}phsp}}\simcidx{Parameters!doing_phsp@\code{doing\_phsp}}.
Upstream SIMC 59ccdf6 lost this reaction in commit e4dd175 (0 successes); SIMC++ runs the code of commit 20d46fc
(FORTRAN\_\allowbreak{}DIFFERENCES U-A1). Physics note 03, A(e,e')A nuclear elastic\footnote{\code{docs/physics/pdf/SIMCpp\_\allowbreak{}Physics\_\allowbreak{}03\_\allowbreak{}nuclear\_\allowbreak{}elastic.\allowbreak{}pdf}}.
\end{itemize}
Limitations: the D(e,e'p) weight has no hadron-angle jacobian while A(e,e'p) has one (a ≤ 1.5 \% effect at the corners of
the generated solid angle); He-3 and He-4 use the C-12 spectral function unless you give a file.

\section{\texorpdfstring{Exclusive pion production (\code{doing\_\allowbreak{}pion})}{Exclusive pion production (doing\_\allowbreak{}pion)}}\label{user_manual-05u5freactions:53-exclusive-pion-production-doingu5fpion}
\code{pion\_\allowbreak{}production} with the \code{pion} model: H(e,e'π+)n, D/A(e,e'π±) quasifree (Fermi motion from momentum distributions),\allowbreak{} Δ
final states and coherent production, selected with \hyperref[param:whichu5fpion]{\code{which\_\allowbreak{}pion}}\simcidx{Parameters!which_pion@\code{which\_pion}} (\hyperref[user_manual-04u5finputs:46-parameter-reference]{4.6}). The cross
section is a parametrisation of the pion electroproduction data; below W = 2 GeV the MAID2007 tables (data directory
\code{pion/}) are also evaluated.

\begin{itemize}
\item \textbf{π0} (\code{doing\_\allowbreak{}pizero\ =\ 1}): the π0 decays to γγ. With the calorimeter arm 7/8 as hadron arm both photons are tracked
to the calorimeter,\allowbreak{} and \hyperref[param:pizerou5fngamma]{\code{pizero\_\allowbreak{}ngamma}}\simcidx{Parameters!pizero_ngamma@\code{pizero\_ngamma}} (1 or 2) says how many must hit (\hyperref[user_manual-06u5fspectrometers:63-the-calorimeter-arms-7-and-8]{6.3}).
The tree gets 12 photon columns.
\item \textbf{Decay in flight} (\code{doing\_\allowbreak{}decay\ =\ 1}, \hyperref[param:ctau]{\code{ctau}}\simcidx{Parameters!ctau@\code{ctau}}): π → μν in the spectrometer; the \code{Mhadron} column is the mass of the
particle that reached the detectors.
\item Not allowed (the run stops or is refused): π− from hydrogen, coherent production on H or D, coherent π− on ³He,
coherent production on \textsuperscript{4}He,\allowbreak{} π0 "from the neutron" (\code{which\_\allowbreak{}pion\ =\ 1}) on hydrogen.
\end{itemize}
Limitations: coherent production uses the nucleon fit with the photon-nucleus W; for the Δ final states a new
Breit-Wigner Δ mass is drawn when the incoming-electron radiation recomputes the event; the MAID amplitudes are single
precision and taken from the nearest grid point.

\section{\texorpdfstring{Kaon electroproduction (\code{doing\_\allowbreak{}kaon})}{Kaon electroproduction (doing\_\allowbreak{}kaon)}}\label{user_manual-05u5freactions:54-kaon-electroproduction-doingu5fkaon}
\code{kaon\_\allowbreak{}production} with the \code{kaon} model: K+Λ, K+Σ0, K+Σ− (on the neutron), quasifree on D and A, and coherent hypernucleus
production (\hyperref[param:whichu5fkaon]{\code{which\_\allowbreak{}kaon}}\simcidx{Parameters!which_kaon@\code{which\_kaon}} 10/11/12). The weight uses the factorized model (\code{sigcm} and \code{factor} columns); the Saghai model
is evaluated for the \code{saghai} column. Tables in the data directory's \code{kaon/}.

\begin{itemize}
\item With \code{doing\_\allowbreak{}decay\ =\ 1} the kaon decays in flight (70 \% μν, 30 \% ππ). With \code{doing\_\allowbreak{}decay\ =\ 0} the weight is multiplied by
the kaon's survival probability to the detectors, which is also written to \code{decdist}.
\end{itemize}
Limitations: a fixed decay model (branching ratios and decay momentum); coherent production uses the nucleon formulas;
the proton mass is used in the Σ− kinematics of the factorized model.

\section{\texorpdfstring{(e,e'p)π in the Delta region (\code{doing\_\allowbreak{}delta})}{(e,e'p)π in the Delta region (doing\_\allowbreak{}delta)}}\label{user_manual-05u5freactions:55-eepu3c0-in-the-delta-region-doingu5fdelta}
\code{delta\_\allowbreak{}production} with the \code{delta} model: p(e,e'p)π with the proton detected and the pion missing (the missing particle
has the charged-pion mass). D and nuclei are supported in the default mode.

\textbf{The event weight contains no cross section} (no Delta model is implemented). The weight is phase space × virtual-photon
flux × the jacobian of the proton angles. The yields, the \code{.\allowbreak{}hist} wtcontr and the \code{Weight} column of a Delta run are
therefore \textbf{not physical rates}; use Delta runs for acceptance and kinematic studies only. The \code{sigcm} column holds a
pion-production fit evaluated at the proton's CM angles, for information; it is not in the weight.

\hyperref[param:whichu5fpion]{\code{which\_\allowbreak{}pion}}\simcidx{Parameters!which_pion@\code{which\_pion}} 2 or 3 also acts on Delta runs: a Breit-Wigner Δ mass then replaces the recoil pion.

\section{\texorpdfstring{Semi-inclusive production (\code{doing\_\allowbreak{}semi})}{Semi-inclusive production (doing\_\allowbreak{}semi)}}\label{user_manual-05u5freactions:56-semi-inclusive-production-doingu5fsemi}
\code{sidis\_\allowbreak{}production} with the \code{sidis} model: (e,e'π±)X with CTEQ5M parton distributions and a 2021 fit of the pion
fragmentation functions (P. Bosted), (e,e'K±)X with fDSS fragmentation,\allowbreak{} π0 as the average of π+ and π−. Deuterium uses the
isospin-symmetric nucleon sum; A ≥ 3 adds the SLAC EMC fit and a hadron-attenuation fit. Tables in the data directory's
\code{sidis/}.

\begin{simctable}{small}{2}{>{\raggedright\arraybackslash}p{\dimexpr 0.476\linewidth-2\tabcolsep\relax}>{\raggedright\arraybackslash}p{\dimexpr 0.477\linewidth-2\tabcolsep\relax}}
%
\toprule\rowcolor{simchead}\simcth{inputs} & \simcth{hadron} \\\midrule\endfirsthead
\multicolumn{2}{@{}l}{\simccontinued}\\\toprule\rowcolor{simchead}\simcth{inputs} & \simcth{hadron} \\\midrule\endhead
\midrule\multicolumn{2}{r@{}}{\simccontinues}\\\endfoot
\bottomrule\endlastfoot
\code{doing\_\allowbreak{}semi=\allowbreak{}1\ doing\_\allowbreak{}pion=\allowbreak{}1\ doing\_\allowbreak{}hplus=\allowbreak{}1} / \code{0} & π+ / π− \\
\code{doing\_\allowbreak{}semi=\allowbreak{}1\ doing\_\allowbreak{}pion=\allowbreak{}1\ doing\_\allowbreak{}pizero=\allowbreak{}1} & π0 (with the calorimeter arm; \hyperref[param:pizerou5fngamma]{\code{pizero\_\allowbreak{}ngamma}}\simcidx{Parameters!pizero_ngamma@\code{pizero\_ngamma}} required) \\
\code{doing\_\allowbreak{}semi=\allowbreak{}1\ doing\_\allowbreak{}kaon=\allowbreak{}1\ doing\_\allowbreak{}hplus=\allowbreak{}1} / \code{0} & K+ / K− \\
\code{doing\_\allowbreak{}semi=\allowbreak{}1} alone & the run stops: \code{I\ don't\ have\ ANY\ idea\ what\ (\allowbreak{}e,\allowbreak{}e'm)X\ we're\ doing!!!} \\
\end{simctable}
Other inputs: \hyperref[param:dou5ffermi]{\code{do\_\allowbreak{}fermi}}\simcidx{Parameters!do_fermi@\code{do\_fermi}} (Fermi motion, D and A only), \hyperref[param:doingu5fdecay]{\code{doing\_\allowbreak{}decay}}\simcidx{Parameters!doing_decay@\code{doing\_decay}}/\hyperref[param:ctau]{\code{ctau}}\simcidx{Parameters!ctau@\code{ctau}} (with decay OFF the survival probability goes
to \code{decdist} and multiplies the weight for kaons only). \hyperref[param:ptu5fbu5fparam]{\code{pt\_\allowbreak{}b\_\allowbreak{}param}}\simcidx{Parameters!pt_b_param@\code{pt\_b\_param}} and \code{sigc\_\allowbreak{}*} are read and ignored. Events where the
parton distributions vanish, or where Fermi motion gives x\_Bj ≤ 0, get weight 0. \hyperref[param:doingu5fsemi]{\code{doing\_\allowbreak{}semi}}\simcidx{Parameters!doing_semi@\code{doing\_semi}} with \hyperref[param:doingu5fphsp]{\code{doing\_\allowbreak{}phsp}}\simcidx{Parameters!doing_phsp@\code{doing\_phsp}} is refused.

Limitations: the 2π / πK threshold uses an approximate missing-mass formula, also for kaons and with Fermi motion; with
radiation the generated hadron-energy range is wide; the hadron-attenuation factor for deuterium is not 1.

\section{\texorpdfstring{Exclusive ρ0 production (\code{doing\_\allowbreak{}rho})}{Exclusive ρ0 production (doing\_\allowbreak{}rho)}}\label{user_manual-05u5freactions:57-exclusive-u3c10-production-doingu5frho}
\code{rho\_\allowbreak{}production} with the \code{rho} model: p(e,e'ρ0)p, ρ0 → π+π−, with one pion detected (\hyperref[param:doingu5fhplus]{\code{doing\_\allowbreak{}hplus}}\simcidx{Parameters!doing_hplus@\code{doing\_hplus}} chooses π+ or π−). The ρ
mass is drawn from a Breit-Wigner and the decay angles from the ρ decay distribution; the \code{Mrho} and \code{Thrho} columns record
them.

\begin{itemize}
\item Targets: H; D and ³He with a nucleon at rest (no Fermi motion, no ³He spectral function). A ≥ 4 is refused.
\item Hadron arm 1, 2, 3, 4 or 5; the calorimeter is refused. \hyperref[param:doingu5frho]{\code{doing\_\allowbreak{}rho}}\simcidx{Parameters!doing_rho@\code{doing\_rho}} together with another reaction flag or with
\hyperref[param:doingu5fpizero]{\code{doing\_\allowbreak{}pizero}}\simcidx{Parameters!doing_pizero@\code{doing\_pizero}} is refused.
\item The ρ is generated over 4π, so the efficiency is low (about 3e-4 accepted events per try): runs of 1e5 events take hours
(use \code{-{}-{}threads}\simcidx{Command-line options!threads@\code{-{}-{}threads}}).
\end{itemize}
Limitations: s without the Fermi terms, t at the ρ pole mass, decay angles in the lab-axis frame rather than the helicity
frame, radiation of the neutral ρ (most of its tail-3 tries fail), and a new ρ drawn when the incoming-electron radiation
recomputes the event.

\section{\texorpdfstring{Phase-space and acceptance runs (\code{doing\_\allowbreak{}phsp})}{Phase-space and acceptance runs (doing\_\allowbreak{}phsp)}}\label{user_manual-05u5freactions:58-phase-space-and-acceptance-runs-doingu5fphsp}
\code{doing\_\allowbreak{}phsp\ =\ 1} sets the cross section to 1 and normfac to 1, for acceptance studies: the weight is then the generation
weight only. Two properties to know:

\begin{itemize}
\item \textbf{On hydrogen the kinematics stay elastic.} The proton is still computed from the elastic kinematics, so the run is an
elastic acceptance run, not a flat phase space.
\item \textbf{Radiation stays on} if \code{using\_\allowbreak{}rad\ =\ 1}, with the photon limits computed for a target mass of 0. At high energy this
can give zero accepted events. Set \code{using\_\allowbreak{}rad\ =\ 0} for a pure acceptance run.
\end{itemize}
\hyperref[param:doingu5fphsp]{\code{doing\_\allowbreak{}phsp}}\simcidx{Parameters!doing_phsp@\code{doing\_phsp}} switches off the pion, kaon, Delta and rho physics but keeps the hadron mass that these flags selected.

To measure the acceptance of one spectrometer,\allowbreak{} or to throw events at a spectrometer without simulating it, see
\hyperref[param:spectu5fmode]{\code{spect\_\allowbreak{}mode}}\simcidx{Parameters!spect_mode@\code{spect\_mode}} (\hyperref[user_manual-06u5fspectrometers:64-spectu5fmode-switching-the-spectrometer-monte-carlo-off]{6.4}).

(How each reaction behaves in the legacy-compatible mode, and what the default mode corrects, is listed in
\code{docs/\allowbreak{}translation/\allowbreak{}markdown/\allowbreak{}user\_\allowbreak{}differences.\allowbreak{}md}.)


% docs/manual/latex/user/06_spectrometers.tex - chapter of the LaTeX edition of the book "user_manual" (SIMC++ documentation).
%   Converted once from docs/manual/markdown/user/06_spectrometers.md by tools/md2latex.py; from then on THIS FILE is the source of the
%   LaTeX edition and may be edited by hand. Re-running tools/md2latex.py --write-sources for this
%   chapter OVERWRITES the hand edits (see docs/README.md).
%   Compile: cmake --build <build> --target pdf, or tools/build_docs_pdf.py (figures, xelatex passes,
%   makeindex with docs/common/latex/simcpp.ist); by hand, see docs/README.md. XeLaTeX, class
%   docs/common/latex/simcpp.cls (packages: fontspec, tikz, xcolor, hyperref, listings, longtable, booktabs,
%   titlesec, titletoc, fancyhdr, makeidx, preview, ...); fonts TeX Gyre Pagella, TeX Gyre Heros, DejaVu Sans
%   Mono. The figures are built first from docs/{manual,translation}/figures/*.tex (preview.sty) into PDFs.
\chapter{\texorpdfstring{Spectrometers}{Spectrometers}}\label{user_manual-06u5fspectrometers}
\section{\texorpdfstring{Spectrometer numbers}{Spectrometer numbers}}\label{user_manual-06u5fspectrometers:61-spectrometer-numbers}
\hyperref[param:electronu5farm]{\code{electron\_\allowbreak{}arm}}\simcidx{Parameters!electron_arm@\code{electron\_arm}} and \hyperref[param:hadronu5farm]{\code{hadron\_\allowbreak{}arm}}\simcidx{Parameters!hadron_arm@\code{hadron\_arm}} take \textbf{one number per spectrometer} (\code{include/\allowbreak{}simc/\allowbreak{}core/\allowbreak{}Spectrometer\allowbreak{}Id.\allowbreak{}h}):

\begin{simctable}{footnotesize}{5}{l>{\raggedright\arraybackslash}p{\dimexpr 0.347\linewidth-2\tabcolsep\relax}>{\raggedright\arraybackslash}p{\dimexpr 0.169\linewidth-2\tabcolsep\relax}>{\raggedright\arraybackslash}p{\dimexpr 0.084\linewidth-2\tabcolsep\relax}>{\raggedright\arraybackslash}p{\dimexpr 0.245\linewidth-2\tabcolsep\relax}}
%
\toprule\rowcolor{simchead}\simcth{code} & \simcth{spectrometer} & \simcth{side of the beam} & \simcth{plug-in} & \simcth{stop counters (\code{.\allowbreak{}geni}, summary)} \\\midrule\endfirsthead
\multicolumn{5}{@{}l}{\simccontinued}\\\toprule\rowcolor{simchead}\simcth{code} & \simcth{spectrometer} & \simcth{side of the beam} & \simcth{plug-in} & \simcth{stop counters (\code{.\allowbreak{}geni}, summary)} \\\midrule\endhead
\midrule\multicolumn{5}{r@{}}{\simccontinues}\\\endfoot
\bottomrule\endlastfoot
1 & HMS & right & \code{hms} & \code{hms} \\
2 & SOS & left & \code{sos} & \code{sos} \\
3 & HRS-R (Hall A) & right & \code{hrsr} & \code{hrsr} \\
4 & HRS-L (Hall A) & left & \code{hrsl} & \code{hrsl} \\
5 & SHMS & left & \code{shms} & \code{shms} \\
7 & calorimeter (NPS-like), HMS side & right & \code{calo} & \code{calo} \\
8 & calorimeter,\allowbreak{} SHMS side & left & \code{calo} & \code{calo} \\
\end{simctable}
\begin{itemize}
\item \textbf{6 is rejected} when the input is read: \code{simc++:\ input\ error:\ .\allowbreak{}.\allowbreak{}.:\ electron\_\allowbreak{}arm\ =\ 6\ was\ the\ retired\ SHMS\ LSA\ tune;\ use\ 5\ =\ SHMS} (exit status 2).
\item Any other number stops the run: \code{I\ dont\ know\ what\ phi\ should\ be\ for\ the\ electron\ arm} (exit status 3).
\item The side of the beam fixes the azimuth of the spectrometer (beam right φ = 3π/2, left φ = π/2) and which tree prefix
holds which particle (\hyperref[user_manual-07u5foutputs:712-arm-prefixes-hs-and-ss]{7.1.2}). The angles \hyperref[param:specu25eu25theta]{\code{spec\%\allowbreak{}e\%\allowbreak{}theta}}\simcidx{Parameters!spec.e.theta@\code{spec\%e\%theta}}, \hyperref[param:specu25pu25theta]{\code{spec\%\allowbreak{}p\%\allowbreak{}theta}}\simcidx{Parameters!spec.p.theta@\code{spec\%p\%theta}} are
taken as absolute values.
\end{itemize}
\section{\texorpdfstring{The magnetic spectrometers}{The magnetic spectrometers}}\label{user_manual-06u5fspectrometers:62-the-magnetic-spectrometers}
For the HMS, SHMS, SOS, HRS-R and HRS-L, each event goes through:

\begin{itemize}
\item the transport through each magnetic element with COSY optics maps, with aperture checks at the element entrance,
middle and exit;
\item decay in flight (π, K) when \code{doing\_\allowbreak{}decay\ =\ 1};
\item multiple scattering in the materials and the wire-chamber resolution when \code{mc\_\allowbreak{}smear\ =\ 1};
\item the detector-hut apertures (chambers, hodoscopes/\allowbreak{}scintillators,\allowbreak{} calorimeter);
\item the reconstruction of delta, yptar, xptar and ytar from the smeared focal plane with the reconstruction matrix (with the
raster position if \code{correct\_\allowbreak{}raster\ =\ 1}).
\end{itemize}
The optics maps are in the data directory's \code{matrices/\allowbreak{}simc\_\allowbreak{}fortran/} (md5-checked). Where a try is lost is counted in the
stop counters (\hyperref[user_manual-07u5foutputs:73-geni-gen-and-genroot]{7.3}).

Collimators: \hyperref[param:usingu5fHMScoll]{\code{using\_\allowbreak{}HMScoll}}\simcidx{Parameters!using_HMScoll@\code{using\_HMScoll}} / \hyperref[param:usingu5fSHMScoll]{\code{using\_\allowbreak{}SHMScoll}}\simcidx{Parameters!using_SHMScoll@\code{using\_SHMScoll}} step pions and muons through the HMS / SHMS collimator (absorption,\allowbreak{} energy
loss, scattering). The \code{.\allowbreak{}geni} then has a \code{Collimator\ stops\ (\allowbreak{}pi/\allowbreak{}mu)} line, and each slit counter counts a ray at most once.

\begin{figure}[htbp]
\simcfigure{spectrometer_arm}
\caption[One spectrometer arm]{One spectrometer arm: from the target through the slit or collimator, the magnetic elements and the detector hut to the reconstructed track; a track that leaves an aperture stops and is counted}\label{fig:spectrometeru5farm}
\end{figure}
\section{\texorpdfstring{The calorimeter arms 7 and 8}{The calorimeter arms 7 and 8}}\label{user_manual-06u5fspectrometers:63-the-calorimeter-arms-7-and-8}
The calorimeter arm (\code{calo}) is a straight drift to a calorimeter face at \hyperref[param:driftu5ftou5fcal]{\code{drift\_\allowbreak{}to\_\allowbreak{}cal}}\simcidx{Parameters!drift_to_cal@\code{drift\_to\_cal}} cm from the target, with the
calorimeter aperture. It has no magnetic field, no smearing and no reconstruction (RECON = the values at the spectrometer
entrance, after the target). It is used for:

\begin{itemize}
\item \textbf{π0 → γγ} (\hyperref[param:doingu5fpizero]{\code{doing\_\allowbreak{}pizero}}\simcidx{Parameters!doing_pizero@\code{doing\_pizero}} with hadron arm 7 or 8): both photons are tracked to the face. \code{pizero\_\allowbreak{}ngamma\ =\ 1} accepts the
event when at least one photon hits, \code{2} when both hit. The tree gets the photon positions on the face (\code{xcal\_\allowbreak{}gamma1/\allowbreak{}2},
\code{ycal\_\allowbreak{}gamma1/\allowbreak{}2}, −1e10 for a miss) and the photon four-momenta. Photons are massless and never decay.
\item \textbf{an electron in the calorimeter} (\hyperref[param:electronu5farm]{\code{electron\_\allowbreak{}arm}}\simcidx{Parameters!electron_arm@\code{electron\_arm}} 7 or 8).
\end{itemize}
Set \hyperref[param:driftu5ftou5fcal]{\code{drift\_\allowbreak{}to\_\allowbreak{}cal}}\simcidx{Parameters!drift_to_cal@\code{drift\_to\_cal}}: it is not checked, and 0 puts the face at the target. The target windows seen by particles going to
the calorimeter are not modelled specifically (the values of the previous particle's trip through the target are reused).

\section{\texorpdfstring{\code{spect\_\allowbreak{}mode}: switching the spectrometer Monte Carlo off}{spect\_\allowbreak{}mode: switching the spectrometer Monte Carlo off}}\label{user_manual-06u5fspectrometers:64-spectu5fmode-switching-the-spectrometer-monte-carlo-off}
This answers the question "do I simulate the acceptance, or do I throw events at the spectrometer?":

\begin{simctable}{small}{3}{>{\raggedright\arraybackslash}p{\dimexpr 0.147\linewidth-2\tabcolsep\relax}>{\raggedright\arraybackslash}p{\dimexpr 0.199\linewidth-2\tabcolsep\relax}>{\raggedright\arraybackslash}p{\dimexpr 0.199\linewidth-2\tabcolsep\relax}}
%
\toprule\rowcolor{simchead}\simcth{\code{spect\_\allowbreak{}mode}} & \simcth{electron arm} & \simcth{hadron arm} \\\midrule\endfirsthead
\multicolumn{3}{@{}l}{\simccontinued}\\\toprule\rowcolor{simchead}\simcth{\code{spect\_\allowbreak{}mode}} & \simcth{electron arm} & \simcth{hadron arm} \\\midrule\endhead
\midrule\multicolumn{3}{r@{}}{\simccontinues}\\\endfoot
\bottomrule\endlastfoot
0 (default) & full Monte Carlo & full Monte Carlo \\
−1 & \textbf{off} & full Monte Carlo \\
−2 & full Monte Carlo & \textbf{off} \\
1 & \textbf{off} & \textbf{off} \\
\end{simctable}
\begin{itemize}
\item \textbf{Arm on}: the event is transported through the spectrometer (apertures, decay, multiple scattering, resolution,
reconstruction). Only events that the spectrometer accepts survive. This is the normal "simulate the measured yield" mode.
\item \textbf{Arm off}: no spectrometer is simulated for that arm. The event is "thrown at the spectrometer" and kept whatever the
spectrometer would have done. The only limits for that arm are the generation limits (the \code{SPedge} window, without the
spectrometer's slop) and the \code{SPedge} window that decides whether an event counts in wtcontr. In the tree, the arm's
\code{*delta}, \code{*yptar}, \code{*xptar} hold the values at the spectrometer entrance (after the target energy loss and multiple
scattering, no resolution). Its focal-plane columns \code{*xfp}, \code{*xpfp}, \code{*yfp}, \code{*ypfp} and its \code{*ytar} are \textbf{0.0 in every
event}: they are not measurements,\allowbreak{} so leave them out of any analysis of such a run. Its \code{.\allowbreak{}geni} counters are 0.
\end{itemize}
So \code{spect\_\allowbreak{}mode\ =\ 1} with \code{doing\_\allowbreak{}phsp\ =\ 1} gives the pure generated phase space, \code{spect\_\allowbreak{}mode\ =\ 0} with \code{doing\_\allowbreak{}phsp\ =\ 1} the
acceptance of both spectrometers (\hyperref[user_manual-05u5freactions:58-phase-space-and-acceptance-runs-doingu5fphsp]{5.8}), and −1/−2 the
acceptance of one spectrometer with the other arm open.

Any other value means "both arms on"; a WARNING is printed.

\section{\texorpdfstring{Polarized-target field (\code{using\_\allowbreak{}tgt\_\allowbreak{}field})}{Polarized-target field (using\_\allowbreak{}tgt\_\allowbreak{}field)}}\label{user_manual-06u5fspectrometers:65-polarized-target-field-usingu5ftgtu5ffield}
With \code{using\_\allowbreak{}tgt\_\allowbreak{}field\ =\ 1}, the particles are tracked through the magnetic field of a polarized target, from the vertex
to the field-free region on the way to the spectrometer,\allowbreak{} and back for the reconstruction. The field map comes from
\hyperref[param:tgtu5ffieldu5ffile]{\code{tgt\_\allowbreak{}field\_\allowbreak{}file}}\simcidx{Parameters!tgt_field_file@\code{tgt\_field\_file}} (\code{'trg\_\allowbreak{}field\_\allowbreak{}map.\allowbreak{}dat'} is in the data directory's \code{field/}; \code{'0'} = no field; \code{'\ '} = a uniform 5 T test
field). Its orientation comes from \hyperref[param:targu5fBangle]{\code{targ\_\allowbreak{}Bangle}}\simcidx{Parameters!targ_Bangle@\code{targ\_Bangle}} and \hyperref[param:targu5fBphi]{\code{targ\_\allowbreak{}Bphi}}\simcidx{Parameters!targ_Bphi@\code{targ\_Bphi}}. The tree gets 8 more columns with the angles between q,
the scattering plane and the polarization (\hyperref[user_manual-07u5foutputs:717-polarized-target-columns-usingu5ftgtu5ffield]{7.1.7}).

\begin{itemize}
\item The calorimeter arms 7/8 cannot be used with the field (the run stops).
\item Approximations to know: the polarization vector is in-plane towards beam left (\hyperref[param:targu5fBphi]{\code{targ\_\allowbreak{}Bphi}}\simcidx{Parameters!targ_Bphi@\code{targ\_Bphi}} and the sign of \hyperref[param:targu5fBangle]{\code{targ\_\allowbreak{}Bangle}}\simcidx{Parameters!targ_Bangle@\code{targ\_Bangle}}
are not used for it); the tracking step limit depends on the side of the arm (about 2.5 \% of the SHMS events of the
polarized-target pion example stop early); the bending constant is 90/E (0.14 \% stronger than c = 89.88).
\end{itemize}
\section{\texorpdfstring{π0 in the HMS or SHMS}{π0 in the HMS or SHMS}}\label{user_manual-06u5fspectrometers:66-u3c00-in-the-hms-or-shms}
With \hyperref[param:doingu5fpizero]{\code{doing\_\allowbreak{}pizero}}\simcidx{Parameters!doing_pizero@\code{doing\_pizero}} and hadron arm 1 or 5, the π0 itself is transported through the magnetic spectrometer as if it were
charged (apertures, optics maps, multiple scattering),\allowbreak{} and RECON = the values at the spectrometer entrance. Read such runs
as "π0 kinematics inside the HMS/SHMS acceptance",\allowbreak{} not as a detector simulation. The π0 does not decay in flight.



\simcpart{Results and reproducibility}
  {What a run writes, how to reproduce a run exactly, and how to make runs faster.}
% docs/manual/latex/user/07_outputs.tex - chapter of the LaTeX edition of the book "user_manual" (SIMC++ documentation).
%   Converted once from docs/manual/markdown/user/07_outputs.md by tools/md2latex.py; from then on THIS FILE is the source of the
%   LaTeX edition and may be edited by hand. Re-running tools/md2latex.py --write-sources for this
%   chapter OVERWRITES the hand edits (see docs/README.md).
%   Compile: cmake --build <build> --target pdf, or tools/build_docs_pdf.py (figures, xelatex passes,
%   makeindex with docs/common/latex/simcpp.ist); by hand, see docs/README.md. XeLaTeX, class
%   docs/common/latex/simcpp.cls (packages: fontspec, tikz, xcolor, hyperref, listings, longtable, booktabs,
%   titlesec, titletoc, fancyhdr, makeidx, preview, ...); fonts TeX Gyre Pagella, TeX Gyre Heros, DejaVu Sans
%   Mono. The figures are built first from docs/{manual,translation}/figures/*.tex (preview.sty) into PDFs.
\chapter{\texorpdfstring{Outputs}{Outputs}}\label{user_manual-07u5foutputs}
\simcabstract{A run writes into the output directory (\code{-{}o\ DIR}, default \code{.}), with \code{<run>} = the input file name without its extension
(or \code{-{}-{}run-{}name}\simcidx{Command-line options!run-name@\code{-{}-{}run-{}name}}):}
\begin{simctable}{small}{3}{>{\raggedright\arraybackslash}p{\dimexpr 0.250\linewidth-2\tabcolsep\relax}>{\raggedright\arraybackslash}p{\dimexpr 0.209\linewidth-2\tabcolsep\relax}>{\raggedright\arraybackslash}p{\dimexpr 0.473\linewidth-2\tabcolsep\relax}}
%
\toprule\rowcolor{simchead}\simcth{file} & \simcth{written by (sink)} & \simcth{content} \\\midrule\endfirsthead
\multicolumn{3}{@{}l}{\simccontinued}\\\toprule\rowcolor{simchead}\simcth{file} & \simcth{written by (sink)} & \simcth{content} \\\midrule\endhead
\midrule\multicolumn{3}{r@{}}{\simccontinues}\\\endfoot
\bottomrule\endlastfoot
\code{<run>\allowbreak{}.\allowbreak{}root} & \code{root\_\allowbreak{}tree} & tree \code{h10}, one entry per accepted event; only if \code{Nntu\ >\ 0} \\
\code{<run>\allowbreak{}.\allowbreak{}hist} & \code{fortran\_\allowbreak{}hist} & the run report \\
\code{<run>\allowbreak{}.\allowbreak{}geni} & \code{fortran\_\allowbreak{}hist} & spectrometer stop counters \\
\code{<run>\allowbreak{}.\allowbreak{}gen} & \code{fortran\_\allowbreak{}hist} & diagnostic histograms (text) \\
\code{<run>\allowbreak{}.\allowbreak{}gen.\allowbreak{}root} & \code{gen\_\allowbreak{}hist} & the same histograms as ROOT TH1D \\
\code{<run>\allowbreak{}.\allowbreak{}summary.\allowbreak{}json}, \code{.\allowbreak{}summary.\allowbreak{}txt} & \code{run\_\allowbreak{}summary} & run summary (schema \code{docs/\allowbreak{}manual/\allowbreak{}markdown/\allowbreak{}reference/\allowbreak{}summary.\allowbreak{}schema.\allowbreak{}json}) \\
\end{simctable}
\code{ngen\ =\ 0} writes nothing. The detailed reference of every value is \code{docs/\allowbreak{}manual/\allowbreak{}markdown/\allowbreak{}reference/\allowbreak{}OUTPUT\_\allowbreak{}FORMAT.\allowbreak{}md}.

\section{\texorpdfstring{The event tree \code{h10}}{The event tree h10}}\label{user_manual-07u5foutputs:71-the-event-tree-h10}
\begin{itemize}
\item One entry per \textbf{accepted} event.
\item Every branch is a \code{Float\_\allowbreak{}t} (computed in double precision, stored as float). The branch names and order are those of
the classic SIMC ntuple, so existing analysis scripts work unchanged.
\item The layout depends on the reaction (\code{NtupleColumnsFor}):
\end{itemize}
\begin{simctable}{small}{3}{>{\raggedright\arraybackslash}p{\dimexpr 0.584\linewidth-2\tabcolsep\relax}>{\raggedright\arraybackslash}p{\dimexpr 0.136\linewidth-2\tabcolsep\relax}>{\raggedright\arraybackslash}p{\dimexpr 0.212\linewidth-2\tabcolsep\relax}}
%
\toprule\rowcolor{simchead}\simcth{run} & \simcth{columns} & \simcth{section} \\\midrule\endfirsthead
\multicolumn{3}{@{}l}{\simccontinued}\\\toprule\rowcolor{simchead}\simcth{run} & \simcth{columns} & \simcth{section} \\\midrule\endhead
\midrule\multicolumn{3}{r@{}}{\simccontinues}\\\endfoot
\bottomrule\endlastfoot
H/\allowbreak{}D/\allowbreak{}A(e,\allowbreak{}e'p) and \hyperref[param:doingu5fphsp]{\code{doing\_\allowbreak{}phsp}}\simcidx{Parameters!doing_phsp@\code{doing\_phsp}} & 46 & \hyperref[user_manual-07u5foutputs:713-eep-layout-46-columns]{7.1.3} \\
pion (\hyperref[param:doingu5fpion]{\code{doing\_\allowbreak{}pion}}\simcidx{Parameters!doing_pion@\code{doing\_pion}}), Delta (\hyperref[param:doingu5fdelta]{\code{doing\_\allowbreak{}delta}}\simcidx{Parameters!doing_delta@\code{doing\_delta}}) & 53; π0: 65 & \hyperref[user_manual-07u5foutputs:714-pion-kaon-and-delta-layouts]{7.1.4} \\
kaon (\hyperref[param:doingu5fkaon]{\code{doing\_\allowbreak{}kaon}}\simcidx{Parameters!doing_kaon@\code{doing\_kaon}}) & 55 & \hyperref[user_manual-07u5foutputs:714-pion-kaon-and-delta-layouts]{7.1.4} \\
SIDIS (\hyperref[param:doingu5fsemi]{\code{doing\_\allowbreak{}semi}}\simcidx{Parameters!doing_semi@\code{doing\_semi}}) & 56; π0: 68 & \hyperref[user_manual-07u5foutputs:715-sidis-layout-56-columns]{7.1.5} \\
rho (\hyperref[param:doingu5frho]{\code{doing\_\allowbreak{}rho}}\simcidx{Parameters!doing_rho@\code{doing\_rho}}) & 59 & \hyperref[user_manual-07u5foutputs:716-rho-layout-59-columns]{7.1.6} \\
any of the above except (e,e'p), with \hyperref[param:usingu5ftgtu5ffield]{\code{using\_\allowbreak{}tgt\_\allowbreak{}field}}\simcidx{Parameters!using_tgt_field@\code{using\_tgt\_field}} & + 8 & \hyperref[user_manual-07u5foutputs:717-polarized-target-columns-usingu5ftgtu5ffield]{7.1.7} \\
\end{simctable}
\subsection{\texorpdfstring{Weights and normalisation}{Weights and normalisation}}\label{user_manual-07u5foutputs:711-weights-and-normalisation}
\begin{itemize}
\item \code{Weight} is the event weight computed for each event: the cross section of the event times the generation
weight (radiative weight, jacobians, spectral-function or survival factors, depending on the reaction).
\item The yield of one event for the simulated charge \hyperref[param:EXPERu25charge]{\code{EXPER\%\allowbreak{}charge}}\simcidx{Parameters!EXPER.charge@\code{EXPER\%charge}} is \code{Weight\ ×\ normfac\ /\ nevent}. \code{normfac} is printed in
the \code{.\allowbreak{}hist} and the summary, and \code{nevent} = Ncontribute when ngen > 0, Ntried when ngen < 0.
\code{normfac\ =\ luminosity\ ×\ genvol\ /\ Ntried\ ×\ nevent}, and 1 for \hyperref[param:doingu5fphsp]{\code{doing\_\allowbreak{}phsp}}\simcidx{Parameters!doing_phsp@\code{doing\_phsp}}.
\item The \code{.\allowbreak{}hist} \code{wtcontr} is the sum of these yields over the accepted events that pass the \code{SPedge} window
(pass\_\allowbreak{}cuts): the expected number of counts.
\item \code{sigcc} / \code{siglab} is the cross section of the event (units in the tables below).
\end{itemize}
\subsection{\texorpdfstring{Arm prefixes \code{hs*} and \code{ss*}}{Arm prefixes hs* and ss*}}\label{user_manual-07u5foutputs:712-arm-prefixes-hs-and-ss}
As in results\_\allowbreak{}write.\allowbreak{}f,\allowbreak{} \code{hs*} is the spectrometer on the beam \textbf{right} and \code{ss*} the one on the \textbf{left}:

\begin{itemize}
\item \hyperref[param:electronu5farm]{\code{electron\_\allowbreak{}arm}}\simcidx{Parameters!electron_arm@\code{electron\_arm}} = 1, 3, 7 (HMS, HRS-R, calorimeter HMS side): \code{hs*} = electron, \code{ss*} = hadron;
\item \hyperref[param:electronu5farm]{\code{electron\_\allowbreak{}arm}}\simcidx{Parameters!electron_arm@\code{electron\_arm}} = 2, 4, 5, 8 (SOS, HRS-L, SHMS, calorimeter SHMS side): \code{hs*} = hadron, \code{ss*} = electron.
\end{itemize}
Slopes (\code{*yptar}, \code{*xptar}, \code{*xpfp}, \code{*ypfp}) are in radians (dx/dz); deltas are in \%.

\subsection{\texorpdfstring{(e,e'p) layout (46 columns)}{(e,e'p) layout (46 columns)}}\label{user_manual-07u5foutputs:713-eep-layout-46-columns}
\begin{simctable}{small}{4}{>{\raggedright\arraybackslash}p{\dimexpr 0.073\linewidth-2\tabcolsep\relax}>{\raggedright\arraybackslash}p{\dimexpr 0.247\linewidth-2\tabcolsep\relax}>{\raggedright\arraybackslash}p{\dimexpr 0.443\linewidth-2\tabcolsep\relax}>{\raggedright\arraybackslash}p{\dimexpr 0.148\linewidth-2\tabcolsep\relax}}
\caption{(e,e'p) layout (46 columns)}\\%
\toprule\rowcolor{simchead}\simcth{\#} & \simcth{branches} & \simcth{content} & \simcth{units} \\\midrule\endfirsthead
\multicolumn{4}{@{}l}{\simccontinued}\\\toprule\rowcolor{simchead}\simcth{\#} & \simcth{branches} & \simcth{content} & \simcth{units} \\\midrule\endhead
\midrule\multicolumn{4}{r@{}}{\simccontinues}\\\endfoot
\bottomrule\endlastfoot
1-4 & \code{hsdelta\ hsyptar\ hsxptar\ hsytar} & right arm, reconstructed & \%, rad, rad, cm \\
5-8 & \code{hsxfp\ hsxpfp\ hsyfp\ hsypfp} & right arm focal plane & cm, rad, cm, rad \\
9-12 & \code{hsdeltai\ hsyptari\ hsxptari\ hsytari} & right arm, generated (at the vertex, before the spectrometer) & \%, rad, rad, cm \\
13-24 & \code{ssdelta\ .\allowbreak{}.\allowbreak{}.\ ssytari} & the same for the left arm &  \\
25-29 & \code{q\ nu\ Q2\ W\ epsilon} & reconstructed kinematics & GeV/c, GeV, GeV², GeV, — \\
30-31 & \code{Em\ Pm} & missing energy and momentum & GeV, GeV/c \\
32-33 & \code{thetapq\ phipq} & angles of the proton w.r.t. q & rad \\
34 & \code{corrsing} & e' momentum minus the elastic value & GeV \\
35-37 & \code{Pmx\ Pmy\ Pmz} & missing-momentum components ("Heepcheck" frame) & GeV/c \\
38-40 & \code{PmPar\ PmPer\ PmOop} & missing momentum parallel, perpendicular,\allowbreak{} out of plane & GeV/c \\
41 & \code{fry} & −(raster vertical position) & cm \\
42 & \code{radphot} & radiated photon energy & GeV \\
43 & \code{sigcc} & cross section (µb/sr for H, µb/MeV/sr² for D/A) &  \\
44 & \code{Weight} & event weight &  \\
45-46 & \code{theta\_\allowbreak{}e\ theta\_\allowbreak{}p} & reconstructed scattering angles & rad \\
\end{simctable}
\subsection{\texorpdfstring{Pion, kaon and Delta layouts}{Pion, kaon and Delta layouts}}\label{user_manual-07u5foutputs:714-pion-kaon-and-delta-layouts}
Columns 1-33 as above, then:

\begin{simctable}{small}{4}{>{\raggedright\arraybackslash}p{\dimexpr 0.073\linewidth-2\tabcolsep\relax}>{\raggedright\arraybackslash}p{\dimexpr 0.227\linewidth-2\tabcolsep\relax}>{\raggedright\arraybackslash}p{\dimexpr 0.430\linewidth-2\tabcolsep\relax}>{\raggedright\arraybackslash}p{\dimexpr 0.181\linewidth-2\tabcolsep\relax}}
\caption{Pion, kaon and Delta layouts}\\%
\toprule\rowcolor{simchead}\simcth{\#} & \simcth{branch} & \simcth{content} & \simcth{units} \\\midrule\endfirsthead
\multicolumn{4}{@{}l}{\simccontinued}\\\toprule\rowcolor{simchead}\simcth{\#} & \simcth{branch} & \simcth{content} & \simcth{units} \\\midrule\endhead
\midrule\multicolumn{4}{r@{}}{\simccontinues}\\\endfoot
\bottomrule\endlastfoot
34 & \code{missmass} & missing mass (on the struck nucleon) & GeV \\
35 & \code{mmnuc} & missing mass on the nucleus & GeV \\
36 & \code{phad} & hadron momentum & GeV/c \\
37 & \code{t} & −t (SIMC sign), reconstructed & GeV² \\
38-40 & \code{pmpar\ pmper\ pmoop} & missing-momentum components & GeV/c \\
41 & \code{fry} & −(raster vertical position) & cm \\
42 & \code{radphot} & radiated photon energy & GeV \\
43 & \code{pfermi} & Fermi momentum along q (signed) & GeV/c \\
44 & \code{siglab} & lab cross section & µb/MeV/sr² \\
45 & \code{sigcm} & CM cross section & model units \\
46 & \code{Weight} & event weight &  \\
47 & \code{decdist} & decay distance (cm); for kaons with decay OFF the survival probability &  \\
48 & \code{Mhadron} & mass of the detected particle (μ or π after a decay) & MeV \\
49 & \code{pdotqhat} & Fermi momentum · q̂ & GeV/c \\
50-53 & \code{Q2i\ Wi\ ti\ phipqi} & vertex (generated) Q², W, t, φpq & GeV², GeV, GeV², rad \\
54-65 & \code{xcal\_\allowbreak{}gamma1\ ycal\_\allowbreak{}gamma1\ Egamma1\ Pgamma1x\ Pgamma1y\ Pgamma1z} and the same for photon 2 & π0 runs only: photon positions on the calorimeter (cm; −1e10 = missed, 0 with HMS/SHMS) and lab photon energy and momentum (MeV) &  \\
\end{simctable}
Kaon runs add 54 \code{saghai} (Saghai-model cross section) and 55 \code{factor}
(factorized model, = \code{sigcm}). Delta runs use the 53 pion columns: \code{missmass} is the missing pion mass, \code{phad} the proton
momentum, \code{siglab} and \code{Weight} contain \textbf{no cross section} (\hyperref[user_manual-05u5freactions:55-eepu3c0-in-the-delta-region-doingu5fdelta]{5.5}),
and \code{decdist}, \code{Mhadron} are 0 and Mp.

\subsection{\texorpdfstring{SIDIS layout (56 columns)}{SIDIS layout (56 columns)}}\label{user_manual-07u5foutputs:715-sidis-layout-56-columns}
Columns 1-33 as (e,e'p), then: 34 \code{missmass} (GeV), 35 \code{ppi} (hadron momentum, GeV/c), 36 \code{t} (GeV²), 37 \code{fry}, 38
\code{radphot}, 39 \code{siglab} (µb/\allowbreak{}MeV²/\allowbreak{}sr²),\allowbreak{} 40 \code{sigcent} (always 0), 41 \code{Weight}, 42 \code{decdist} (cm, or the survival probability
with decay OFF), 43 \code{Mhadron} (MeV), 44-45 \code{z\ zi}, 46-47 \code{pt2\ pt2i} (GeV²), 48-49 \code{xbj\ xbji}, 50 \code{thqi} (rad), 51 \code{sighad}
(GeV\textsuperscript{−}²),\allowbreak{} 52 \code{jacobian}, 53 \code{centjac} (always 0), 54 \code{pfermi} (GeV/c), 55 \code{xfermi}, 56 \code{phipqi} (rad). π0 SIDIS appends the 12 photon columns (57-68). The \code{i}
suffix means vertex (generated) values.

\subsection{\texorpdfstring{Rho layout (59 columns)}{Rho layout (59 columns)}}\label{user_manual-07u5foutputs:716-rho-layout-59-columns}
The 56 SIDIS columns with the rho meanings, then 57 \code{Mrho} (Breit-Wigner ρ mass of the event, MeV), 58 \code{Thrho} (decay
angle in the ρ rest frame w.r.t. the lab z axis, rad) and 59 \code{mmnuc} (GeV). The hadron-arm columns describe the detected
decay pion; \code{t} is the pion t; \code{siglab} (µb/MeV/sr) is the ρ cross section; \code{sigcent}, \code{zi}, \code{pt2i}, \code{xfermi} are 0.

\subsection{\texorpdfstring{Polarized-target columns (\code{using\_\allowbreak{}tgt\_\allowbreak{}field})}{Polarized-target columns (using\_\allowbreak{}tgt\_\allowbreak{}field)}}\label{user_manual-07u5foutputs:717-polarized-target-columns-usingu5ftgtu5ffield}
Eight columns follow \code{phipqi} (pion/\allowbreak{}kaon/\allowbreak{}Delta: 54-61; SIDIS/rho: 57-64), in radians: \code{th\_\allowbreak{}tarq} (angle between q and the
polarization,\allowbreak{} NaN when \code{targ\_\allowbreak{}pol\ =\ 0}), \code{phitarq}, \code{beta}, \code{phis} (Sivers angle φpq − φtarg), \code{phic} (Collins angle φpq

\begin{itemize}
\item φtarg), and the vertex values \code{betai}, \code{phisi}, \code{phici}. The kaon columns \code{saghai}/\code{factor} and the rho columns
\code{Mrho}/\code{Thrho}/\code{mmnuc} move after them. The shipped polarized-target pion example has 61 columns.
\end{itemize}
\subsection{\texorpdfstring{Values that are not measurements}{Values that are not measurements}}\label{user_manual-07u5foutputs:718-values-that-are-not-measurements}
\begin{itemize}
\item With an arm switched off (\hyperref[param:spectu5fmode]{\code{spect\_\allowbreak{}mode}}\simcidx{Parameters!spect_mode@\code{spect\_mode}} ≠ 0), that arm's focal-plane columns and \code{*ytar} are 0.0 in every event
(\hyperref[user_manual-06u5fspectrometers:64-spectu5fmode-switching-the-spectrometer-monte-carlo-off]{6.4}).
\item Other constant or NaN columns: \hyperref[user_manual-10u5ffaq:nan-or-constant-columns-in-the-tree]{10}.
\end{itemize}
\section{\texorpdfstring{\code{.\allowbreak{}hist}}{.hist}}\label{user_manual-07u5foutputs:72-hist}
The \code{.\allowbreak{}hist} is the run report, in the classic SIMC layout: the run title and
time stamps, the FLAGS block, the kinematics, the target, the central event and its cross section (\code{CENTRAL.\allowbreak{}sigcc}), the
radiative settings, the event counters (\code{Ngen\ (\allowbreak{}request)}, \code{Ntried}, \code{Ncontribute}, \code{Nco\_\allowbreak{}no\_\allowbreak{}rad\_\allowbreak{}prot}), the
normalisation (\code{wtcontr}, \code{AVERAGE.\allowbreak{}sigcc}, \code{charge}, \code{targetfac}, luminosity, \code{genvol}, \code{normfac}), the resolutions
(RECON SUMMARY), the input spectrometer limits (printed \textbf{after} the slop widening) and the limit
tables. Lines 2-3 are the BEGIN/END times; everything else depends only on the input, the seed and the mode.

\section{\texorpdfstring{\code{.\allowbreak{}geni}, \code{.\allowbreak{}gen} and \code{.\allowbreak{}gen.\allowbreak{}root}}{.geni, .gen and .gen.root}}\label{user_manual-07u5foutputs:73-geni-gen-and-genroot}
\begin{itemize}
\item \code{.\allowbreak{}geni}: for each spectrometer used, the number of trials and where they were stopped (slits, magnet entrances,
middles and exits, hut, chambers, scintillators,\allowbreak{} calorimeter),\allowbreak{} and the successes. Counter keys per spectrometer are
listed in \code{docs/\allowbreak{}manual/\allowbreak{}markdown/\allowbreak{}reference/\allowbreak{}OUTPUT\_\allowbreak{}FORMAT.\allowbreak{}md} §2. With \hyperref[param:usingu5fHMScoll]{\code{using\_\allowbreak{}HMScoll}}\simcidx{Parameters!using_HMScoll@\code{using\_HMScoll}}/\hyperref[param:usingu5fSHMScoll]{\code{using\_\allowbreak{}SHMScoll}}\simcidx{Parameters!using_SHMScoll@\code{using\_SHMScoll}} a \code{Collimator\ stops\ (\allowbreak{}pi/\allowbreak{}mu)} line is added.
\item \code{.\allowbreak{}gen}: diagnostic histograms (50 bins) of delta, yptar, xptar of both arms, Em and Pm: \code{geni} (all tries,
vertex), \code{gen} (accepted, vertex) and \code{recon} (accepted, reconstructed,\allowbreak{} weighted). \code{.\allowbreak{}gen.\allowbreak{}root} has the same histograms
as TH1D named \code{<set>\allowbreak{}\_\allowbreak{}<arm>\allowbreak{}\_\allowbreak{}<var>} (for example \code{gen\_\allowbreak{}e\_\allowbreak{}delta}, \code{recon\_\allowbreak{}Em}), filled with the \code{.\allowbreak{}gen} binning (no
under/\allowbreak{}overflow). \code{gen} and \code{geni} use −xptar and \code{recon} uses +xptar, and \code{gen\_\allowbreak{}Pm} is never filled.
\end{itemize}
\section{\texorpdfstring{Run summary (\code{.\allowbreak{}summary.\allowbreak{}json}, \code{.\allowbreak{}summary.\allowbreak{}txt})}{Run summary (.summary.\allowbreak{}json,\allowbreak{} .summary.\allowbreak{}txt)}}\label{user_manual-07u5foutputs:74-run-summary-summaryjson-summarytxt}
A JSON document (\code{"schema":\ "simc\_\allowbreak{}run\_\allowbreak{}summary"}, version 1, \code{docs/\allowbreak{}manual/\allowbreak{}markdown/\allowbreak{}reference/\allowbreak{}summary.\allowbreak{}schema.\allowbreak{}json}) with the run identification,\allowbreak{} the
counters, the normalisation,\allowbreak{} all FLAGS, the stop counters, the kinematics, the central event, the limits and the tree
location. \code{.\allowbreak{}summary.\allowbreak{}txt} is the same in readable form. the comparison tool \code{simc\_\allowbreak{}compare} reads it;
\code{hist2json} converts a \code{.\allowbreak{}hist}/\code{.\allowbreak{}geni} into the same format. A threaded run adds \code{"run":\ \{"threads":\ N\}}.

\section{\texorpdfstring{Exit status}{Exit status}}\label{user_manual-07u5foutputs:75-exit-status}
\begin{simctable}{small}{2}{l>{\raggedright\arraybackslash}p{\dimexpr 0.869\linewidth-2\tabcolsep\relax}}
%
\toprule\rowcolor{simchead}\simcth{status} & \simcth{meaning} \\\midrule\endfirsthead
\multicolumn{2}{@{}l}{\simccontinued}\\\toprule\rowcolor{simchead}\simcth{status} & \simcth{meaning} \\\midrule\endhead
\midrule\multicolumn{2}{r@{}}{\simccontinues}\\\endfoot
\bottomrule\endlastfoot
0 & success \\
1 & usage error (unknown option, missing input file, bad \code{-{}-{}threads}\simcidx{Command-line options!threads@\code{-{}-{}threads}} value) \\
2 & input or configuration error: unreadable file, rejected value (arm code 6, π0 from the neutron on H, rho on A ≥ 4, decay of a non-π/K hadron, \code{-{}-{}threads}\simcidx{Command-line options!threads@\code{-{}-{}threads}} with \code{-{}-{}fortran-{}compat}\simcidx{Command-line options!fortran-compat@\code{-{}-{}fortran-{}compat}} or RANLUX, ...) \\
3 & run error: a set-up check stopped the run (the message says \code{STOP\ during\ setup\ .\allowbreak{}.\allowbreak{}.}), a missing data file, an I/O error \\
\end{simctable}

% docs/manual/latex/user/08_running_and_reproducibility.tex - chapter of the LaTeX edition of the book "user_manual" (SIMC++ documentation).
%   Converted once from docs/manual/markdown/user/08_running_and_reproducibility.md by tools/md2latex.py; from then on THIS FILE is the source of the
%   LaTeX edition and may be edited by hand. Re-running tools/md2latex.py --write-sources for this
%   chapter OVERWRITES the hand edits (see docs/README.md).
%   Compile: cmake --build <build> --target pdf, or tools/build_docs_pdf.py (figures, xelatex passes,
%   makeindex with docs/common/latex/simcpp.ist); by hand, see docs/README.md. XeLaTeX, class
%   docs/common/latex/simcpp.cls (packages: fontspec, tikz, xcolor, hyperref, listings, longtable, booktabs,
%   titlesec, titletoc, fancyhdr, makeidx, preview, ...); fonts TeX Gyre Pagella, TeX Gyre Heros, DejaVu Sans
%   Mono. The figures are built first from docs/{manual,translation}/figures/*.tex (preview.sty) into PDFs.
\chapter{\texorpdfstring{Running modes and reproducibility}{Running modes and reproducibility}}\label{user_manual-08u5frunningu5fandu5freproducibility}
\section{\texorpdfstring{Seeds}{Seeds}}\label{user_manual-08u5frunningu5fandu5freproducibility:81-seeds}
The seed comes from, in this order: the pipeline \code{"rng":\ \{"seed":\ N\}}, then \code{-{}-{}seed\ N}\simcidx{Command-line options!seed@\code{-{}-{}seed}}, then the input's \hyperref[param:randomu5fseed]{\code{random\_\allowbreak{}seed}}\simcidx{Parameters!random_seed@\code{random\_seed}}.
\code{-{}-{}seed}\simcidx{Command-line options!seed@\code{-{}-{}seed}} overrides \hyperref[param:randomu5fseed]{\code{random\_\allowbreak{}seed}}\simcidx{Parameters!random_seed@\code{random\_seed}}, and the \code{.\allowbreak{}hist} prints the overridden value. With the default generator:

\begin{itemize}
\item \code{random\_\allowbreak{}seed\ >\ 0}: \code{std::\allowbreak{}mt19937\_\allowbreak{}64} seeded with it;
\item \code{random\_\allowbreak{}seed\ =\ 0} or absent (and other values ≤ 0 except −1): the fixed seed 5489, so the run is still reproducible;
\item \code{random\_\allowbreak{}seed\ =\ −1}: the current time. The log line \code{\ Use\ random\ seed\ =\ .\allowbreak{}.\allowbreak{}.} shows the value.
\end{itemize}
A given input, seed, mode and number of threads always produce the same files.

\section{\texorpdfstring{Random-number generators: \code{mt64} and \code{ranlux}}{Random-number generators: mt64 and ranlux}}\label{user_manual-08u5frunningu5fandu5freproducibility:82-random-number-generators-mt64-and-ranlux}
\begin{simctable}{footnotesize}{3}{>{\raggedright\arraybackslash}p{\dimexpr 0.169\linewidth-2\tabcolsep\relax}>{\raggedright\arraybackslash}p{\dimexpr 0.282\linewidth-2\tabcolsep\relax}>{\raggedright\arraybackslash}p{\dimexpr 0.488\linewidth-2\tabcolsep\relax}}
%
\toprule\rowcolor{simchead}\simcth{generator} & \simcth{pipeline} & \simcth{properties} \\\midrule\endfirsthead
\multicolumn{3}{@{}l}{\simccontinued}\\\toprule\rowcolor{simchead}\simcth{generator} & \simcth{pipeline} & \simcth{properties} \\\midrule\endhead
\midrule\multicolumn{3}{r@{}}{\simccontinues}\\\endfoot
\bottomrule\endlastfoot
\code{mt64} (default) & \code{"rng":\ "mt64"} or \code{\{"type":\ "mt64",\ "seed":\ N\}} & std::\allowbreak{}mt19937\_\allowbreak{}64; fast; allows \code{-{}-{}threads}\simcidx{Command-line options!threads@\code{-{}-{}threads}} \\
\code{ranlux} & \code{"rng":\ \{"type":\ "ranlux"\}}, optionally \code{"seed":\ N} or \code{"state":\ [25\ integers]} & RANLUX at luxury level 3, the generator of the legacy SIMC, seeded from \hyperref[param:randomu5fseed]{\code{random\_\allowbreak{}seed}}\simcidx{Parameters!random_seed@\code{random\_seed}} the same way. About 25 \% slower. \code{"state"} takes the 25 integers of a saved generator state (a legacy \code{*\_\allowbreak{}start\_\allowbreak{}random\_\allowbreak{}state.\allowbreak{}dat} file) \\
\end{simctable}
\section{\texorpdfstring{Default and legacy-compatible mode}{Default and legacy-compatible mode}}\label{user_manual-08u5frunningu5fandu5freproducibility:83-default-and-legacy-compatible-mode}
SIMC++ is a translation of the legacy Fortran SIMC. The legacy program has a number of known bugs (for example, the cut
that decides whether an event counts in the \code{.\allowbreak{}hist} yield uses the electron delta limit for the hadron). \textbf{The default
mode corrects them. The legacy-compatible mode, \code{-{}-{}fortran-{}compat}\simcidx{Command-line options!fortran-compat@\code{-{}-{}fortran-{}compat}}, reproduces the legacy behaviour exactly}, bugs
included:

\begin{simctable}{small}{3}{>{\raggedright\arraybackslash}p{\dimexpr 0.387\linewidth-2\tabcolsep\relax}>{\raggedright\arraybackslash}p{\dimexpr 0.115\linewidth-2\tabcolsep\relax}>{\raggedright\arraybackslash}p{\dimexpr 0.430\linewidth-2\tabcolsep\relax}}
%
\toprule\rowcolor{simchead}\simcth{} & \simcth{default} & \simcth{\code{-{}-{}fortran-{}compat}} \\\midrule\endfirsthead
\multicolumn{3}{@{}l}{\simccontinued}\\\toprule\rowcolor{simchead}\simcth{} & \simcth{default} & \simcth{\code{-{}-{}fortran-{}compat}} \\\midrule\endhead
\midrule\multicolumn{3}{r@{}}{\simccontinues}\\\endfoot
\bottomrule\endlastfoot
known bugs of the legacy SIMC & corrected & kept \\
with \code{ranlux}: identical output files to a legacy SIMC run of the same input & no & yes \\
\code{-{}-{}threads\ N\ >\ 1}\simcidx{Command-line options!threads@\code{-{}-{}threads}} & yes & refused \\
use it for & physics & reproducing or checking an existing legacy SIMC result \\
\end{simctable}
In a JSON input: \code{"run":\ \{"fortran\_\allowbreak{}compat":\ true\}}. The log says which mode is active
(\code{\ SIMC++:\ fortran\_\allowbreak{}compat\ =\ true\ (\allowbreak{}exact\ Fortran\ behaviour,\ known\ Fortran\ bugs\ kept)}).

\begin{simclegacy}
what each correction changes in the results, and which run types are affected, is listed in
\code{docs/\allowbreak{}translation/\allowbreak{}markdown/\allowbreak{}user\_\allowbreak{}differences.\allowbreak{}md}; how to compare SIMC++ with legacy SIMC runs statistically is in
\code{docs/\allowbreak{}translation/\allowbreak{}markdown/\allowbreak{}user\_\allowbreak{}comparing\_\allowbreak{}with\_\allowbreak{}fortran.\allowbreak{}md}.
\end{simclegacy}
\begin{figure}[htbp]
\simcfigure{compat_switch}
\caption[The two behaviour modes]{The two behaviour modes: what each Category of finding does by default and with -{}-fortran-compat,\allowbreak{} and how the switch reaches the code}\label{fig:compatu5fswitch}
\end{figure}
\section{\texorpdfstring{Reproducing legacy SIMC results exactly}{Reproducing legacy SIMC results exactly}}\label{user_manual-08u5frunningu5fandu5freproducibility:84-reproducing-legacy-simc-results-exactly}
Run the same \code{.\allowbreak{}inp} with \code{ranlux}, \code{-{}-{}fortran-{}compat}\simcidx{Command-line options!fortran-compat@\code{-{}-{}fortran-{}compat}} and one thread. The \code{.\allowbreak{}hist} (apart from its two time lines), the
\code{.\allowbreak{}geni}, the \code{.\allowbreak{}gen} and every entry of the tree are then identical to those of the legacy program. Example with a stored
legacy run (\code{\$SIMC\_\allowbreak{}REFERENCE/\allowbreak{}runs/\allowbreak{}.\allowbreak{}.\allowbreak{}.} are the project's reference runs; 1000 events, with an \hyperref[param:extrau5fdbaseu5ffile]{\code{extra\_\allowbreak{}dbase\_\allowbreak{}file}}\simcidx{Parameters!extra_dbase_file@\code{extra\_dbase\_file}}):

\begin{simcshell}[basicstyle=\ttfamily\tiny]
echo '{"pipeline": {"rng": {"type": "ranlux"}}}' > ranlux.json
R=$SIMC_REFERENCE/runs/eep_h_test_pctnames            # the legacy run: .inp, extra file, .hist, .geni, .gen, .root
simc++ $R/eep_h_test_pctnames.inp -o exact --fortran-compat --pipeline ranlux.json > exact.log
diff <(sed 2,3d exact/eep_h_test_pctnames.hist) <(sed 2,3d $R/eep_h_test_pctnames.hist) && echo ".hist: identical apart from the time lines"
cmp exact/eep_h_test_pctnames.geni $R/eep_h_test_pctnames.geni && cmp exact/eep_h_test_pctnames.gen $R/eep_h_test_pctnames.gen && echo ".geni, .gen: identical"
simc_root_identity exact/eep_h_test_pctnames.root $R/eep_h_test_pctnames.root | tail -1
\end{simcshell}
Output:

\begin{simcplain}[basicstyle=\ttfamily\tiny]
.hist: identical apart from the time lines
.geni, .gen: identical
exact/eep_h_test_pctnames.root vs .../reference/runs/eep_h_test_pctnames/eep_h_test_pctnames.root: IDENTICAL
\end{simcplain}
Notes:

\begin{itemize}
\item Use the same input, including its \hyperref[param:randomu5fseed]{\code{random\_\allowbreak{}seed}}\simcidx{Parameters!random_seed@\code{random\_seed}} (none = the RANLUX default seed). Bit-for-bit identity also needs the
same C math library (glibc \code{libm}) as the machine that made the legacy run.
\item Without \code{-{}-{}fortran-{}compat}\simcidx{Command-line options!fortran-compat@\code{-{}-{}fortran-{}compat}}, the \code{.\allowbreak{}hist} differs in the corrected lines (here two lines where the legacy program writes
integers as \code{T}/\code{F}).
\item This holds for every example input distributed with SIMC that the legacy program can run (CTest label
\code{regression\_\allowbreak{}shipped}).
\end{itemize}
\section{\texorpdfstring{Threads: \code{-{}-{}threads\ N}}{Threads: -{}-threads N}}\label{user_manual-08u5frunningu5fandu5freproducibility:85-threads---threads-n}
\code{-{}-{}threads\ N}\simcidx{Command-line options!threads@\code{-{}-{}threads}} (or \code{"pipeline":\ \{"threads":\ N\}}; the option wins) runs the event loop on N threads
(\code{docs/\allowbreak{}manual/\allowbreak{}markdown/\allowbreak{}reference/\allowbreak{}PERFORMANCE.\allowbreak{}md} §3):

\begin{itemize}
\item each thread has its own copy of every stage and its own mt64 random stream, derived from (seed, thread number);
\item the events are split: thread k makes ngen/N events (+1 for k < ngen mod N), so the total is exactly ngen (successes
for ngen > 0, tries for ngen < 0);
\item the outputs are merged in a fixed thread order: the tree entries thread after thread, the histograms and counters
summed, the normalisation computed from the merged totals;
\item \textbf{a given seed and N always give the same files}, but they are statistically equivalent to, not equal to, the
one-thread run, or a run with another N. \code{-{}-{}threads\ 1}\simcidx{Command-line options!threads@\code{-{}-{}threads}} is exactly the sequential run;
\item the summary JSON records \code{"run":\ \{"threads":\ N\}};
\item the stage warnings of the threads may be interleaved in the log (\code{-{}-{}quiet}\simcidx{Command-line options!quiet@\code{-{}-{}quiet}} silences them).
\end{itemize}
\begin{simcshell}[basicstyle=\ttfamily\tiny]
$ simc++ run.json -o out_t4 --ngen 5000 --threads 4 > out_t4.log; grep -E "Ntried|Ncontribute" out_t4/run.hist
            Ntried         =     102516
            Ncontribute    =       5000
$ simc++ run.json -o out_t4b --ngen 5000 --threads 4 > /dev/null && simc_root_identity out_t4/run.root out_t4b/run.root | tail -1
out_t4/run.root vs out_t4b/run.root: IDENTICAL
\end{simcshell}
\begin{simcwarning}
\code{-{}-{}threads\ N\ >\ 1}\simcidx{Command-line options!threads@\code{-{}-{}threads}} does not work with \code{-{}-{}fortran-{}compat}\simcidx{Command-line options!fortran-compat@\code{-{}-{}fortran-{}compat}} or RANLUX: the exact reproduction of a legacy
run needs one random stream consumed in the legacy order, and N threads use N streams. Both combinations are refused
before the run starts (exit status 2):
\end{simcwarning}
\begin{simcshell}[basicstyle=\ttfamily\tiny]
$ simc++ run.json -o out_compat --ngen 5000 --threads 4 --fortran-compat
simc++: --threads 4 (more than one event-loop thread) cannot be used with --fortran-compat: the exact Fortran reproduction needs the single random stream in the Fortran order (run with --threads 1)
\end{simcshell}
Statistical comparisons of 8-thread and one-thread runs agree (\code{docs/\allowbreak{}manual/\allowbreak{}markdown/\allowbreak{}reference/\allowbreak{}PERFORMANCE.\allowbreak{}md} §3).

\begin{figure}[htbp]
\simcfigure{parallel_mode}
\caption[The parallel event loop]{The parallel event loop: the master builds N workers with their own stages and random streams, the workers run independently,\allowbreak{} and the results are merged in a fixed order}\label{fig:parallelu5fmode}
\end{figure}
\section{\texorpdfstring{Other run options}{Other run options}}\label{user_manual-08u5frunningu5fandu5freproducibility:86-other-run-options}
\begin{itemize}
\item \code{-{}-{}ngen\ N}\simcidx{Command-line options!ngen@\code{-{}-{}ngen}} overrides \hyperref[param:ngen]{\code{ngen}}\simcidx{Parameters!ngen@\code{ngen}} (use a small N for a test; a negative N gives |N| tries, which always ends).
\item \code{-{}-{}quiet}\simcidx{Command-line options!quiet@\code{-{}-{}quiet}} silences the stage messages (the progress lines stay).
\item \code{-{}-{}print-{}pipeline}\simcidx{Command-line options!print-pipeline@\code{-{}-{}print-{}pipeline}} prints the plug-ins chosen and continues with the run.
\item \code{-{}-{}run-{}name}\simcidx{Command-line options!run-name@\code{-{}-{}run-{}name}}, \code{-{}o}: output naming (\hyperref[user_manual-04u5finputs:422-run]{4.2.2}).
\end{itemize}

% docs/manual/latex/user/09_performance.tex - chapter of the LaTeX edition of the book "user_manual" (SIMC++ documentation).
%   Converted once from docs/manual/markdown/user/09_performance.md by tools/md2latex.py; from then on THIS FILE is the source of the
%   LaTeX edition and may be edited by hand. Re-running tools/md2latex.py --write-sources for this
%   chapter OVERWRITES the hand edits (see docs/README.md).
%   Compile: cmake --build <build> --target pdf, or tools/build_docs_pdf.py (figures, xelatex passes,
%   makeindex with docs/common/latex/simcpp.ist); by hand, see docs/README.md. XeLaTeX, class
%   docs/common/latex/simcpp.cls (packages: fontspec, tikz, xcolor, hyperref, listings, longtable, booktabs,
%   titlesec, titletoc, fancyhdr, makeidx, preview, ...); fonts TeX Gyre Pagella, TeX Gyre Heros, DejaVu Sans
%   Mono. The figures are built first from docs/{manual,translation}/figures/*.tex (preview.sty) into PDFs.
\chapter{\texorpdfstring{Performance tips}{Performance tips}}\label{user_manual-09u5fperformance}
\simcabstract{Measured on the development machine (Intel i7-3770, 4 cores / 8 hardware threads, GCC 16.2 \code{-{}O3}; \code{docs/\allowbreak{}manual/\allowbreak{}markdown/\allowbreak{}reference/\allowbreak{}PERFORMANCE.\allowbreak{}md}).}
\section{\texorpdfstring{Speed}{Speed}}\label{user_manual-09u5fperformance:91-speed}
One thread, whole program including the start-up of about 0.5 s:

\begin{simctable}{small}{2}{>{\raggedright\arraybackslash}p{\dimexpr 0.272\linewidth-2\tabcolsep\relax}>{\raggedright\arraybackslash}p{\dimexpr 0.136\linewidth-2\tabcolsep\relax}}
%
\toprule\rowcolor{simchead}\simcth{case} & \simcth{µs per try} \\\midrule\endfirsthead
\multicolumn{2}{@{}l}{\simccontinued}\\\toprule\rowcolor{simchead}\simcth{case} & \simcth{µs per try} \\\midrule\endhead
\midrule\multicolumn{2}{r@{}}{\simccontinues}\\\endfoot
\bottomrule\endlastfoot
H(e,e'p) radiative & 16.3 \\
π+ H, decay on & 27.1 \\
SIDIS & 7.2 \\
rho & 14.2 \\
π0 calorimeter (NPS) & 5.8 \\
polarized-target field & 40.7 \\
\end{simctable}
The cost of a run is (number of tries) × (cost per
try). The number of tries per accepted event is Ntried/\allowbreak{}Ncontribute in the \code{.\allowbreak{}hist}: about 21 for the H(e,e'p) example,
about 3000 for rho.

\section{\texorpdfstring{Tips}{Tips}}\label{user_manual-09u5fperformance:92-tips}
\begin{itemize}
\item \textbf{Use threads.} \code{-{}-{}threads\ 4}\simcidx{Command-line options!threads@\code{-{}-{}threads}} gave 2.8-3.4× and \code{-{}-{}threads\ 8}\simcidx{Command-line options!threads@\code{-{}-{}threads}} 2.8-4.3× on the 4-core machine (\code{docs/\allowbreak{}manual/\allowbreak{}markdown/\allowbreak{}reference/\allowbreak{}PERFORMANCE.\allowbreak{}md}
§3). Beyond the number of physical cores only hyper-threading helps. Threads are not allowed with \code{-{}-{}fortran-{}compat}\simcidx{Command-line options!fortran-compat@\code{-{}-{}fortran-{}compat}} or
RANLUX (\hyperref[user_manual-08u5frunningu5fandu5freproducibility:85-threads---threads-n]{8.5}).
\item \textbf{Memory with threads.} Each thread holds its own tables; a pion run on 8 threads uses about 450 MB. Each thread also
builds its stages at the start (0.1-0.3 s per thread), which dominates very short runs.
\item \textbf{Set tight \code{SPedge} limits.} The generation volume comes from them, plus the slop. Wide-open defaults (±1e10) waste
almost every try, and the run may never end (\hyperref[user_manual-10u5ffaq:the-run-never-ends]{10}).
\item \textbf{Test with negative \hyperref[param:ngen]{\code{ngen}}\simcidx{Parameters!ngen@\code{ngen}}.} \code{-{}-{}ngen\ -{}100000}\simcidx{Command-line options!ngen@\code{-{}-{}ngen}} makes exactly 100000 tries, whatever the acceptance.
\item \textbf{mt64 is faster than RANLUX} (by about 25 \%). Use RANLUX only to reproduce legacy SIMC results.
\item \textbf{Leave out outputs you do not need.} For example, a \code{pipeline.\allowbreak{}sinks} list without \code{gen\_\allowbreak{}hist} (or \code{Nntu\ =\ 0} for no tree).
The \code{.\allowbreak{}hist} is written by \code{fortran\_\allowbreak{}hist}.
\item \textbf{Do not build with \code{-{}march=\allowbreak{}native} or \code{-{}ffast-{}math}.} They change the floating-point results; the bitwise tests and
the exact reproduction of legacy results then fail.
\item Radiative runs spend most of their time in the radiative-correction set-up of each try (the gamma function), the
target energy loss and the COSY maps. Runs without radiation are faster.
\end{itemize}


\simcpart{Help}
  {Common messages and problems, and what to do about them.}
% docs/manual/latex/user/10_faq.tex - chapter of the LaTeX edition of the book "user_manual" (SIMC++ documentation).
%   Converted once from docs/manual/markdown/user/10_faq.md by tools/md2latex.py; from then on THIS FILE is the source of the
%   LaTeX edition and may be edited by hand. Re-running tools/md2latex.py --write-sources for this
%   chapter OVERWRITES the hand edits (see docs/README.md).
%   Compile: cmake --build <build> --target pdf, or tools/build_docs_pdf.py (figures, xelatex passes,
%   makeindex with docs/common/latex/simcpp.ist); by hand, see docs/README.md. XeLaTeX, class
%   docs/common/latex/simcpp.cls (packages: fontspec, tikz, xcolor, hyperref, listings, longtable, booktabs,
%   titlesec, titletoc, fancyhdr, makeidx, preview, ...); fonts TeX Gyre Pagella, TeX Gyre Heros, DejaVu Sans
%   Mono. The figures are built first from docs/{manual,translation}/figures/*.tex (preview.sty) into PDFs.
\chapter{\texorpdfstring{FAQ and troubleshooting}{FAQ and troubleshooting}}\label{user_manual-10u5ffaq}
\section*{\texorpdfstring{The run stops with \code{bad\ targ.\allowbreak{}can\ value}}{The run stops with bad targ.can value}}\addcontentsline{toc}{section}{\texorpdfstring{The run stops with \code{bad\ targ.\allowbreak{}can\ value}}{The run stops with bad targ.can value}}\label{user_manual-10u5ffaq:the-run-stops-with-bad-targcan-value}
The input uses old \code{.} names (\code{targ.\allowbreak{}A}, \code{spec.\allowbreak{}e.\allowbreak{}P}, ...). They are ignored, so the target is left at its defaults.
Convert the file and its \hyperref[param:extrau5fdbaseu5ffile]{\code{extra\_\allowbreak{}dbase\_\allowbreak{}file}}\simcidx{Parameters!extra_dbase_file@\code{extra\_dbase\_file}} (\hyperref[user_manual-04u5finputs:413-old--names]{4.1.3}). The warning above the STOP lists the ignored names.
For a liquid target (Z < 2.4), \hyperref[param:targu25can]{\code{targ\%\allowbreak{}can}}\simcidx{Parameters!targ.can@\code{targ\%can}} must be 1, 2 or 3.

\section*{\texorpdfstring{\code{simc++:\ STOP\ during\ setup\ (\allowbreak{}legacy\ SIMC\ behaviour):\ STOP\ .\allowbreak{}.\allowbreak{}.} (exit 3)}{simc++: STOP during setup (legacy SIMC behaviour): STOP ... (exit 3)}}\addcontentsline{toc}{section}{\texorpdfstring{\code{simc++:\ STOP\ during\ setup\ (\allowbreak{}legacy\ SIMC\ behaviour):\ STOP\ .\allowbreak{}.\allowbreak{}.} (exit 3)}{simc++: STOP during setup (legacy SIMC behaviour): STOP ... (exit 3)}}\label{user_manual-10u5ffaq:simc-stop-during-setup-legacy-simc-behaviour-stop--exit-3}
The set-up checks the inputs (the messages are those of the classic SIMC). Common ones:

\begin{simctable}{small}{2}{>{\raggedright\arraybackslash}p{\dimexpr 0.555\linewidth-2\tabcolsep\relax}>{\raggedright\arraybackslash}p{\dimexpr 0.398\linewidth-2\tabcolsep\relax}}
%
\toprule\rowcolor{simchead}\simcth{message} & \simcth{cause} \\\midrule\endfirsthead
\multicolumn{2}{@{}l}{\simccontinued}\\\toprule\rowcolor{simchead}\simcth{message} & \simcth{cause} \\\midrule\endhead
\midrule\multicolumn{2}{r@{}}{\simccontinues}\\\endfoot
\bottomrule\endlastfoot
\code{I\ dont\ know\ what\ phi\ should\ be\ for\ the\ electron\ arm} (or hadron) & \hyperref[param:electronu5farm]{\code{electron\_\allowbreak{}arm}}\simcidx{Parameters!electron_arm@\code{electron\_arm}}/\hyperref[param:hadronu5farm]{\code{hadron\_\allowbreak{}arm}}\simcidx{Parameters!hadron_arm@\code{hadron\_arm}} missing or not one of 1-5, 7, 8 \\
\code{Pi-{}\ production\ from\ Hydrogen\ not\ allowed!}, \code{Coherent\ production\ from\ Hydrogen/\allowbreak{}Deuterium\ not\ allowed!}, \code{Sigma-{}\ production\ from\ Hydrogen\ not\ allowed!} & reaction and target do not fit (\hyperref[param:whichu5fpion]{\code{which\_\allowbreak{}pion}}\simcidx{Parameters!which_pion@\code{which\_pion}}, \hyperref[param:whichu5fkaon]{\code{which\_\allowbreak{}kaon}}\simcidx{Parameters!which_kaon@\code{which\_kaon}}, \hyperref[param:targu25A]{\code{targ\%\allowbreak{}A}}\simcidx{Parameters!targ.A@\code{targ\%A}}) \\
\code{pizero\_\allowbreak{}ngamma\ not\ set\ correctly(\allowbreak{}must\ be\ 1\ or\ 2),\ stopping} & π0 run without \hyperref[param:pizerou5fngamma]{\code{pizero\_\allowbreak{}ngamma}}\simcidx{Parameters!pizero_ngamma@\code{pizero\_ngamma}} \\
\code{I\ don't\ have\ ANY\ idea\ what\ (\allowbreak{}e,\allowbreak{}e'm)X\ we're\ doing!!!} & \hyperref[param:doingu5fsemi]{\code{doing\_\allowbreak{}semi}}\simcidx{Parameters!doing_semi@\code{doing\_semi}} without \hyperref[param:doingu5fpion]{\code{doing\_\allowbreak{}pion}}\simcidx{Parameters!doing_pion@\code{doing\_pion}} or \hyperref[param:doingu5fkaon]{\code{doing\_\allowbreak{}kaon}}\simcidx{Parameters!doing_kaon@\code{doing\_kaon}} \\
\code{Moron!\ one\_\allowbreak{}tail>\allowbreak{}3\ turns\ radiation\ off,\ but\ using\_\allowbreak{}rad\ wants\ it\ on.} & \hyperref[param:oneu5ftail]{\code{one\_\allowbreak{}tail}}\simcidx{Parameters!one_tail@\code{one\_tail}} and \hyperref[param:usingu5frad]{\code{using\_\allowbreak{}rad}}\simcidx{Parameters!using_rad@\code{using\_rad}} \\
\code{Idiot!\ check\ setting\ of\ de\allowbreak{}Forest\_\allowbreak{}flag} & |deForest\_\allowbreak{}flag| > 1 \\
\code{BAD\ targ.\allowbreak{}angle\ .\allowbreak{}.\allowbreak{}.} & solid target turned by more than about 58° \\
\end{simctable}
\section*{\texorpdfstring{\code{simc++:\ input\ error} or \code{configuration\ error} (exit 2)}{simc++: input error or configuration error (exit 2)}}\addcontentsline{toc}{section}{\texorpdfstring{\code{simc++:\ input\ error} or \code{configuration\ error} (exit 2)}{simc++: input error or configuration error (exit 2)}}\label{user_manual-10u5ffaq:simc-input-error-or-configuration-error-exit-2}
The input was refused before the run: arm code 6 (use 5 for the SHMS), π0 from the neutron on hydrogen, \hyperref[param:doingu5fdecay]{\code{doing\_\allowbreak{}decay}}\simcidx{Parameters!doing_decay@\code{doing\_decay}} for a
hadron that is not a π or K, rho on A ≥ 4 or combined with another reaction, \hyperref[param:usingu5ftgtu5ffield]{\code{using\_\allowbreak{}tgt\_\allowbreak{}field}}\simcidx{Parameters!using_tgt_field@\code{using\_tgt\_field}} without \hyperref[param:tgtu5ffieldu5ffile]{\code{tgt\_\allowbreak{}field\_\allowbreak{}file}}\simcidx{Parameters!tgt_field_file@\code{tgt\_field\_file}},
\code{-{}-{}threads}\simcidx{Command-line options!threads@\code{-{}-{}threads}} with \code{-{}-{}fortran-{}compat}\simcidx{Command-line options!fortran-compat@\code{-{}-{}fortran-{}compat}} or RANLUX, an unreadable JSON file. The message says which.

\section*{\texorpdfstring{The run never ends}{The run never ends}}\addcontentsline{toc}{section}{\texorpdfstring{The run never ends}{The run never ends}}\label{user_manual-10u5ffaq:the-run-never-ends}
With \code{ngen\ >\ 0}, SIMC++ runs until it has \hyperref[param:ngen]{\code{ngen}}\simcidx{Parameters!ngen@\code{ngen}} \textbf{accepted} events. If almost nothing is accepted
(wrong angles or momenta, wide-open \code{SPedge} defaults, a kinematics outside the acceptance),\allowbreak{} it runs for a very long time.
Test with a negative \hyperref[param:ngen]{\code{ngen}}\simcidx{Parameters!ngen@\code{ngen}} (\code{-{}-{}ngen\ -{}10000}\simcidx{Command-line options!ngen@\code{-{}-{}ngen}} = 10000 tries) and look at Ncontribute and the \code{.\allowbreak{}geni} stop counters. The
progress lines ("... successes so far") show the rate.

\section*{\texorpdfstring{All weights or the luminosity are 0}{All weights or the luminosity are 0}}\addcontentsline{toc}{section}{\texorpdfstring{All weights or the luminosity are 0}{All weights or the luminosity are 0}}\label{user_manual-10u5ffaq:all-weights-or-the-luminosity-are-0}
Check \hyperref[param:targu25abundancy]{\code{targ\%\allowbreak{}abundancy}}\simcidx{Parameters!targ.abundancy@\code{targ\%abundancy}} (in \%, 100 for a pure target: the default 0 gives zero luminosity),\allowbreak{} \hyperref[param:EXPERu25charge]{\code{EXPER\%\allowbreak{}charge}}\simcidx{Parameters!EXPER.charge@\code{EXPER\%charge}}, and for
A(e,e'p) with Benhar spectral functions \hyperref[param:transparency]{\code{transparency}}\simcidx{Parameters!transparency@\code{transparency}}. Delta runs have no cross section at all
(\hyperref[user_manual-05u5freactions:55-eepu3c0-in-the-delta-region-doingu5fdelta]{5.5}).

\section*{\texorpdfstring{Data files not found}{Data files not found}}\addcontentsline{toc}{section}{\texorpdfstring{Data files not found}{Data files not found}}\label{user_manual-10u5ffaq:data-files-not-found}
\code{simc++\ -{}-{}print-{}data-{}dir} shows where the tables are looked for. Set \code{SIMC\_\allowbreak{}DATA\_\allowbreak{}DIR} to a directory with the layout of
\code{data/} (\code{matrices/}, \code{pion/}, \code{kaon/}, \code{sidis/}, \code{spectral/}, \code{field/}), or use \code{-{}-{}data-{}dir}\simcidx{Command-line options!data-dir@\code{-{}-{}data-{}dir}} for the spectral and field
files (\hyperref[user_manual-02u5finstallation:25-where-the-data-files-are-found]{2.5}). An installed \code{simc++} uses \code{\$SIMCPP\_\allowbreak{}PREFIX/\allowbreak{}share/\allowbreak{}simc/\allowbreak{}data}.

\section*{\texorpdfstring{\code{error\ while\ loading\ shared\ libraries} or crashes in ROOT at start-up}{error while loading shared libraries or crashes in ROOT at start-up}}\addcontentsline{toc}{section}{\texorpdfstring{\code{error\ while\ loading\ shared\ libraries} or crashes in ROOT at start-up}{error while loading shared libraries or crashes in ROOT at start-up}}\label{user_manual-10u5ffaq:error-while-loading-shared-libraries-or-crashes-in-root-at-start-up}
\code{LD\_\allowbreak{}LIBRARY\_\allowbreak{}PATH} must contain the \code{lib} of the ROOT 6.40 that SIMC++ was built with, and \textbf{not} a directory with an older
ROOT or libtbb (for example \code{/\allowbreak{}usr/\allowbreak{}local/\allowbreak{}lib}) (\hyperref[user_manual-02u5finstallation:22-requirements]{2.2}).

\section*{\texorpdfstring{NaN or constant columns in the tree}{NaN or constant columns in the tree}}\addcontentsline{toc}{section}{\texorpdfstring{NaN or constant columns in the tree}{NaN or constant columns in the tree}}\label{user_manual-10u5ffaq:nan-or-constant-columns-in-the-tree}
\begin{simctable}{small}{3}{>{\raggedright\arraybackslash}p{\dimexpr 0.305\linewidth-2\tabcolsep\relax}>{\raggedright\arraybackslash}p{\dimexpr 0.272\linewidth-2\tabcolsep\relax}>{\raggedright\arraybackslash}p{\dimexpr 0.355\linewidth-2\tabcolsep\relax}}
%
\toprule\rowcolor{simchead}\simcth{column} & \simcth{when} & \simcth{why} \\\midrule\endfirsthead
\multicolumn{3}{@{}l}{\simccontinued}\\\toprule\rowcolor{simchead}\simcth{column} & \simcth{when} & \simcth{why} \\\midrule\endhead
\midrule\multicolumn{3}{r@{}}{\simccontinues}\\\endfoot
\bottomrule\endlastfoot
\code{*xfp\ *xpfp\ *yfp\ *ypfp\ *ytar} = 0 & \hyperref[param:spectu5fmode]{\code{spect\_\allowbreak{}mode}}\simcidx{Parameters!spect_mode@\code{spect\_mode}} ≠ 0, for the switched-off arm & no spectrometer simulated (\hyperref[user_manual-06u5fspectrometers:64-spectu5fmode-switching-the-spectrometer-monte-carlo-off]{6.4}) \\
\code{th\_\allowbreak{}tarq} NaN & \hyperref[param:usingu5ftgtu5ffield]{\code{using\_\allowbreak{}tgt\_\allowbreak{}field}}\simcidx{Parameters!using_tgt_field@\code{using\_tgt\_field}} with \code{targ\_\allowbreak{}pol\ =\ 0} & no polarization direction \\
\code{sigcent}, \code{centjac} = 0 & SIDIS/rho & not computed (central cross section disabled) \\
\code{zi}, \code{pt2i}, \code{xfermi} = 0 & rho & not defined for rho \\
\code{xcal\_\allowbreak{}gamma*} = −1e10 & π0 runs & the photon missed the calorimeter \\
\code{Weight}, \code{siglab} without cross section & Delta runs & no Delta model (\hyperref[user_manual-05u5freactions:55-eepu3c0-in-the-delta-region-doingu5fdelta]{5.5}) \\
\end{simctable}
With \code{-{}-{}fortran-{}compat}\simcidx{Command-line options!fortran-compat@\code{-{}-{}fortran-{}compat}} a few more columns carry the legacy program's NaN or meaningless values (\code{pfermi}, Delta \code{sigcm},
kaon \code{saghai}; see \code{docs/\allowbreak{}translation/\allowbreak{}markdown/\allowbreak{}user\_\allowbreak{}differences.\allowbreak{}md}).

\code{simc\_\allowbreak{}compare} counts and excludes non-finite entries.

\section*{\texorpdfstring{My numbers differ from a legacy SIMC run}{My numbers differ from a legacy SIMC run}}\addcontentsline{toc}{section}{\texorpdfstring{My numbers differ from a legacy SIMC run}{My numbers differ from a legacy SIMC run}}\label{user_manual-10u5ffaq:my-numbers-differ-from-a-legacy-simc-run}
The default mode corrects known bugs of the legacy program and uses another random-number generator, so a default run is
statistically equivalent, not identical, apart from the corrections. To reproduce the legacy result exactly use \code{ranlux},
\code{-{}-{}fortran-{}compat}\simcidx{Command-line options!fortran-compat@\code{-{}-{}fortran-{}compat}} and one thread (\hyperref[user_manual-08u5frunningu5fandu5freproducibility:84-reproducing-legacy-simc-results-exactly]{8.4}).
The statistical comparison and the list of expected differences are in \code{docs/\allowbreak{}translation/}.

\section*{\texorpdfstring{Two runs with \code{-{}-{}threads} give different numbers than without}{Two runs with -{}-threads give different numbers than without}}\addcontentsline{toc}{section}{\texorpdfstring{Two runs with \code{-{}-{}threads} give different numbers than without}{Two runs with -{}-threads give different numbers than without}}\label{user_manual-10u5ffaq:two-runs-with---threads-give-different-numbers-than-without}
That is expected: N threads use N random streams. A given seed and N are reproducible; \code{-{}-{}threads\ 1}\simcidx{Command-line options!threads@\code{-{}-{}threads}} equals the
sequential run (\hyperref[user_manual-08u5frunningu5fandu5freproducibility:85-threads---threads-n]{8.5}).

\section*{\texorpdfstring{Warnings about \code{pt\_\allowbreak{}b\_\allowbreak{}param}, \code{sigc\_\allowbreak{}*}, \code{random\_\allowbreak{}state\_\allowbreak{}file}, unknown names}{Warnings about pt\_\allowbreak{}b\_\allowbreak{}param,\allowbreak{} sigc\_*, random\_\allowbreak{}state\_\allowbreak{}file,\allowbreak{} unknown names}}\addcontentsline{toc}{section}{\texorpdfstring{Warnings about \code{pt\_\allowbreak{}b\_\allowbreak{}param}, \code{sigc\_\allowbreak{}*}, \code{random\_\allowbreak{}state\_\allowbreak{}file}, unknown names}{Warnings about pt\_\allowbreak{}b\_\allowbreak{}param,\allowbreak{} sigc\_*, random\_\allowbreak{}state\_\allowbreak{}file,\allowbreak{} unknown names}}\label{user_manual-10u5ffaq:warnings-about-ptu5fbu5fparam-sigcu5f-randomu5fstateu5ffile-unknown-names}
These inputs are read and ignored (\hyperref[param:randomu5fstateu5ffile]{\code{random\_\allowbreak{}state\_\allowbreak{}file}}\simcidx{Parameters!random_state_file@\code{random\_state\_file}} is not supported; use the \code{ranlux} \code{"state"} option). Remove them from the
input to silence the warnings. Unknown names are usually typos or legacy names.

\section*{\texorpdfstring{The \code{.\allowbreak{}hist} "Input Spectrometer Limits" differ from my input}{The .hist "Input Spectrometer Limits" differ from my input}}\addcontentsline{toc}{section}{\texorpdfstring{The \code{.\allowbreak{}hist} "Input Spectrometer Limits" differ from my input}{The .hist "Input Spectrometer Limits" differ from my input}}\label{user_manual-10u5ffaq:the-hist-input-spectrometer-limits-differ-from-my-input}
They are printed \textbf{after} widening by the spectrometer slop.

\section*{\texorpdfstring{Can I use my own spectrometer,\allowbreak{} reaction or cross section?}{Can I use my own spectrometer,\allowbreak{} reaction or cross section?}}\addcontentsline{toc}{section}{\texorpdfstring{Can I use my own spectrometer,\allowbreak{} reaction or cross section?}{Can I use my own spectrometer,\allowbreak{} reaction or cross section?}}\label{user_manual-10u5ffaq:can-i-use-my-own-spectrometer-reaction-or-cross-section}
Yes: every stage is a plug-in. See the developer manual (\code{docs/\allowbreak{}manual/\allowbreak{}markdown/\allowbreak{}developer/}) and \code{docs/\allowbreak{}manual/\allowbreak{}markdown/\allowbreak{}reference/\allowbreak{}EXTENDING.\allowbreak{}md}; the
\code{examples/\allowbreak{}ideal\_\allowbreak{}box} and \code{examples/\allowbreak{}out\_\allowbreak{}of\_\allowbreak{}tree\_\allowbreak{}plugin} directories are complete examples.



\simcappendix
% docs/manual/latex/user/VERIFIED_EXAMPLES.tex - chapter of the LaTeX edition of the book "user_manual" (SIMC++ documentation).
%   Converted once from docs/manual/markdown/user/VERIFIED_EXAMPLES.md by tools/md2latex.py; from then on THIS FILE is the source of the
%   LaTeX edition and may be edited by hand. Re-running tools/md2latex.py --write-sources for this
%   chapter OVERWRITES the hand edits (see docs/README.md).
%   Compile: cmake --build <build> --target pdf, or tools/build_docs_pdf.py (figures, xelatex passes,
%   makeindex with docs/common/latex/simcpp.ist); by hand, see docs/README.md. XeLaTeX, class
%   docs/common/latex/simcpp.cls (packages: fontspec, tikz, xcolor, hyperref, listings, longtable, booktabs,
%   titlesec, titletoc, fancyhdr, makeidx, preview, ...); fonts TeX Gyre Pagella, TeX Gyre Heros, DejaVu Sans
%   Mono. The figures are built first from docs/{manual,translation}/figures/*.tex (preview.sty) into PDFs.
\chapter{\texorpdfstring{Verified examples (transcript)}{Verified examples (transcript)}}\label{user_manual-VERIFIEDu5fEXAMPLES}
\simcabstract{Every command shown in this manual was run as recorded below. Code and manual text: commit \code{b045ce6} (branch
\code{translation}); this transcript file was added in the next commit. Toolchain: GCC 16.2.1, ROOT 6.40.04, glibc 2.43,
Release build installed with \code{cmake\ -{}-{}install}.}
\textbf{Later wording change.} Three diagnostic messages were reworded afterwards ("Fortran STOP during the setup" → "STOP during setup (legacy SIMC behaviour)",\allowbreak{} "IGNORED, as by the Fortran" → "IGNORED (legacy SIMC behaviour; ...)", and the corrected-mode banner). Those lines below show the current wording; they were re-run on the final commit. No result changed.

\textbf{Paths.} The outputs are real, but the machine-specific paths of the run were rewritten into the layout of
\hyperref[user_manual-02u5finstallation:21-directory-layout-used-in-this-manual]{2.1}: \code{\textasciitilde{}/\allowbreak{}src/\allowbreak{}simcpp} (\$SIMCPP\_\allowbreak{}SRC),\allowbreak{} \code{\textasciitilde{}/\allowbreak{}build/\allowbreak{}simcpp}
(\$SIMCPP\_\allowbreak{}BUILD),\allowbreak{} \code{\textasciitilde{}/\allowbreak{}build/\allowbreak{}simcpp-{}validation} (the comparison tool's build), \code{\textasciitilde{}/\allowbreak{}opt/\allowbreak{}simcpp} (\$SIMCPP\_\allowbreak{}PREFIX),\allowbreak{}
\code{\textasciitilde{}/\allowbreak{}src/\allowbreak{}simc\_\allowbreak{}gfortran} (\$SIMC\_\allowbreak{}FORTRAN),\allowbreak{} \code{\textasciitilde{}/\allowbreak{}simc\_\allowbreak{}reference} (\$SIMC\_\allowbreak{}REFERENCE),\allowbreak{} \code{\textasciitilde{}/\allowbreak{}opt/\allowbreak{}root} (\$ROOTSYS) and \code{\textasciitilde{}/\allowbreak{}work}
(the working directory). The environment was:

\begin{simcplain}[basicstyle=\ttfamily\footnotesize]
export ROOTSYS=~/opt/root LD_LIBRARY_PATH=$ROOTSYS/lib
export SIMCPP_SRC=~/src/simcpp SIMCPP_BUILD=~/build/simcpp SIMCPP_PREFIX=~/opt/simcpp
export SIMC_FORTRAN=~/src/simc_gfortran SIMC_REFERENCE=~/simc_reference
export EX=$SIMCPP_PREFIX/share/simc/docs/manual/markdown/user/examples
export PATH=$SIMCPP_PREFIX/bin:$ROOTSYS/bin:/usr/bin:/bin
\end{simcplain}
Each block ran in a fresh \code{bash\ -{}c} in one working directory; its output and non-zero exit status follow it. The build
commands of chapter 2 were run with this layout without compiler warnings. The full test suite is run by
\code{tools/\allowbreak{}check\_\allowbreak{}all.\allowbreak{}sh} on the final commit. The \code{ctest\ -{}L\ docs} run at the end was made before this file existed, so the
link test then reported only the not-yet-existing \code{VERIFIED\_\allowbreak{}EXAMPLES.\allowbreak{}md} links, and \code{docs\_\allowbreak{}no\_\allowbreak{}machine\_\allowbreak{}paths} reported
files of other documents that were still being edited. The \code{\$SIMC\_\allowbreak{}FORTRAN} and \code{\$SIMC\_\allowbreak{}REFERENCE} commands (chapter 8.4
and \code{docs/\allowbreak{}translation/\allowbreak{}markdown/\allowbreak{}user\_\allowbreak{}comparing\_\allowbreak{}with\_\allowbreak{}fortran.\allowbreak{}md}) need the optional legacy data.

\begin{simcshell}[basicstyle=\ttfamily\tiny]
$ simc++ --help
usage: simc++ <file.inp|file.json> [-o OUTDIR] [--seed N] [--ngen N] [--pipeline FILE]
               [--run-name NAME] [--data-dir DIR] [--infiles-dir DIR] [--quiet] [--print-pipeline]
               [--fortran-compat] [--threads N] [--print-data-dir]
  --fortran-compat  exact Fortran behaviour (known Fortran bugs kept; bitwise validation)
  --threads N       event loop on N threads (independent mt64 streams per thread; a given seed and N
                    always give the same output; not with --fortran-compat or rng ranlux). Default 1.
  --print-data-dir  print the default data directory (env SIMC_DATA_DIR overrides it) and exit

$ simc++ --print-data-dir
~/opt/simcpp/share/simc/data

$ SIMC_DATA_DIR=/data/simc simc++ --print-data-dir
/data/simc

$ simc++ $EX/h_elastic_rad.inp -o out --ngen 5000 > out.log; echo "exit status $?"
exit status 0

$ ls out
h_elastic_rad.gen
h_elastic_rad.geni
h_elastic_rad.gen.root
h_elastic_rad.hist
h_elastic_rad.root
h_elastic_rad.summary.json
h_elastic_rad.summary.txt

$ grep -E "fortran_compat|random seed" out.log
 Use random seed = 12345
 SIMC++: fortran_compat = false (clear legacy-SIMC bugs corrected, docs/translation/markdown/06_catalogue.md)

$ grep -E "^ *(Ngen|Ntried|Ncontribute)|wtcontr|normfac" out/h_elastic_rad.hist
            Ngen (request) =       5000
            Ntried         =     102995
            Ncontribute    =       5000
               MeV: wtcontr=   0.20906661E+03
                   normfac =     0.407176E+06       

$ cat > peek_h10.C <<'MACRO'
  // peek_h10.C: open a SIMC++ tree, print its size and the weighted Em distribution
  void peek_h10(const char* file = "out/h_elastic_rad.root") {
    TFile f(file);
    TTree* t = (TTree*)f.Get("h10");
    std::cout << t->GetEntries() << " events, " << t->GetListOfBranches()->GetEntries() << " branches" << std::endl;
    t->Scan("hsdelta:ssdelta:W:Em:Weight", "", "", 5);
    TH1D h("h", "E_{m} (weighted);E_{m} [GeV]", 100, -0.02, 0.2);
    t->Project("h", "Em", "Weight");
    std::cout << "sum of weights in [-0.02, 0.2] GeV: " << h.Integral() << ", mean Em " << h.GetMean() << " GeV" << std::endl;
  }
  MACRO
  root -l -b -q peek_h10.C
Processing peek_h10.C...
5000 events, 46 branches
************************************************************************
*    Row   *   hsdelta *   ssdelta *         W *        Em *    Weight *
************************************************************************
*        0 * 6.8000464 * -2.967325 * 1.0069015 * 0.0488890 * 0.0002927 *
*        1 * 2.4381718 * -0.939918 * 0.9532427 * 0.0063532 * 0.0004068 *
*        2 * -2.744465 * -9.402965 * 1.4675965 * 0.8642277 * 0.0004998 *
*        3 * -1.449645 * 0.4658077 * 0.9453506 * 0.0016173 * 0.0005580 *
*        4 * 0.9987937 * -1.216251 * 0.9978361 * 0.0712616 * 0.0004664 *
************************************************************************
sum of weights in [-0.02, 0.2] GeV: 2.24626, mean Em 0.0176472 GeV

$ inp2json -o run.json $EX/h_elastic_rad.inp && head -c 300 run.json; echo
{
  "_source": {
    "files": [
      "~/opt/simcpp/share/simc/docs/manual/markdown/user/examples/h_elastic_rad.inp"
    ],
    "generator": "inp2json (SIMC++ P1.3)",
    "not_registered": {},
    "note": "parameters = every dbase.f:regallvars variable, in .inp unit

$ simc++ run.json -o out_json --ngen 5000 > out_json.log && diff <(sed 2,3d out/h_elastic_rad.hist) <(sed 2,3d out_json/run.hist) && echo ".hist identical (apart from the time lines)"
.hist identical (apart from the time lines)

$ simc_root_identity out/h_elastic_rad.root out_json/run.root
  h10                      TTree  5000 entries x 46 leaves: identical
out/h_elastic_rad.root vs out_json/run.root: IDENTICAL

$ simc++ $SIMCPP_PREFIX/share/simc/data/config/h_elastic_rad.json -o out_cfg --ngen 5000 > out_cfg.log; echo "exit status $?"; ls out_cfg | head -3
exit status 0
h_elastic_rad.gen
h_elastic_rad.geni
h_elastic_rad.gen.root

$ $SIMCPP_BUILD/bin/simc++ $SIMCPP_SRC/docs/manual/markdown/user/examples/h_elastic_rad.inp -o out_readme --ngen 5000 > /dev/null; echo "exit status $?"
exit status 0

$ simc++ run.json -o out_t4 --ngen 5000 --threads 4 > out_t4.log; echo "exit status $?"; grep -E "Ntried|Ncontribute" out_t4/run.hist
exit status 0
            Ntried         =     102516
            Ncontribute    =       5000

$ simc++ run.json -o out_t4b --ngen 5000 --threads 4 > /dev/null && simc_root_identity out_t4/run.root out_t4b/run.root | tail -1
out_t4/run.root vs out_t4b/run.root: IDENTICAL

$ simc++ run.json -o out_compat --ngen 5000 --threads 4 --fortran-compat
simc++: --threads 4 (more than one event-loop thread) cannot be used with --fortran-compat: the exact Fortran reproduction needs the single random stream in the Fortran order (run with --threads 1)
[exit 2]

$ simc++ run.json -o out_ideal --ngen 200 --pipeline $SIMCPP_PREFIX/share/simc/examples/ideal_box/hadron_arm.json --print-pipeline --quiet | sed -n '/"spectrometers"/,/^  }/p'
  "spectrometers": {
    "electron": {
      "type": ""
    },
    "hadron": {
      "delta_max": 8.0,
      "delta_min": -8.0,
      "type": "ideal_box",
      "xptar_max": 0.05,
      "xptar_min": -0.05,
      "yptar_max": 0.03,
      "yptar_min": -0.03
    }
  },

$ simc++ $EX/legacy_names.inp -o legacy > legacy.log 2> legacy.err; echo "exit status $?"; cat legacy.err
exit status 3
simc++: WARNING: 46 input variable(s) are not SIMC variables and are IGNORED (legacy SIMC behaviour; their defaults are used):
  EXPER.charge, spec.e.P, spec.e.theta, spec.p.P, spec.p.theta, targ.A
  targ.Z, targ.mass_amu, targ.mrec_amu, targ.rho, targ.thick, targ.angle
  targ.abundancy, targ.can, SPedge.e.delta.min, SPedge.e.delta.max, SPedge.e.yptar.min, SPedge.e.yptar.max
  SPedge.e.xptar.min, SPedge.e.xptar.max, SPedge.p.delta.min, SPedge.p.delta.max, SPedge.p.yptar.min, SPedge.p.yptar.max
  SPedge.p.xptar.min, SPedge.p.xptar.max, gen.xwid, gen.ywid, targ.fr_pattern, targ.fr1
  targ.fr2, targ.xoffset, targ.yoffset, targ.zoffset, spec.e.offset.x, spec.e.offset.y
  spec.e.offset.z, spec.e.offset.xptar, spec.e.offset.yptar, spec.p.offset.x, spec.p.offset.y, spec.p.offset.z
  spec.p.offset.xptar, spec.p.offset.yptar, cuts.Em.min, cuts.Em.max
  Legacy '.' names (targ.A, spec.e.P, ...) must be converted to '%' names with simc_gfortran/infiles/convert_inputfile.sh.
simc++: STOP during setup (legacy SIMC behaviour): STOP bad targ.can value
simc++: note: the input has 46 ignored variable(s) (listed above), so their defaults were used — the likely cause of this STOP. Convert legacy '.' names with simc_gfortran/infiles/convert_inputfile.sh.

$ sed -E ':a; s/^(\s*(targ|spec|SPedge|gen|cuts|EXPER)[A-Za-z_%]*)\./\1%/; ta' $EX/legacy_names.inp > converted.inp
  simc++ converted.inp -o converted_out > converted.log 2>&1; echo "exit status $?"; grep -c "IGNORED" converted.log; grep -E "Ntried|Ncontribute" converted_out/converted.hist
exit status 0
0
            Ntried         =    4158458
            Ncontribute    =     200000

$ echo '{"pipeline": {"rng": {"type": "ranlux"}}}' > ranlux.json
  R=$SIMC_REFERENCE/runs/eep_h_test_pctnames
  simc++ $R/eep_h_test_pctnames.inp -o exact --fortran-compat --pipeline ranlux.json > exact.log; echo "exit status $?"
  diff <(sed 2,3d exact/eep_h_test_pctnames.hist) <(sed 2,3d $R/eep_h_test_pctnames.hist) && echo ".hist: identical apart from the time lines"
  cmp exact/eep_h_test_pctnames.geni $R/eep_h_test_pctnames.geni && cmp exact/eep_h_test_pctnames.gen $R/eep_h_test_pctnames.gen && echo ".geni, .gen: identical"
  simc_root_identity exact/eep_h_test_pctnames.root $R/eep_h_test_pctnames.root | tail -1
exit status 0
.hist: identical apart from the time lines
.geni, .gen: identical
exact/eep_h_test_pctnames.root vs ~/simc_reference/runs/eep_h_test_pctnames/eep_h_test_pctnames.root: IDENTICAL

$ R=$SIMC_REFERENCE/runs/eep_h_test_pctnames
  simc++ $R/eep_h_test_pctnames.inp -o exact_corr --pipeline ranlux.json > /dev/null; diff <(sed 2,3d exact_corr/eep_h_test_pctnames.hist) <(sed 2,3d $R/eep_h_test_pctnames.hist) | head -12
65c65
<                using_Eloss= T        using_Coulomb= F        deForest_flag= 0
---
>                using_Eloss= T        using_Coulomb= F        deForest_flag= F
79c79
<       use_expon=  0   include_hard=  T    calc_spence=  T
---
>       use_expon=  F   include_hard=  T    calc_spence=  T

$ cat > arm6.json <<'JSON'
  {"parameters": {"electron_arm": 6, "hadron_arm": 1}}
  JSON
  simc++ arm6.json -o arm6
simc++: input error: arm6.json: parameters: electron_arm = 6 was the retired SHMS LSA tune; use 5 = SHMS
[exit 2]

$ python3 -c '
  import json
  j = json.load(open("run.json"))
  j["parameters"]["ngen"] = 2000
  j["parameters"]["Ngne"] = 5           # a typo: reported and ignored
  j["run"] = {"name": "edited", "output_dir": "edited_out", "fortran_compat": False}
  json.dump(j, open("edited.json", "w"), indent=1)'
  simc++ edited.json > edited.log; echo "exit status $?"; ls edited_out
simc++: warning: unknown parameter 'Ngne' (not registered in dbase.f:regallvars), ignored
exit status 0
edited.gen
edited.geni
edited.gen.root
edited.hist
edited.root
edited.summary.json
edited.summary.txt

$ simc++ $SIMC_FORTRAN/infiles/eep_h_elastic_rad_A.inp -o l2 --seed 4242 --threads 4 > l2.log; echo "exit status $?"
exit status 0

$ $SIMCPP_BUILD-validation/simc_compare --a $SIMC_REFERENCE/runs/eep_h_elastic_rad_A --b l2/eep_h_elastic_rad_A.summary.json --out cmp \
                         --config $SIMCPP_SRC/validation/config/compare_default.json > cmp.log; echo "exit status $?"
  ls cmp; sed -n '1,25p' cmp/summary.txt
exit status 0
compare.pdf
runA.summary.json
runB.summary.json
summary.json
summary.txt
simc_compare  A=A  B=B
A: ~/simc_reference/runs/eep_h_elastic_rad_A/eep_h_elastic_rad_A.hist  [fortran-hist]  tree ~/simc_reference/runs/eep_h_elastic_rad_A/eep_h_elastic_rad_A.root (200000 entries)
B: l2/eep_h_elastic_rad_A.summary.json  [summary-json]  tree l2/eep_h_elastic_rad_A.root (200000 entries)
normalisation: A Weight/Ntried, B Weight/Ntried
VERDICT: PASS   tests counted N=214  failed (p<0.01) F=0  expected 2.14  fraction 0.0000  P(X>=F|N,alpha)=1.000e+00 (need >= 1.000e-03)  gross failures=0 (det mismatch or p<1.000e-05)
worst p = 1.485e-02  (ssdeltai [chi2_ww])   largest |pull| = 1.859  (hms.scin/Ntried)

RUN-LEVEL TESTS
quantity                                      method                   A       errA             B       errB     pull          p  status
acceptance Ncontribute/Ntried                 binomial         0.0481844   0.000105     0.0481577   0.000105    0.179      0.858  ok
yield sum(Weight) [Weight/Ntried]             pull           2.48294e-05   5.64e-08   2.48231e-05   5.64e-08   0.0798      0.936  ok
wtcontr                                       pull               207.354      0.472       207.306      0.472   0.0726      0.942  ok  sigma = wtcontr * relerr: A pass_cuts subset, compat window; B pass_cuts subset, corrected window
AVERAGE.sigcc                                 pull           0.000594383   3.67e-07   0.000594476   3.67e-07   -0.179      0.858  ok
normfac                                       pull                404142        882        403918        881     0.18      0.858  ok (info)  Var(Ntried)/Ntried^2=(1-acc)/Ncontribute; not counted (duplicates acceptance)
sigcc_central                                 det            0.000579704   1.08e-09   0.000579704          -        -          1  ok  tol 1.08e-09
charge_mC                                     det                      1      6e-06             1          -        -          1  ok  tol 6e-06
targetfac                                     det            2.67954e-09   7.68e-15   2.67954e-09          -        -          1  ok  tol 7.68e-15
luminosity_ub                                 det            3.73199e+08        873   3.73199e+08          -        -          1  ok  tol 873
luminosity_GeV2                               det                1453.16    0.00645       1453.16          -        -          1  ok  tol 0.00645
genvol                                        det              0.0224744   7.25e-08     0.0224744          -        -          1  ok  tol 7.25e-08
Ngen (request)                                det                 200000          -        200000          -        -          1  ok
hms.trials/Ntried                             binomial          0.874289   0.000163      0.874455   0.000163   -0.721      0.471  ok
hms.slit_hor/Ntried                           binomial          0.671944    0.00023      0.671843    0.00023    0.308      0.758  ok
hms.slit_vert/Ntried                          binomial          0.146327   0.000173      0.146589   0.000174    -1.07      0.286  ok
hms.slit_oct/Ntried                           binomial        0.00695011   4.08e-05    0.00697492   4.08e-05    -0.43      0.667  ok

$ (cd $SIMCPP_BUILD && ctest -L docs)
Test project ~/build/simcpp
    Start 434: docs_user_manual_parameters
1/7 Test #434: docs_user_manual_parameters ......   Passed    0.11 sec
    Start 435: docs_user_manual_links
2/7 Test #435: docs_user_manual_links ...........   Passed    0.04 sec
    Start 436: docs_no_machine_paths
3/7 Test #436: docs_no_machine_paths ............***Failed    0.03 sec
    Start 437: docs_dev_manual
4/7 Test #437: docs_dev_manual ..................   Passed    0.04 sec
    Start 438: docs_dev_manual_links
5/7 Test #438: docs_dev_manual_links ............   Passed    0.04 sec
    Start 439: docs_translation
6/7 Test #439: docs_translation .................   Passed    0.19 sec
    Start 440: docs_translation_links
7/7 Test #440: docs_translation_links ...........   Passed    0.04 sec

86% tests passed, 1 tests failed out of 7

Label Time Summary:
docs    =   0.48 sec*proc (7 tests)

Total Test time (real) =   0.57 sec

The following tests FAILED:
	436 - docs_no_machine_paths (Failed)                    docs
Errors while running CTest
Output from these tests are in: ~/build/simcpp/Testing/Temporary/LastTest.log
Use "--rerun-failed --output-on-failure" to re-run the failed cases verbosely.
[exit 8]
\end{simcshell}

\simcbackmatter
\end{document}
